\documentclass[UTF8,11pt,a4paper,openany,fontset=windows]{ctexbook}

% ============================ 字体 ============================
\setCJKmainfont{SimSun}[BoldFont=SimHei, ItalicFont=KaiTi, BoldItalicFont=SimHei]
\setCJKsansfont{Microsoft YaHei}
\setCJKmonofont{FangSong}
\setmainfont{Palatino Linotype}
\setsansfont{Segoe UI}
\setmonofont{Consolas}[Scale=0.86]

% ============================ 版式 ============================
\usepackage[a4paper,left=2.5cm,right=2.5cm,top=2.6cm,bottom=2.4cm]{geometry}
\usepackage{fancyhdr}
\usepackage{titlesec}
\usepackage[dvipsnames]{xcolor}
\usepackage{tcolorbox}
\tcbuselibrary{skins,breakable,listings}
\usepackage{amsmath,amssymb,amsthm,mathtools,bm}
\usepackage{booktabs,longtable,multirow,array,tabularx}
\usepackage{enumitem}
\usepackage{graphicx}
\graphicspath{{../reports/figures/}{./figures/}}
\usepackage{tikz}
\usepackage{listings}
\usepackage[hidelinks,bookmarksnumbered]{hyperref}
\usepackage{caption}
\xeCJKDeclareCharClass{CJK}{"2192, "2190, "00D7, "2248, "2265, "2264, "00B1, "2212}
\captionsetup{font=small,labelfont={sf,bf},labelsep=quad}

\definecolor{ink}{HTML}{1F2937}
\definecolor{accent}{HTML}{2F5D8A}
\definecolor{accent2}{HTML}{7A4E8C}
\definecolor{soft}{HTML}{EEF3F8}
\definecolor{warn}{HTML}{B4532A}
\definecolor{good}{HTML}{2E6B4F}
\definecolor{warms}{HTML}{FBF2EC}

\titleformat{\chapter}[block]
  {\sffamily\bfseries\raggedright\color{accent}}
  {\fontsize{42}{42}\selectfont\thechapter\hspace{14pt}}
  {0pt}
  {\huge\color{ink}}
  [\vspace{-2pt}{\color{accent}\rule{\textwidth}{1.5pt}}]
\titlespacing*{\chapter}{0pt}{-14pt}{22pt}

\titleformat{\section}{\sffamily\bfseries\large\color{accent}}{\thesection}{0.7em}{}
\titleformat{\subsection}{\sffamily\bfseries\normalsize\color{ink}}{\thesubsection}{0.6em}{}
\titleformat{\subsubsection}{\sffamily\bfseries\small\color{ink!85}}{\thesubsubsection}{0.5em}{}

\setlength{\headheight}{23pt}
\pagestyle{fancy}
\fancyhf{}
\fancyhead[L]{\sffamily\footnotesize\color{gray!80} AI 量化炒股}
\fancyhead[R]{\sffamily\footnotesize\color{gray!80}\nouppercase{\leftmark}}
\fancyfoot[C]{\sffamily\footnotesize\color{gray!70}\thepage}
\renewcommand{\headrulewidth}{0.3pt}
\renewcommand{\headrule}{\hbox to\headwidth{\color{gray!35}\leaders\hrule height \headrulewidth\hfill}}

% ============================ 盒子 ============================
\newtcolorbox{keybox}[1][]{
  enhanced, breakable, colback=soft, colframe=accent,
  boxrule=0pt, leftrule=3.2pt, arc=2pt, left=10pt, right=10pt, top=7pt, bottom=7pt,
  fonttitle=\sffamily\bfseries\small\color{accent}, coltitle=accent,
  attach boxed title to top left={xshift=6pt,yshift=-2pt},
  boxed title style={colback=white,colframe=white,boxrule=0pt},
  title={#1}}

\newtcolorbox{warnbox}[1][]{
  enhanced, breakable, colback=warms, colframe=warn,
  boxrule=0pt, leftrule=3.2pt, arc=2pt, left=10pt, right=10pt, top=7pt, bottom=7pt,
  fonttitle=\sffamily\bfseries\small\color{warn}, coltitle=warn,
  attach boxed title to top left={xshift=6pt,yshift=-2pt},
  boxed title style={colback=warms,colframe=warms,boxrule=0pt},
  title={#1}}

\lstnewenvironment{pycode}[1][]{%
  \lstset{language=Python, basicstyle=\ttfamily\scriptsize,
    breaklines=true, columns=fullflexible, keepspaces=true,
    frame=single, framerule=0.4pt, rulecolor=\color{gray!45},
    backgroundcolor=\color{gray!3}, xleftmargin=4pt, xrightmargin=4pt,
    keywordstyle=\color{accent}\bfseries, commentstyle=\color{gray!75},
    stringstyle=\color{good}, showstringspaces=false, tabsize=2,
    title={\sffamily\bfseries\scriptsize\color{gray!30!black} #1}}%
}{}

\newtcolorbox{plainbox}[1][]{
  enhanced, breakable, colback=white, colframe=gray!40,
  boxrule=0.5pt, arc=1pt, left=10pt, right=10pt, top=7pt, bottom=7pt,
  fonttitle=\sffamily\bfseries\small\color{gray!30!black}, title={#1}}

\setlength{\parskip}{0.32em}
\renewcommand{\arraystretch}{1.2}
\setlist{itemsep=1pt,topsep=3pt,parsep=2pt}

\newtheorem{proposition}{命题}[chapter]
\newtheorem{definition}{定义}[chapter]
\newtheorem{example}{例}[chapter]

\newcommand{\SR}{\mathrm{SR}}
\newcommand{\IC}{\mathrm{IC}}
\newcommand{\Var}{\mathrm{Var}}
\newcommand{\E}{\mathbb{E}}
\newcommand{\N}{\mathcal{N}}
\newcommand{\ph}[1]{\textcolor{gray!60}{#1}}

% ============================ 封面 ============================
\newcommand{\AuditDSRStress}{0.8220}
\newcommand{\AuditDSRHundred}{0.9847}
\newcommand{\AuditPBO}{0.107}
\newcommand{\AuditCommonN}{402}
\newcommand{\AuditCorrMin}{0.716}
\newcommand{\AuditCorrMax}{0.994}
\newcommand{\AuditFlowMfr}{-0.0108}
\newcommand{\AuditFlowPos}{+0.0246}
\newcommand{\AuditFlowPosT}{5.41}
\newcommand{\AuditQCfoN}{371}
\newcommand{\AuditQCfoCount}{1609}
\newcommand{\AuditQCfoT}{3.27}
\newcommand{\AuditQCfoNaive}{6.82}
\newcommand{\AuditBPLower}{0.71}
\newcommand{\AuditBPUpper}{2.27}

\begin{document}
\frontmatter
\thispagestyle{empty}
\begin{center}
\vspace*{3.2cm}
{\sffamily\color{accent}\fontsize{13}{16}\selectfont 量化研究工作手册 · 第 1 版}

\vspace{1.4cm}
{\sffamily\bfseries\color{ink}\fontsize{34}{42}\selectfont AI 量化炒股}

\vspace{6pt}
{\sffamily\bfseries\color{ink}\fontsize{22}{30}\selectfont 原理、实现与过拟合诊断}

\vspace{1.6cm}
{\color{accent}\rule{0.62\textwidth}{1.2pt}}

\vspace{0.9cm}
{\large\color{ink!85}
从市场可预测性的来源，到可复现的因子库、\\[4pt]
本地数据库与多轮重复验证的完整链路}

\vfill
{\small\color{gray}
数据区间：2010--2026 年 · 全 A 股日频\\
外部意见交叉核验：Grok · Gemini（2026-10-31）\\
主要计算可用配套代码复现；未解决的审计项随报告披露}
\vspace{2.2cm}
\end{center}
\clearpage

% ============================ 前言 ============================
\chapter*{写在前面：这本书想解决什么问题}
\addcontentsline{toc}{chapter}{写在前面}
\markboth{写在前面}{写在前面}

关于"用 AI 炒股"，市面上有两类内容。一类是营销：把一条漂亮的净值曲线当作能力的证明，
绝口不提这条曲线是从多少次尝试里挑出来的。另一类是学术：公式严谨，但默认你有
Wind 终端和一台 64 核的服务器，落到个人电脑上寸步难行。

这本书想走第三条路：\textbf{在个人电脑上，用免费数据，做一次经得起自己反复质疑的研究}。
它有三条硬约束，贯穿全书：

\begin{enumerate}
  \item \textbf{可复现}。每一个数字都对应仓库里的一段代码，换台机器、换个时间重跑，
        结论应当一样（或至少差异可解释）。
  \item \textbf{可证伪}。任何"这个因子有效"的说法，都必须给出它在什么条件下会失效，
        以及用什么统计量判断它已经失效。
  \item \textbf{可执行}。不考虑涨跌停不可成交、不考虑 T+1、不考虑 100 股整手、
        不考虑印花税的结论，不叫结论，叫练习。
\end{enumerate}

\section*{你会得到什么}

全书共 13 章与 7 个附录。配套仓库（\texttt{output/ai-quant-2026/}）包含四件可以独立使用的东西：

\begin{center}
\begin{tabular}{@{}lll@{}}
\toprule
交付物 & 位置 & 说明 \\
\midrule
本地数据库 & \texttt{data/quant.db} & 全 A 股日线、PIT 财报、估值快照、指数 \\
因子库 & \texttt{aq/factors.py} & 39 个主因子，带公式、方向、回看期与注册表 \\
研究框架 & \texttt{aq/*.py} & 面板、中性化、评价、CV、DSR/PBO、回测 \\
拟合精度报告 & \texttt{reports/} & 多协议重复测试、过拟合诊断、代码清单 \\
\bottomrule
\end{tabular}
\end{center}

\section*{怎么读}

如果你只有两小时：读第 1、9、10、11 章，你会知道"AI 炒股"这个词里哪部分是统计学、
哪部分是叙事，以及一份诚实的实证报告长什么样。
如果你要动手：按第~\ref{ch:roadmap} 章的路线图走，每完成一步都有明确的"可以继续"判据。
如果你在读别人的研报：直接翻到第 10 章，用那里的四个统计量给对方的结论打折。

\begin{keybox}[关于本书的立场]
本书不提供任何投资建议，也不承诺任何策略能赚钱。它提供的是一套
\textbf{判断"某个策略结论是否可信"的方法}。在这个行业里，
会算夏普比率的人很多，会在算完之后给自己扣分的人很少。
\end{keybox}

\clearpage
\tableofcontents
\clearpage

\mainmatter
\chapter{可预测性从哪里来}\label{ch:source}

\begin{keybox}[本章要回答的问题]
股票价格里到底有没有可被系统利用的规律？如果有，它来自哪里、能持续多久、
为什么没有被套利干净？以及——机器学习在其中究竟贡献了什么？
\end{keybox}

\section{一个思想实验：随机也能做出漂亮的回测}

在讲任何原理之前，先做一个实验。假设我们对 500 只股票的未来收益完全没有任何信息，
但做 1000 次独立的"策略研发"：每次随机给股票打分、按分数排序、取前 20\% 做多、
后 20\% 做空，回测五年。

结果会令你吃惊：这 1000 条净值曲线里，最好的一条年化收益可以轻松超过 20\%，
夏普比率超过 1.5。而且它看起来"很稳定"——月度胜率 60\% 以上，最大回撤不到 10\%。

这不是数据挖掘的偶然，而是极值分布的必然。若 1000 条纯噪声策略的夏普比率近似服从
$\N(0, \sigma^2)$，那么最大值的期望约为

\begin{equation}
  \E\!\left[\max_{1\le n\le N} \SR_n\right]
  \;\approx\; \sigma\left[(1-\gamma)\Phi^{-1}\!\left(1-\tfrac1N\right)
  + \gamma\,\Phi^{-1}\!\left(1-\tfrac{1}{Ne}\right)\right],
  \qquad \gamma \approx 0.5772,
  \label{eq:emax}
\end{equation}

其中 $\Phi^{-1}$ 是标准正态分位数。$N=1000$、$\sigma=1/\sqrt{T}$（$T$ 为样本天数）时，
即便真实夏普为 $0$，观测到的最大值也会显著大于零。\textbf{你看到的"最好策略"，
本质上是一个极值统计量，而不是一个能力证明。}

\begin{warnbox}[本书的第一个纪律]
任何回测结果，在知道"它是在多少次尝试中被选出来的"之前，都不构成证据。
第~\ref{ch:overfit} 章会把这个直觉变成可计算的检验（Deflated Sharpe Ratio 与 PBO）。
\end{warnbox}

\section{可预测性的五类真实来源}

抛开叙事，股价的可预测性只有有限的几个来源。理解它们，你才能判断一个"新因子"
是新发现，还是老故事的换皮。

\subsection{风险补偿：你赚的是承担风险的报酬}

有效市场理论并不否认"超额收益"，它否认的是"无风险的超额收益"。价值股长期跑赢，
可能只是因为它们在经济下行时跌得更惨，投资者要求补偿。这条逻辑链上有几类经典异象：

\begin{itemize}
  \item \textbf{价值}：账面市值比（BP）、盈利收益率（EP）高的股票，长期收益更高。
  \item \textbf{低波动}：低波动股票的风险调整后收益反而更高（杠杆约束与彩票偏好共同导致）。
  \item \textbf{质量}：高 ROE、低应计、现金流扎实的公司更抗跌。
  \item \textbf{流动性}：流动性差的股票需要更高的预期收益来补偿变现成本。
\end{itemize}

\textbf{关键点}：风险补偿型因子的收益是"周期性"的。它们在危机中会大幅回撤，
而且当越来越多资金涌向同一类股票时，补偿会被压缩——这就是"因子拥挤"。

\subsection{行为偏差：你赚的是别人犯错的代价}

A 股散户成交占比高，行为偏差的痕迹比成熟市场更明显：

\begin{itemize}
  \item \textbf{过度反应与反转}：短期（1--20 日）涨幅过大的股票，后续倾向回落。
        这是 A 股最强的横截面效应之一，本书实测的 5 日反转 Rank IC 在 $0.05$ 量级。
  \item \textbf{处置效应}：投资者倾向卖出盈利股、持有亏损股，造成动量在中长期偏弱。
  \item \textbf{彩票偏好}：正偏度（偶尔暴涨）的股票被系统性高估，
        因此"最大单日收益"是一个负向因子。
  \item \textbf{有限注意}：换手率骤升、登上龙虎榜的股票，短期被过度关注，随后回落。
\end{itemize}

\subsection{微观结构：交易规则本身制造了规律}

这一类最容易被忽视，却最"可执行"，因为它直接来自规则：

\begin{itemize}
  \item \textbf{T+1}：当日买入不能当日卖出，日内动量策略在 A 股基本不可行，
        但隔夜与次日开盘的行为因此更有规律。
  \item \textbf{涨跌停}：涨停日买不进、跌停日卖不出。任何假设"信号出现即可成交"的回测，
        在高换手策略上会系统性高估收益。
  \item \textbf{最小报价单位与整手}：低价股、高价股的实际可交易性与理论权重不同。
  \item \textbf{流动性摩擦}：Amihud 非流动性、买卖价差，直接影响可实现收益。
\end{itemize}

\subsection{信息扩散：你知道得比别人早，或者理解得比别人深}

财报公布后的漂移（PEAD）、分析师预期修正、公告文本的情绪变化，都属于这一类。
它的前提是\textbf{严格的时间点对齐}：你必须在信息真正公开的那一刻之后才能使用它。
用报告期（而不是公告日）对齐财务数据，是这一族因子里最常见的致命错误。

\subsection{拥挤与衰减：alpha 会自己消失}

任何被公开、被产品化、被资金追逐的规律，都会因为交易而衰减。这不是理论推测：

\begin{itemize}
  \item 小市值因子在 2017 年之前是 A 股最强的因子之一，2017 年之后随着监管与
        注册制扩容大幅削弱；
  \item 经典的中长期动量在 A 股长期偏弱甚至为负，而在美股相对稳定；
  \item 高频量价因子在 2021 年量化规模快速扩张后被显著压缩。
\end{itemize}

\begin{keybox}[判断一个因子能不能用的第一问]
它的收益来自\textbf{承担了什么风险}、\textbf{纠正了谁的错误}、还是\textbf{利用了哪条规则}？
如果三个都答不上来，那它大概率只是你在历史上挖到的一段噪声。
\end{keybox}

\section{机器学习到底贡献了什么}

这是全书最需要说清楚的一点。很多人以为机器学习能"发现人看不见的规律"，
但在日频横截面选股这个问题上，它的真实贡献要朴素得多：

\begin{enumerate}
  \item \textbf{非线性组合}。把 40 个因子线性加权，等于假设"每个因子的边际贡献恒定"。
        GBDT 能学到"低波因子在小市值股票里更有效"这类交互，从而提升一点 IC。
  \item \textbf{自动特征工程}。树模型能自动构造阈值与交叉项，省去人工试错。
  \item \textbf{对噪声的鲁棒性}。正则化与集成让模型不至于被单个极端样本带偏。
\end{enumerate}

它\textbf{不能}做的事同样清楚：

\begin{itemize}
  \item 它不能凭空创造新的经济来源。如果 40 个因子本身没有信息，模型也变不出来。
  \item 它不能对抗标签的噪声。日频收益的信噪比极低，模型很容易学到噪声。
  \item 它不能免除验证的责任。恰恰相反，模型越复杂，过拟合的自由度越大，
        需要的验证越严格。
\end{itemize}

\begin{center}
\begin{tabular}{@{}p{0.44\textwidth}p{0.46\textwidth}@{}}
\toprule
\textbf{机器学习真正擅长的} & \textbf{它做不到的} \\
\midrule
特征交互与非线性 & 无中生有的 alpha \\
多因子加权与降噪 & 替代严格的时间点对齐 \\
高频数据的模式识别 & 消除交易成本 \\
大规模横截面的批量预测 & 保证样本外有效 \\
\bottomrule
\end{tabular}
\end{center}

业界常引用的一个经验量级：在严格执行样本外验证之后，日频全 A 股的
模型 Rank IC 大致落在 $0.03$--$0.08$ 区间，年化 ICIR 在 $1$--$4$ 之间。
超过 $0.15$ 的日频 Rank IC，你首先应该怀疑的是数据泄漏，而不是庆祝。

\section{A 股的结构性特征}

把美股的那套方法直接搬到 A 股，会在四个地方翻车：

\subsection{T+1 与涨跌停}

T+1 意味着"当日买入不能当日卖出"，因此：
\begin{itemize}
  \item 标签不能用"当日收盘买、当日收盘卖"构造；
  \item 组合的换手约束更强，日频调仓的实际成本高于理论值。
\end{itemize}

涨跌停（主板 $\pm10\%$、创业板与科创板 $\pm20\%$、北交所 $\pm30\%$、ST $\pm5\%$）
意味着：
\begin{itemize}
  \item 涨停当天大概率买不进（尤其一字板），跌停当天卖不出；
  \item 如果不模拟这一点，反转类策略的回测收益会被系统性高估——
        因为"跌得最惨的股票"往往正是跌停卖不出的那些。
\end{itemize}

\subsection{散户结构与情绪波动}

A 股个人投资者成交占比长期偏高，导致：
\begin{itemize}
  \item 短期反转、换手率、情绪类因子比成熟市场更有效；
  \item 因子收益的波动更大，尾部风险更厚。
\end{itemize}

\subsection{注册制扩容与退市常态化}

2019 年科创板、2020 年创业板注册制、2021 年北交所，使得：
\begin{itemize}
  \item 股票数量从 2010 年的约 2000 只增长到 2026 年的约 5600 只；
  \item 小市值不再稀缺，"壳价值"逻辑失效；
  \item 退市股增多，\textbf{幸存者偏差}的杀伤力比过去更大。
\end{itemize}

\subsection{量化监管与交易约束}

2023 年之后对程序化交易的监管趋严，高频策略的容量被压缩。
这对个人研究者的启示是：\textbf{不要把研究建立在"极高换手"的假设上}，
那条路在 A 股的制度环境里越来越窄。

\section{哪些信号还在，哪些已经衰减}

把上文压缩成一张可执行的判断表。这是本书后续所有实证工作的先验：

\begin{center}
\begin{longtable}{@{}p{0.20\textwidth}p{0.13\textwidth}p{0.55\textwidth}@{}}
\toprule
\textbf{因子族} & \textbf{当前状态} & \textbf{判断依据} \\
\midrule
\endhead
短期反转（1--20 日） & 仍有效 & A 股最强的横截面效应，散户过度反应的直接结果 \\
低波动 / 特质波动 & 仍有效 & 杠杆约束 + 彩票偏好，跨市场稳健 \\
换手率 / 流动性 & 仍有效 & 拥挤度与关注度的代理变量 \\
价值（EP/BP） & 中等 & 大中盘较稳，小盘受"价值陷阱"干扰 \\
质量（ROE/现金流） & 中等 & 需要配合估值使用，单独使用收益有限 \\
规模（小市值） & 显著衰减 & 2017 年后壳价值消失，注册制进一步压缩 \\
中长期动量 & 弱甚至为负 & A 股中长期动量长期不稳定 \\
高频量价 & 衰减中 & 量化规模扩张 + 监管趋严 \\
\bottomrule
\end{longtable}
\end{center}

\begin{warnbox}[方向不是常识，是实证问题]
很多教程会写"动量因子：过去涨得多的股票未来继续涨"。在 A 股日频上，
这句话很可能是反的。本书第~\ref{ch:factor} 章的实测会显示：20 日动量的 IC 为负，
即\textbf{反转主导}。任何关于因子方向的断言，都必须在自己的样本上重新验证一次。
\end{warnbox}

\section{本章小结}

\begin{enumerate}
  \item 漂亮的回测曲线是极值统计量的自然产物，没有"尝试次数"信息的夏普比率不构成证据。
  \item 可预测性只有五类来源：风险补偿、行为偏差、微观结构、信息扩散、拥挤衰减。
        说不出来源的因子，大概率是噪声。
  \item 机器学习的贡献是"更好的组合与降噪"，不是"发现新规律"。
        它同时放大了过拟合的风险。
  \item A 股的 T+1、涨跌停、散户结构、注册制扩容，决定了哪些方法可行、哪些不可行。
  \item 因子方向必须实证，不能照搬教科书。
\end{enumerate}

\chapter{数据：一切结论的地基}\label{ch:data}

\begin{keybox}[本章要回答的问题]
免费能不能拿到可用的 A 股数据？复权到底怎么算？Point-in-Time 为什么是生死线？
一个能长期维护的本地数据库应该长什么样？
\end{keybox}

\section{一次真实的数据源排雷记录}

2026 年 10 月 31 日，本项目在真实网络环境下依次尝试了五个数据源，过程值得完整记录，
因为它比任何"数据源对比表"都更接近你将来会遇到的情况。

\begin{center}
\begin{longtable}{@{}p{0.16\textwidth}p{0.30\textwidth}p{0.42\textwidth}@{}}
\toprule
\textbf{源} & \textbf{实测结果} & \textbf{失败方式与教训} \\
\midrule
\endhead
东方财富 \texttt{push2his} & 前 50 次请求正常，之后全部失败 & \texttt{RemoteDisconnected}。
IP 级封锁，直连与系统代理同时失效。\textbf{教训：批量抓取必须在几分钟内完成，
或者必须限速。} \\
\midrule
腾讯 \texttt{fqkline} & 前约 200 次请求正常，之后全部 501 & 返回
\texttt{waf.tencent.com} 的挑战页。\textbf{教训：501 不一定代表"接口不存在"，
可能是 WAF。} \\
\midrule
baostock & 单只 3--28 秒，并发一高即被限速 & TCP 长连接协议，字段最全
（含 \texttt{isST}、\texttt{tradeStatus}），但不适合大批量。 \\
\midrule
新浪 \texttt{CN\_MarketData} & 稳定，单只 0.3--0.9 秒 & \texttt{datalen} 参数可放大到
$10^4$，一次返回全历史；配套 \texttt{hfq.js} 给出累计后复权因子。\textbf{本项目最终选它。} \\
\midrule
同花顺 \texttt{v6/line} & 稳定，单只约 1 秒返回全部 6416 根 & 但数据是差分编码的私有格式，
需要逆向解码，工程成本高，未采用。 \\
\bottomrule
\end{longtable}
\end{center}

\begin{warnbox}[这段经历本身就是结论]
\textbf{没有任何单一免费数据源是可靠的。}生产级做法只有三条：
\begin{enumerate}
  \item 写\textbf{适配器层}，把"取数据"和"用数据"解耦，换源只改一个文件；
  \item 写\textbf{断点续抓}，每只股票的成功与否记进日志表，失败可以单独重跑；
  \item 写\textbf{交叉校验}，用第二个源抽查 1\% 的样本，核对收盘价与复权因子。
\end{enumerate}
\end{warnbox}

\section{复权：三种口径与一个必须记住的结论}

假设某股票在除权前收盘 10 元，实施"10 送 10"，除权后理论开盘价 5 元。
如果不做任何处理，行情数据上会出现一根 $-50\%$ 的"暴跌"，而你的动量因子会认为
这只股票"跌惨了，应该反弹"——完全错误。

\subsection{三种复权方式}

设第 $t$ 日的复权因子为 $f_t$，原始价格为 $P_t$。

\begin{description}
  \item[后复权（hfq）] $\tilde P_t = P_t \cdot f_t$，且 $f$ 单调不减，最新一日 $f$ 最大。
        历史价格不变，新数据追加时旧值不动。
  \item[前复权（qfq）] 以上市首日为基准，最新价等于真实价，历史价被等比缩小。
        \textbf{每次除权，全部历史价格都会变}。
  \item[不复权（bfq）] 就是原始价格，除权日出现跳空。
\end{description}

\subsection{为什么回测必须用后复权}

\begin{enumerate}
  \item \textbf{前复权会破坏历史}。你今天下载的前复权序列，明天公司分红后整条序列都会变。
        如果因子值依赖历史价格，你的因子值也会变——研究结果无法复现。
  \item \textbf{不复权会把分红送转当成暴跌}。所有基于收益率的因子（反转、动量、波动）
        都会在除权日附近产生极端值。
  \item \textbf{后复权是唯一"只增不改"的口径}。追加新数据不会修改旧值。
\end{enumerate}

\begin{keybox}[库表设计的对应做法]
数据库里\textbf{同时存}不复权价格与累计复权因子 $f_t$，任何时刻都能还原
$\tilde P_t = P_t f_t$。这样既保留了"真实成交价"（算市值、算涨跌停、算 PE 都需要它），
又保留了"总收益序列"（算因子、算净值需要它）。
\end{keybox}

\subsection{复权因子的构造}

后复权因子在每次除权除息时跳变。若某次分红为每股 $D$ 元、送转为每股 $S$ 股，
除权前收盘 $P_{t^-}$，则

\begin{equation}
  f_{t} = f_{t^-}\cdot\frac{P_{t^-}}{P_{t^-} - D + S\cdot P_{t^-}/(1+S)}
  \;\;\approx\;\; f_{t^-}\cdot\frac{1}{1 - \frac{D}{P_{t^-}}},
  \label{eq:adjfactor}
\end{equation}

工程上更稳妥的做法是直接使用数据源给出的累计因子序列（新浪 \texttt{hfq.js}、
baostock 的 \texttt{backAdjustFactor}），然后用"相邻两日总收益是否连续"来校验：
除权日附近不应出现超过涨跌停幅度的收益率跳变。本书的校验脚本
\texttt{scripts/check\_db.py} 会检查复权因子是否单调。

\section{Point-in-Time：最隐蔽的前视偏差}

假设某公司 2024 年一季报显示净利润同比增长 $300\%$，报告期是 2024-03-31，
但公告日是 2024-04-28。

错误做法：在 2024-03-31 就用这个数据算因子。
后果：你的模型在 3 月 31 日就知道了 4 月 28 日才公开的消息，样本外表现会好得离谱，
而这个"优势"在实盘中完全不存在。

正确做法：用 \texttt{notice\_date} 做 as-of 连接，即对每个交易日 $t$，
只能使用 $\{\text{notice\_date} \le t\}$ 的最近一期财报。

\begin{pycode}[PIT as-of 连接]
df = pd.merge_asof(
    df.sort_values(["date", "code"]),          # 行 = 交易日
    fin.sort_values(["notice_date", "code"]),  # 行 = 财报公告
    left_on="date", right_on="notice_date",
    by="code", direction="backward",           # 只往回找，绝不往前
)
\end{pycode}

\begin{warnbox}[三个容易踩的变体]
\begin{enumerate}
  \item \textbf{用报告期代替公告日}——最常见的错误，导致 1--4 个月的前视。
  \item \textbf{用"最新一期"数据回填全部历史}——把今天才知道的财报喂给 2015 年的模型。
  \item \textbf{用修正后的财务数据}——财报会重述，正确的做法是用当时的原始值。
        免费数据源通常只提供最新值，这是必须承认的局限。
\end{enumerate}
\end{warnbox}

\section{幸存者偏差：只在活着的股票上做研究}

如果你今天下载"全部 A 股"，你拿到的是\textbf{今天还活着}的 5600 只。
2010 年以来退市的几百只股票不在里面。而退市的股票往往是表现最差的那些。

后果：任何"买入低估值股票"的策略，回测收益会被系统性高估，
因为你永远买不到后来退市的那批。

\subsection{怎么缓解}

\begin{itemize}
  \item 尽量使用\textbf{含退市股}的数据源（baostock 覆盖较好，东财数据中心的历史财报也能还原已退市公司）；
  \item 用\textbf{时点股票池}：每一天只用"当天真实存在且可交易"的股票，
        而不是"今天的股票池"回测历史；
  \item 明确记录局限：如果数据源确实缺退市股，就要在报告里写清楚，
        并对结论打一个折扣。
\end{itemize}

本项目的数据构建方式（用历年财报期反推存续公司）天然包含已退市公司，
但行情接口对已退市股票的覆盖并不完整，这一点在第~\ref{ch:result} 章的数据质量表中如实披露。

\section{可交易性：把"买不到"变成显式变量}

很多回测的隐含假设是"信号出现就能按收盘价成交"。在 A 股，这个假设在四种情况下不成立：

\begin{center}
\begin{tabular}{@{}lll@{}}
\toprule
\textbf{情形} & \textbf{判定方式} & \textbf{回测处理} \\
\midrule
停牌 & 当日无成交（K 线缺行） & 不可买、不可卖 \\
涨停 & $\text{收盘} \ge \text{涨停价} - \epsilon$ & 当日不可买入 \\
跌停 & $\text{收盘} \le \text{跌停价} + \epsilon$ & 当日不可卖出 \\
一字板 & $\text{最高} = \text{最低}$ & 全天无法成交 \\
ST & 简称含 \texttt{ST} & 剔除（或单独建模） \\
新股 & 上市未满 $N$ 日 & 剔除 \\
流动性不足 & 20 日均成交额 $< \theta$ & 剔除 \\
\bottomrule
\end{tabular}
\end{center}

涨停价的判定必须按板块区分：

\begin{pycode}[涨跌停幅度]
def limit_pct(code, is_st, date):
    chinext = code.startswith(("300", "301"))   # 创业板
    star    = code.startswith(("688", "689"))   # 科创板
    bse     = code.startswith(("4", "8", "92")) # 北交所
    pct = np.where(chinext | star, 0.20, 0.10)
    pct = np.where(bse, 0.30, pct)
    pct = np.where(is_st & ~(chinext | star | bse), 0.05, pct)
    early = (date < "2020-08-24") & chinext       # 注册制改革前为 10%
    pct = np.where(early, np.where(is_st, 0.05, 0.10), pct)
    return pct
\end{pycode}

\begin{warnbox}[ST 状态必须用历史，不能用现在]
用"当前简称含 ST"去筛历史，会把 2015 年正常交易、2020 年才戴帽的公司错误剔除，
也会把早已摘帽的公司错误保留。正确做法是构造\textbf{简称时间线}：
每期财报的披露简称 + 公告日，as-of 连接回每个交易日。
\end{warnbox}

\section{数据库设计}

\subsection{表结构}

\begin{center}
\begin{longtable}{@{}p{0.18\textwidth}p{0.16\textwidth}p{0.54\textwidth}@{}}
\toprule
\textbf{表} & \textbf{主键} & \textbf{内容与设计意图} \\
\midrule
\endhead
\texttt{stock} & code & 证券主表：交易所、行业、上市/退市日 \\
\texttt{bar} & (code, date) & 日线：不复权 OHLC + 复权因子 + 成交量/额 + 停牌标记 \\
\texttt{fin} & (code, report\_date) & 财报：含 \texttt{notice\_date}（PIT 关键）、简称历史 \\
\texttt{valuation} & (code, date) & 月度股本/市值/PE/PB 快照 \\
\texttt{index\_bar} & (code, date) & 指数行情（基准与归因） \\
\texttt{trade\_calendar} & date & 交易日历（由指数日期生成） \\
\texttt{fetch\_log} & (task, key) & 抓取日志，支持断点续抓 \\
\bottomrule
\end{longtable}
\end{center}

\subsection{一条设计纪律：派生量不进库}

"ST 状态""换手率""净利润 TTM"这三列，看起来应该存在 \texttt{bar} 表里，
但本书坚持\textbf{不在库里硬写}，而是在构建面板时实时生成。理由有三条：

\begin{enumerate}
  \item \textbf{可重算}。口径一变（比如 ST 判定规则改了），历史值能追溯重算，
        而不是被固化在几千万行里。
  \item \textbf{可维护}。每次更新财务/估值都要全表重算几千万行，代价太高。
  \item \textbf{防误用}。一旦把"当前简称"写进历史行，前视偏差就有了滋生的土壤。
\end{enumerate}

\begin{keybox}[数据库的分层原则]
\textbf{原始事实层}（行情、财报、估值）只放数据源给的东西，一行都不加工；
\textbf{派生特征层}（面板、因子、标签）每次研究时重新计算，可以随时改口径。
两者之间用 \texttt{date} 与 \texttt{code} 对齐，靠 as-of 连接保证时点正确。
\end{keybox}

\section{数据质量检查清单}

一个负责任的研究，在跑任何因子之前，应该先回答下面这些问题：

\begin{enumerate}
  \item 每个交易日的股票数量是否合理？是否有整段缺失？
  \item 复权因子是否单调不减？除权日附近的收益率是否超过涨跌停幅度？
  \item 财务数据的公告日是否都晚于报告期？
  \item 可交易股票占比是否在合理区间（一般 70\%--90\%）？
  \item 涨跌停天数占比是否与市场常识相符（极端行情年份会显著上升）？
  \item 换手率覆盖率是多少？缺失的股票是不是都是已退市/长期停牌？
  \item 随机抽查 20 只股票，用第二个数据源核对收盘价与复权因子。
\end{enumerate}

对应的实现是 \texttt{scripts/check\_db.py}，它输出的表会直接放进第~\ref{ch:result}
章的报告里。\textbf{不先看这张表就开始挖因子，等于在流沙上盖楼}。

\section{本章小结}

\begin{enumerate}
  \item 免费数据源的可靠期以"分钟到小时"计，必须写适配器 + 断点续抓 + 交叉校验。
  \item 回测只用后复权；库里同时保留不复权价与复权因子，两者各司其职。
  \item PIT 是生死线：用公告日，不用报告期；用当时的原始值，不用修正值。
  \item 幸存者偏差靠"时点股票池"缓解，无法完全消除时要如实披露。
  \item 可交易性必须变成显式布尔变量，而不是隐含假设。
  \item 原始事实与派生特征分层存储，是长期可维护的关键。
\end{enumerate}

\chapter{标签：你预测什么，决定你能得到什么}\label{ch:label}

\begin{keybox}[本章要回答的问题]
为什么"预测明天涨跌"是一个坏问题？横截面排序、波动率归一化、三重障碍法各自
适合什么场景？重叠标签为什么会让你的 t 值虚高？A 股的 T+1 与涨跌停如何改变标签设计？
\end{keybox}

\section{一个被低估的事实：多数模型失败在标签，而不是模型}

研究新手最常见的路径是：拿到数据 $\rightarrow$ 构造几十个因子 $\rightarrow$ 上 LightGBM $\rightarrow$ 调参 $\rightarrow$
发现样本外不灵 $\rightarrow$ 换更复杂的模型。

他们很少检查一件事：\textbf{标签本身是否可预测}。

日频个股收益的信噪比极低。以本书的数据为例，某只股票某日的收益，
其横截面标准差约为 $2.5\%$；而任何模型能解释的部分，通常不到这个数字的 $5\%$。
在这种信噪比下，标签的设计方式（预测什么、怎么归一化、如何处理重叠）
对最终结果的影响，往往大于模型选择。

\section{三种标签形态}

\subsection{回归：预测未来 $h$ 日收益}

\begin{equation}
  y_{i,t} = \frac{\tilde P_{i,t+1+h}}{\tilde P_{i,t+1}} - 1
  \label{eq:label-ret}
\end{equation}

其中 $\tilde P$ 是后复权价，$t+1$ 表示"次日开盘买入"。

\textbf{优点}：保留了收益的幅度信息，与最终的组合权重直接对应。
\textbf{缺点}：厚尾。A 股单日涨跌 $20\%$ 的样本不少，少数极端值会主导平方损失。

\begin{warnbox}[一个常见错误：用收盘价构造标签]
\texttt{fwd\_ret = close[t+h]/close[t] - 1} 隐含了"$t$ 日收盘买入"的假设。
但你的信号是用 $t$ 日收盘数据算出来的，实际上根本来不及在 $t$ 日收盘成交。
正确做法是 \texttt{close[t+h]/open[t+1] - 1}——\textbf{信号与成交之间必须隔一个可执行的时点}。
\end{warnbox}

\subsection{横截面排序：预测相对强弱}

\begin{equation}
  y_{i,t} = \mathrm{rank}_{i}\!\left(r_{i,t+1\rightarrow t+1+h}\right) - \tfrac12
  \;\in\; \left(-\tfrac12,\ \tfrac12\right]
  \label{eq:label-rank}
\end{equation}

\textbf{优点}：
\begin{itemize}
  \item 天然抗极值（排名不受 $+50\%$ 与 $+5\%$ 的差距影响）；
  \item 与选股目标一致——我们只关心"买哪些"，不关心"市场整体涨多少"；
  \item 自动剔除市场 beta，模型学到的纯粹是横截面信息。
\end{itemize}

\textbf{缺点}：丢失幅度信息，无法区分"小幅跑赢"与"大幅跑赢"。

本书的主标签采用这一形式。这是业界的标准选择，原因是：日频收益的幅度部分
几乎全是噪声，而排序部分是相对稳定的。

\subsection{分类：涨 / 跌 / 前 $30\%$}

把问题变成二分类或多分类。

\textbf{优点}：可以直接用分类模型，输出概率有明确含义。
\textbf{缺点}：阈值是人为的；信息损失最大；概率校准需要额外工作。

除非你有明确的"只买前 $30\%$"的业务约束，否则不建议作为主标签，
可以用来做辅助诊断（例如分析模型在极端分位上的表现）。

\section{波动率归一化：让不同股票可比}

同样是 $+3\%$ 的周收益，对一只年化波动 $20\%$ 的银行股是"大涨"，
对一只年化波动 $80\%$ 的题材股是"平常"。

\begin{equation}
  y_{i,t} = \frac{r_{i,t+1\rightarrow t+1+h}}{\sigma_{i,t}\sqrt{h}},
  \qquad \sigma_{i,t} = \mathrm{std}(r_{i,t-59\rightarrow t})
  \label{eq:label-vol}
\end{equation}

除以 $\sigma$ 之后，标签近似服从同分布，模型的残差也更接近同方差。
\textbf{代价}：低波动股票的权重被放大，组合会向低波倾斜；
如果你同时用了低波因子，就会产生双重暴露。

\section{三重障碍法：什么时候值得用}

López de Prado 提出的三重障碍法（Triple Barrier）给每个样本定义三个出口：

\begin{itemize}
  \item 上障碍：$P_{i,t}(1 + u\sigma_i)$，止盈；
  \item 下障碍：$P_{i,t}(1 - d\sigma_i)$，止损；
  \item 垂直障碍：$t + h$ 个交易日，到期平仓。
\end{itemize}

标签取\textbf{先被触及的那个障碍}：上障碍 $+1$，下障碍 $-1$，到期 $0$。

\subsection{它解决了什么}

它把"路径依赖"纳入了标签：一个先跌 $8\%$ 再涨 $5\%$ 的样本，
在固定持有期标签下是 $+5\%$，在三重障碍下是 $-1$（先触发止损）。
如果你的实际交易规则带止损，三重障碍标签与真实盈亏更一致。

\subsection{它的代价}

\begin{enumerate}
  \item \textbf{标签不平衡}：多数样本落在 $0$（到期平仓），分类器会退化为常数预测。
  \item \textbf{参数敏感}：$u,d$ 的选择直接影响标签分布，而你又会用它来调模型——
        这是双重过拟合的入口。
  \item \textbf{重叠更严重}：路径依赖让相邻样本的标签相关性更高，有效样本数进一步下降。
  \item \textbf{计算量大}：需要对每个样本扫描未来 $h$ 天的路径。
\end{enumerate}

\begin{keybox}[本书的判断]
\textbf{日频选股不需要三重障碍法。}
它的设计目标是"交易执行"层面的路径依赖，更适合高频或单标的择时。
日频横截面选股的核心是"排序准不准"，固定持有期 + 横截面排序标签已经足够，
而且统计性质更干净。
\end{keybox}

\section{Meta-labeling：把方向和仓位分开}

Meta-labeling 的思路是两级模型：

\begin{enumerate}
  \item \textbf{主模型}决定方向（买 / 不买），可以是一条简单规则或一个粗模型；
  \item \textbf{副模型}预测"主模型这次判断对不对"，输出一个概率 $p$；
  \item 最终仓位 $\propto (2p - 1)$ 或只保留 $p > \theta$ 的样本。
\end{enumerate}

它的价值在于\textbf{降低假阳性}：主模型给出 100 个买入信号，副模型只保留其中 40 个
把握较大的。这在信号本身很稀疏（比如事件驱动）时特别有效。

在日频全市场选股中，因为每天都有几千个候选，"稀疏"不是问题，
meta-labeling 的边际收益有限，反而增加了复杂度。

\section{重叠标签：为什么你的 $t$ 值虚高}

这是本书最重要的一节之一。

假设你用 $h=5$ 的持有期标签，并且\textbf{每个交易日}都训练一个样本。那么：
\begin{itemize}
  \item $t$ 日的标签覆盖 $[t+1, t+6]$；
  \item $t+1$ 日的标签覆盖 $[t+2, t+7]$；
  \item 两者有 4 天的重叠。
\end{itemize}

后果是相邻样本高度相关。若你用普通 $t$ 检验判断 IC 是否显著，
标准误会严重低估，$t$ 值被夸大 $2$--$4$ 倍。

\begin{proposition}[有效样本数]
若标签持有期为 $h$ 个交易日、采样间隔为 $1$ 天，则\textbf{独立}样本数约为
观测数的 $1/h$，即
\begin{equation}
  n_{\text{eff}} \approx \frac{n}{h}.
\end{equation}
更精细的估计可以用 IC 序列的自相关来修正：
\begin{equation}
  n_{\text{eff}} = \frac{n}{1 + 2\sum_{k=1}^{h-1}\left(1-\frac{k}{h}\right)\rho_k},
\end{equation}
其中 $\rho_k$ 是 IC 序列的 $k$ 阶自相关。
\end{proposition}

工程上有两个正解：

\begin{enumerate}
  \item \textbf{间隔采样}：每 $h$ 个交易日取一个样本，让标签不重叠。
        本书的模型训练采用这一做法（每 5 个交易日一个样本）。
  \item \textbf{Newey-West 修正}：保留全部样本，但在算 $t$ 值时用 HAC 标准误，
        滞后阶数取 $\ge h-1$。本书的因子评价采用这一做法。
\end{enumerate}

\section{A 股的特殊约束}

\subsection{T+1}

当日买入不能当日卖出，因此：
\begin{itemize}
  \item 最短持有期是 1 个交易日（今天买、明天卖）；
  \item 标签的起点必须是 $t+1$（次日），不能是 $t$。
\end{itemize}

\subsection{涨跌停}

如果你的标签是"次日开盘到第 $h$ 日开盘的收益"，那么：
\begin{itemize}
  \item 次日一字涨停买不进 $\rightarrow$ 该样本应当被标记为不可执行；
  \item 到期日一字跌停卖不出 $\rightarrow$ 该样本的收益实际上无法实现。
\end{itemize}

这两种情况都集中在极端样本上，而极端样本恰恰是反转类因子最关注的对象。
不处理它们，反转因子的回测收益会被显著高估。

本书的处理方式：在样本层面保留，但在\textbf{组合构建}时剔除
\texttt{can\_buy\_next = False} 的股票（见第~\ref{ch:portfolio} 章）。

\section{标签与持有期的匹配}

持有期 $h$ 不是一个可以随便设的参数，它由三个因素共同决定：

\begin{enumerate}
  \item \textbf{信号的衰减速度}。反转信号 5 日内衰减大半，20 日基本消失；
        质量因子可以持有 3--6 个月。
  \item \textbf{交易成本}。$h$ 越小，换手越高，成本吃掉的比例越大。
        在单边 $30$ bp 的成本下，$h=1$ 的策略需要极高的 IC 才能覆盖成本。
  \item \textbf{容量}。$h$ 越小，策略容量越小。
\end{enumerate}

\begin{center}
\begin{tabular}{@{}lccc@{}}
\toprule
持有期 $h$ & 年换手（倍） & 单边成本要求 & 适合的信号类型 \\
\midrule
1 日 & 250 & 极高 & 微观结构（个人基本不可行） \\
5 日 & 50 & 高 & 短期反转、量价 \\
20 日 & 12 & 中 & 动量、波动、质量 \\
60 日 & 4 & 低 & 价值、质量、成长 \\
\bottomrule
\end{tabular}
\end{center}

本书选择 $h=5$ 作为主标签：它在"信号强度"与"成本可承受性"之间取得平衡，
并且能把日频数据的重叠问题控制在可管理范围内。

\section{本章小结}

\begin{enumerate}
  \item 标签设计比模型选择更重要，日频收益的信噪比极低。
  \item 横截面排序标签是选股任务的标准选择；回归标签保留幅度但受厚尾干扰。
  \item 波动率归一化让不同股票可比，但会引入低波暴露。
  \item 三重障碍法解决路径依赖，代价是标签不平衡与参数敏感；日频选股不必使用。
  \item 重叠标签会让 $t$ 值虚高，必须用间隔采样或 Newey-West 修正。
  \item A 股的 T+1 与涨跌停改变了标签的可执行性，必须在样本与组合两层处理。
\end{enumerate}

\chapter{因子：从经济假设到可计算表达式}\label{ch:factor}

\begin{keybox}[本章要回答的问题]
什么样的因子值得算？40 个因子背后的经济逻辑分别是什么？为什么"因子越多越好"
是错的？一个可维护的因子库应该如何组织？
\end{keybox}

\section{好因子的四个判据}

在动手算之前，先用四条标准筛一遍：

\begin{enumerate}
  \item \textbf{经济逻辑}：能说清它为什么应该预测收益（第~\ref{ch:source} 章的五类来源）。
  \item \textbf{稳健性}：换一个样本区间、换一个股票池、换一种中性化方式，方向和量级不变。
  \item \textbf{可计算性}：只用公开可得的数据，且能严格做时点对齐。
  \item \textbf{容量}：在可承受的换手下，扣除成本仍有剩余。
\end{enumerate}

\begin{warnbox}["因子越多越好"是最贵的错觉]
每多一个因子，你就多了一次多重检验。挖 500 个因子、留下 20 个，
看起来是"筛选"，实际上是在做 500 次假设检验而不做任何校正。
第~\ref{ch:eval} 章会给出 Benjamini-Hochberg 的完整做法。
本书的因子库只保留 40 个，每个都有明确的经济含义，这不是偷懒，是纪律。
\end{warnbox}

\section{因子库的组织方式}

每个因子登记六件事，缺一不可：

\begin{center}
\begin{tabular}{@{}ll@{}}
\toprule
\textbf{字段} & \textbf{作用} \\
\midrule
\texttt{name} & 唯一名，小写下划线，例如 \texttt{rev\_5} \\
\texttt{category} & 家族，便于做相关性聚类与暴露分析 \\
\texttt{formula} & 人能读懂的公式，必须与代码逐项对应 \\
\texttt{direction} & $+1$ 表示数值越大预期收益越高，$-1$ 反之 \\
\texttt{window} & 最长回看交易日，用于构造 PIT 样本与 purge 区间 \\
\texttt{func} & 接收宽表面板、返回同形状矩阵 \\
\bottomrule
\end{tabular}
\end{center}

\begin{pycode}[因子注册表机制]
@factor("rev_5", "反转", "-(C_t / C_{t-5} - 1)", +1, 7,
        "周度反转：A 股最强的短周期效应之一")
def _rev_5(w: Wide):
    c = w("adj_close")
    return -(c / c.shift(5) - 1)
\end{pycode}

\texttt{direction} 字段看似多余，其实很关键：它把"我的预期"写进了代码。
当你算出的 IC 与 \texttt{direction} 相反时，程序会提醒你——这可能是发现了反例，
也可能是代码写错了。两种都值得停下来看一眼。

\section{反转族}

\subsection{经济逻辑}

散户对短期信息过度反应，推高（或压低）价格后逐步修正。A 股散户成交占比高，
这个效应比成熟市场更明显。

\subsection{核心因子}

\begin{align}
  \text{rev\_1} &= -\left(\frac{C_t}{C_{t-1}} - 1\right), \\
  \text{rev\_5} &= -\left(\frac{C_t}{C_{t-5}} - 1\right), \\
  \text{rev\_10} &= -\left(\frac{C_t}{C_{t-10}} - 1\right).
\end{align}

其中 $C_t$ 为后复权收盘价。负号表示"跌得多的，预期反弹"，因此因子值越大越好。

\subsection{实测与解读}

本书在全 A 股日频上实测，5 日反转的 Rank IC 约为 $+0.05$，年化 ICIR 在 $3$--$7$ 之间，
是所有因子里最稳定的一个。但它的换手极高（单边约 $35\%$／5 日），
扣掉成本后收益会大幅缩水。

\begin{warnbox}[反转因子的两个陷阱]
\begin{enumerate}
  \item \textbf{涨停/跌停不可成交}。反转因子选出的往往是跌幅最大的股票，
        其中相当一部分是跌停板——你根本买不到，或者买到了也卖不掉。
  \item \textbf{ST 与退市}。跌得最惨的股票里混着大量 ST 与准退市股。
        不剔除它们，反转因子的收益会被显著高估。
\end{enumerate}
\end{warnbox}

\section{动量族}

\subsection{经济逻辑}

信息扩散缓慢，价格对新信息反应不足，因此过去涨的继续涨。这个逻辑在美股中长期成立，
在 A 股日频上\textbf{往往反向}——因为短期反转的强度压过了中期动量。

\subsection{核心因子}

\begin{align}
  \text{mom\_20} &= \frac{C_t}{C_{t-20}} - 1, \\
  \text{mom\_60} &= \frac{C_t}{C_{t-60}} - 1, \\
  \text{mom\_12\_1} &= \frac{C_{t-20}}{C_{t-240}} - 1, \\
  \text{resid\_mom\_60} &= \text{mom\_60} - \beta_{120}\cdot \text{mom}^{\text{mkt}}_{60}.
\end{align}

\texttt{mom\_12\_1} 跳过最近一个月，是为了避开短期反转的污染——这是动量研究的
标准做法（Jegadeesh \& Titman 的 12-1 动量）。

\texttt{resid\_mom\_60} 剥离市场 beta，只保留个股特质动量，比裸动量稳定。

\subsection{实测：一个必须面对的结论}

本书实测 20 日动量的 Rank IC 为\textbf{负}（约 $-0.05$ 到 $-0.08$）。
也就是说，在 A 股日频上，"过去 20 天涨得多的股票"后续表现更差。

这不是本书的独特发现，而是 A 股长期存在的现象。它提醒我们：
\textbf{任何因子的方向都必须在自己的样本上重新验证}，教科书的符号不可直接照搬。

\section{波动族}

\subsection{经济逻辑}

低波动异象有两个解释：一是杠杆约束（不能加杠杆的投资者只能通过买高波动股来追求高收益，
推高了高波动股的估值）；二是彩票偏好（投资者愿意为"可能暴涨"付溢价）。

\subsection{核心因子}

\begin{align}
  \text{vol\_20} &= -\mathrm{std}(r, 20), \\
  \text{vol\_60} &= -\mathrm{std}(r, 60), \\
  \text{ivol\_60} &= -\mathrm{std}(r - \beta r^{\text{mkt}}, 60), \\
  \text{skew\_20} &= -\mathrm{skew}(r, 20), \\
  \text{max\_ret\_20} &= -\max(r, 20), \\
  \text{amp\_20} &= -\mathrm{mean}\!\left(\frac{H-L}{C}, 20\right), \\
  \text{range\_60} &= -\frac{\max(H,60) - \min(L,60)}{C}.
\end{align}

\texttt{ivol\_60}（特质波动率）比总波动更干净，因为它剥离了市场 beta 的影响。
\texttt{max\_ret\_20}（最大单日收益）是彩票偏好最直接的代理变量，实测效果常常优于
单纯的波动率。

\section{流动性族}

\subsection{经济逻辑}

高换手意味着高分歧、高投机、高拥挤。流动性差的股票需要更高预期收益来补偿变现成本。

\subsection{核心因子}

\begin{align}
  \text{turn\_20} &= -\mathrm{mean}(\text{Turnover}, 20), \\
  \text{turn\_std\_20} &= -\mathrm{std}(\text{Turnover}, 20), \\
  \text{turn\_chg} &= -\left(\frac{\mathrm{mean}(\text{Turn},5)}{\mathrm{mean}(\text{Turn},60)} - 1\right), \\
  \text{amihud\_20} &= \mathrm{mean}\!\left(\frac{|r|}{\text{Amount}}, 20\right)\times 10^8, \\
  \text{ln\_amount\_20} &= -\ln\big(\mathrm{mean}(\text{Amount}, 20)\big).
\end{align}

其中 Amihud 非流动性指标的含义是"单位成交额引起的价格冲击"：
数值越大，说明这只股票越难在不推动价格的情况下买卖。

\begin{warnbox}[换手率的口径问题]
换手率 $=$ 成交量 $/$ 流通股本。流通股本在送转、增发时会变，
而免费数据源通常只给"当前股本"。用当前股本算历史换手率，
会在送转当年产生系统性偏差。本书的做法是用\textbf{月度股本快照前向填充}，
把误差限制在一个月内。
\end{warnbox}

\section{量价族}

\begin{align}
  \text{pv\_corr\_20} &= -\mathrm{corr}(r, \Delta V/V, 20), \\
  \text{gap\_20} &= -\mathrm{mean}\!\left(\frac{O_t}{C_{t-1}} - 1, 20\right), \\
  \text{hl\_20} &= -\mathrm{mean}\!\left(\frac{H}{L} - 1, 20\right).
\end{align}

\texttt{pv\_corr\_20} 衡量"量价是否同步"：放量上涨的股票持续性更差，
因此取负号后因子值越大越好。

\section{价值族}

\subsection{经济逻辑}

估值低的股票长期跑赢，可能因为承担了困境风险，也可能因为市场对坏消息过度反应。

\subsection{核心因子}

\begin{align}
  \text{bp} &= \frac{\text{BPS}}{C}, \qquad
  \text{ep\_ttm} = \frac{\text{EPS}_{\text{TTM}}}{C}, \\
  \text{sp\_ttm} &= \frac{\text{Revenue}_{\text{TTM}}}{\text{MktCap}}, \qquad
  \text{cfp} = \frac{\text{CFO\_PS}_{\text{TTM}}}{C}.
\end{align}

\subsection{TTM 的构造}

财报是累计口径：一季报的净利润是 1--3 月的累计，年报是 1--12 月的累计。
直接比较两者毫无意义。TTM（滚动十二个月）的还原方式是：

\begin{equation}
  \text{TTM}_t =
  \begin{cases}
    \text{累计}_t, & \text{若为年报}, \\[4pt]
    \text{累计}_t + \text{上年年报} - \text{上年同期累计}, & \text{否则}.
  \end{cases}
  \label{eq:ttm}
\end{equation}

\section{质量族}

\begin{align}
  \text{roe} &= \text{加权 ROE}, \\
  \text{gross\_margin} &= \text{销售毛利率}, \\
  \text{ocf\_to\_eps} &= \frac{\text{CFO\_PS}_{\text{TTM}}}{|\text{EPS}_{\text{TTM}}|}, \\
  \text{accrual} &= -\frac{\text{EPS}_{\text{TTM}} - \text{CFO\_PS}_{\text{TTM}}}{|\text{EPS}_{\text{TTM}}|}.
\end{align}

\texttt{accrual}（应计项）衡量"会计利润里有多少还没变成现金"。
应计越高，利润质量越差，后续越容易向下修正。这是 Sloan（1996）以来
最稳健的会计异象之一。

\section{成长族}

\begin{align}
  \text{rev\_yoy} &= \text{营业总收入同比}, \\
  \text{np\_yoy} &= \text{归母净利润同比}, \\
  \text{rev\_accel} &= \text{rev\_yoy}_t - \text{rev\_yoy}_{t-60}.
\end{align}

\texttt{rev\_accel}（成长加速度）是一个二阶量：市场对"增速在加快"的定价
往往比对"增速高"更敏感。

\section{规模与风险}

\begin{align}
  \text{size} &= -\ln(\text{流通市值}), \\
  \text{beta\_60} &= -\beta_{120}.
\end{align}

\begin{warnbox}[规模因子：不是衰减，是只在不可执行的区段里存在]
小市值在 2012--2016 年是 A 股最强的因子。2017 年之后，随着监管趋严、
壳价值消失、注册制扩容，它的效力大幅削弱，甚至在某些区间反向。
如果你在 2015 年之前的数据上挖因子，会得到一个"小市值无敌"的结论，
而这个结论在样本外会失效。
\end{warnbox}

\section{因子相关性：你在重复计算同一个故事}

40 个因子里，真正独立的信息可能只有 8--12 个维度。检验方式有两种：

\begin{enumerate}
  \item \textbf{截面相关性}：每天算一次因子之间的 Spearman 相关，再对时间取平均。
        相关系数 $> 0.7$ 的两个因子，基本可以视为同一个。
  \item \textbf{IC 序列相关性}：比较两个因子的 IC 时间序列是否同步。
        这衡量的是"它们在什么时候有效"是否一致。
\end{enumerate}

\begin{pycode}[因子相关性矩阵]
C = factor_corr({name: factor_matrix[name] for name in names})
# 距离 = 1 - |corr|，用层次聚类把 40 个因子压成 10 个簇
# 每簇只保留经济含义最清晰的一个代表
\end{pycode}

\begin{keybox}[因子库的正确用法]
不是"把 40 个因子一起喂给模型"，而是：
先聚类 → 每簇选代表 → 用代表因子做组合 → 用簇内其他因子做稳健性检验。
模型输入 10 个正交性较好的因子，比输入 40 个高度相关的因子更稳。
\end{keybox}

\section{本书因子库全表}

\begin{center}
\small
\begin{longtable}{@{}p{0.14\textwidth}p{0.09\textwidth}p{0.42\textwidth}p{0.06\textwidth}p{0.08\textwidth}@{}}
\toprule
\textbf{因子} & \textbf{家族} & \textbf{公式} & \textbf{方向} & \textbf{回看} \\
\midrule
\endhead
\multicolumn{5}{@{}l}{\textit{反转 / 动量}}\\
\texttt{rev\_1} & 反转 & $-(C_t/C_{t-1}-1)$ & $+$ & 2 \\
\texttt{rev\_5} & 反转 & $-(C_t/C_{t-5}-1)$ & $+$ & 7 \\
\texttt{rev\_10} & 反转 & $-(C_t/C_{t-10}-1)$ & $+$ & 12 \\
\texttt{mom\_20} & 动量 & $C_t/C_{t-20}-1$ & $+$ & 22 \\
\texttt{mom\_60} & 动量 & $C_t/C_{t-60}-1$ & $+$ & 62 \\
\texttt{mom\_120} & 动量 & $C_t/C_{t-120}-1$ & $+$ & 122 \\
\texttt{mom\_12\_1} & 动量 & $C_{t-20}/C_{t-240}-1$ & $+$ & 242 \\
\texttt{resid\_mom\_60} & 动量 & $\text{mom}_{60}-\beta\,\text{mom}^{\text{mkt}}_{60}$ & $+$ & 130 \\
\midrule
\multicolumn{5}{@{}l}{\textit{波动 / 风险}}\\
\texttt{vol\_20} & 波动 & $-\mathrm{std}(r,20)$ & $+$ & 22 \\
\texttt{vol\_60} & 波动 & $-\mathrm{std}(r,60)$ & $+$ & 62 \\
\texttt{ivol\_60} & 波动 & $-\mathrm{std}(r-\beta r^{m},60)$ & $+$ & 125 \\
\texttt{skew\_20} & 波动 & $-\mathrm{skew}(r,20)$ & $+$ & 22 \\
\texttt{kurt\_20} & 波动 & $-\mathrm{kurt}(r,20)$ & $+$ & 22 \\
\texttt{max\_ret\_20} & 波动 & $-\max(r,20)$ & $+$ & 22 \\
\texttt{amp\_20} & 波动 & $-\mathrm{mean}((H-L)/C,20)$ & $+$ & 22 \\
\texttt{range\_60} & 波动 & $-(\max H_{60}-\min L_{60})/C$ & $+$ & 62 \\
\texttt{beta\_60} & 风险 & $-\beta_{120}$ & $+$ & 122 \\
\midrule
\multicolumn{5}{@{}l}{\textit{流动性}}\\
\texttt{turn\_20} & 流动性 & $-\mathrm{mean}(\text{Turn},20)$ & $+$ & 22 \\
\texttt{turn\_std\_20} & 流动性 & $-\mathrm{std}(\text{Turn},20)$ & $+$ & 22 \\
\texttt{turn\_chg} & 流动性 & $-\big(\bar{T}_5/\bar{T}_{60}-1\big)$ & $+$ & 62 \\
\texttt{amihud\_20} & 流动性 & $\mathrm{mean}(|r|/\text{Amt},20)\!\times\!10^8$ & $+$ & 22 \\
\texttt{ln\_amount\_20} & 流动性 & $-\ln(\mathrm{mean}(\text{Amt},20))$ & $+$ & 22 \\
\texttt{vol\_ratio} & 流动性 & $-\big(\bar{V}_5/\bar{V}_{60}-1\big)$ & $+$ & 62 \\
\midrule
\multicolumn{5}{@{}l}{\textit{量价}}\\
\texttt{pv\_corr\_20} & 量价 & $-\mathrm{corr}(r,\Delta V/V,20)$ & $+$ & 22 \\
\texttt{gap\_20} & 量价 & $-\mathrm{mean}(O_t/C_{t-1}-1,20)$ & $+$ & 22 \\
\texttt{hl\_20} & 量价 & $-\mathrm{mean}(H/L-1,20)$ & $+$ & 22 \\
\midrule
\multicolumn{5}{@{}l}{\textit{价值}}\\
\texttt{bp} & 价值 & $\text{BPS}/C$ & $+$ & 1 \\
\texttt{ep\_ttm} & 价值 & $\text{EPS}_{\text{TTM}}/C$ & $+$ & 1 \\
\texttt{sp\_ttm} & 价值 & $\text{Rev}_{\text{TTM}}/\text{MktCap}$ & $+$ & 1 \\
\texttt{cfp} & 价值 & $\text{CFO\_PS}_{\text{TTM}}/C$ & $+$ & 1 \\
\texttt{div\_yield} & 价值 & 股息率 & $+$ & 1 \\
\midrule
\multicolumn{5}{@{}l}{\textit{质量 / 成长}}\\
\texttt{roe} & 质量 & 加权 ROE & $+$ & 1 \\
\texttt{gross\_margin} & 质量 & 销售毛利率 & $+$ & 1 \\
\texttt{ocf\_to\_eps} & 质量 & $\text{CFO\_PS}/|\text{EPS}|$ & $+$ & 1 \\
\texttt{accrual} & 质量 & $-(\text{EPS}-\text{CFO\_PS})/|\text{EPS}|$ & $+$ & 1 \\
\texttt{rev\_yoy} & 成长 & 营收同比 & $+$ & 1 \\
\texttt{np\_yoy} & 成长 & 净利同比 & $+$ & 1 \\
\texttt{rev\_accel} & 成长 & $\Delta$ 营收同比 & $+$ & 62 \\
\midrule
\multicolumn{5}{@{}l}{\textit{规模}}\\
\texttt{size} & 规模 & $-\ln(\text{流通市值})$ & $+$ & 1 \\
\multicolumn{5}{@{}l}{\quad\itshape 附录~\ref{app:strict}：剔除最小 30\% 市值后 IC 翻负，判定为不可用}\\
\bottomrule
\end{longtable}
\end{center}

\section{本章小结}

\begin{enumerate}
  \item 好因子要同时满足经济逻辑、稳健性、可计算性、容量四条。
  \item 因子库必须登记公式、方向与回看期；方向是"我的预期"，与实测不符时要停下来查。
  \item 反转、低波、流动性是 A 股最稳的三个家族；中长期动量在 A 股日频上往往为负。
  \item 规模因子已显著衰减，用它之前必须分区间验证。
  \item 40 个因子的有效维度可能只有 10 个，聚类与正交化是必做功课。
\end{enumerate}

\chapter{因子预处理：为什么你的"新因子"只是小盘股的换皮}\label{ch:prep}

\begin{keybox}[本章要回答的问题]
去极值该用 3$\sigma$ 还是 MAD？标准化用 z-score 还是排名？中性化到底在剥离什么？
为什么"中性化后 IC 一定更高"是错的？
\end{keybox}

\section{为什么必须预处理}

原始因子值直接使用会遇到三个问题：

\begin{enumerate}
  \item \textbf{极端值主导}。A 股有大量"净资产为负 → BP 爆炸""净利润接近零 → EP 爆炸"
        的样本。一个 BP 等于 500 的股票会把线性回归的系数完全带偏。
  \item \textbf{量纲不可比}。换手率是百分数（0--30），ROE 是百分数（$-$50--50），
        市值是元（$10^8$--$10^{12}$）。不加处理无法相加。
  \item \textbf{风格暴露}。你的"新因子"可能只是"小盘股"或"医药行业"的换皮。
\end{enumerate}

\section{去极值}

\subsection{三种方法}

\begin{description}
  \item[3$\sigma$ 法] 把超过 $\mu \pm 3\sigma$ 的值截断。问题是 $\mu,\sigma$ 本身
        被极端值污染——厚尾分布下 $\sigma$ 会被拉大，导致"该截的没截掉"。
  \item[百分位法] 把超过 $1\%$ 与 $99\%$ 分位的值截断。简单稳健，但在样本很少时
        分位数估计不稳。
  \item[MAD 法] 用中位数绝对偏差：
        \begin{equation}
          \text{MAD} = \mathrm{median}\big(|x_i - \mathrm{median}(x)|\big), \qquad
          \hat\sigma = 1.4826\,\text{MAD}.
        \end{equation}
        常数 $1.4826$ 使 $\hat\sigma$ 在正态分布下与标准差一致。
        截断区间取 $[\mathrm{median} - n\hat\sigma,\ \mathrm{median} + n\hat\sigma]$，
        $n$ 常取 5。
\end{description}

\begin{keybox}[本书的选择]
截面 MAD 法，$n=5$。理由：中位数与 MAD 的崩溃点是 50\%，
即使有 20\% 的极端样本也不影响估计；而 $3\sigma$ 法的崩溃点是 0\%。
\end{keybox}

\section{标准化}

\subsection{z-score 与排名}

\begin{align}
  z_{i} &= \frac{x_i - \mathrm{mean}(x)}{\mathrm{std}(x)}, \\
  r_{i} &= \mathrm{rank}(x_i)/N - \tfrac12 \;\in\; \left(-\tfrac12, \tfrac12\right].
\end{align}

\begin{center}
\begin{tabular}{@{}lll@{}}
\toprule
 & \textbf{z-score} & \textbf{排名} \\
\midrule
抗极值 & 弱（受尾部影响） & 强 \\
保留幅度 & 是 & 否 \\
分布形状 & 近似正态 & 均匀 \\
适合 & 线性模型、组合权重 & 稳健评价、IC 计算 \\
\bottomrule
\end{tabular}
\end{center}

实践中常见的折中：\textbf{先去极值，再 z-score}。这样既保留了幅度，又不被极端值主导。

\section{中性化：剥离你不想要的风险}

\subsection{在剥离什么}

假设你构造了一个"研发投入强度"因子，回测发现 IC 很高。但仔细一看：
高研发投入的公司集中在创业板，而创业板在样本期内整体跑赢。
那么你的因子赚的其实是"创业板 beta"，而不是"研发"本身。

中性化就是把这个"其实不是你的"部分剥离掉。

\subsection{行业中性化}

最简单的做法是行业内 demean：

\begin{equation}
  \tilde x_{i} = x_i - \mathrm{mean}_{j \in \text{行业}(i)}(x_j).
\end{equation}

等价于对行业哑变量做无截距回归取残差。

\subsection{行业 + 市值中性化}

更完整的做法是截面回归：

\begin{equation}
  x_{i,t} = \alpha_t + \sum_{k=1}^{K-1}\beta_{k,t}\,\mathbb{1}[i\in \text{行业}_k]
  + \gamma_t \ln(\text{MktCap}_{i,t}) + \varepsilon_{i,t},
  \label{eq:neutral}
\end{equation}

取残差 $\varepsilon_{i,t}$ 作为中性化后的因子值。

\begin{proposition}[残差的正交性]
由最小二乘的正规方程，残差 $\hat\varepsilon$ 与设计矩阵的每一列都正交：
\begin{equation}
  \sum_i \hat\varepsilon_i = 0, \qquad
  \sum_i \hat\varepsilon_i \mathbb{1}[i\in k] = 0\ \forall k, \qquad
  \sum_i \hat\varepsilon_i \ln(\text{MktCap}_i) = 0.
\end{equation}
因此中性化后的因子在截面上\textbf{行业均值为零、与对数市值零相关}。
这正是"剥离"的数学含义。
\end{proposition}

\subsection{为什么不用分组 demean}

分组 demean 只能处理离散分组（行业），无法同时控制连续的市值。
而且它无法处理"行业内市值差异"带来的残余暴露。
回归可以一次性处理两者，且残差严格正交，是更干净的做法。

\subsection{几个工程细节}

\begin{enumerate}
  \item \textbf{市值用对数}。市值分布严重右偏，取对数后接近正态；
        直接用原始市值会让极少数巨型股主导回归。
  \item \textbf{哑变量丢一列}。$K$ 个行业要用 $K-1$ 个哑变量加截距，
        否则完全共线。
  \item \textbf{截面做，不时序做}。因子有效性是横截面问题，
        每天独立回归一次。
  \item \textbf{样本太少时放弃}。当天有效样本 $<30$ 时，回归系数不稳，宁可不中性化。
  \item \textbf{二次项通常不需要}。线性项已经足够，加二次项会增加自由度与过拟合风险；
        除非你明确怀疑存在 U 形关系且样本量足够。
\end{enumerate}

\section{一个反直觉的结论：中性化不总是提升 IC}

很多教程写"中性化后 IC 会提高"。这是错的。

\begin{itemize}
  \item 如果因子的信息部分来自行业轮动（例如"景气度向上行业的股票"），
        行业中性化会把真实 alpha 一起剥掉，IC 反而下降。
  \item 如果因子本身就是一个风格因子（例如规模因子），
        对市值中性化等于把它归零。
\end{itemize}

\begin{warnbox}[正确的判断方式]
中性化前后各算一遍 IC，看\textbf{降低幅度}。
\begin{itemize}
  \item IC 几乎不变：因子本身与行业/市值无关，中性化是"买个保险"。
  \item IC 明显下降：因子有相当一部分收益来自风格暴露，中性化后剩下的才是真 alpha。
  \item IC 反而上升：说明原始因子被风格噪声污染，中性化起了降噪作用。
\end{itemize}
三种情况都是有用信息，但只有第三种能说"中性化提升了因子"。
\end{warnbox}

\section{实现：向量化的截面回归}

逐日循环回归在小样本上够用，但在 $250$ 天 $\times$ 40 个因子上会明显变慢。
一个有效的优化是：\textbf{同一天的所有因子共用同一个设计矩阵}，
因此对 $X$ 做一次 QR 分解，就能把整块因子矩阵投影掉。

\begin{pycode}[批量残差化]
def neutralize_block(Y, X):
    """Y: (n, m) 因子矩阵；X: (n, k) 设计矩阵。返回残差 (n, m)。"""
    Q, _ = np.linalg.qr(X)        # 只分解一次
    return Y - Q @ (Q.T @ Y)      # 一次性投影掉所有因子
\end{pycode}

复杂度从 $O(m\,n k^2)$ 降到 $O(nk^2 + mnk)$，在 $m=40$、$n=4000$、$k=32$ 时
大约快 10 倍。

\section{预处理流程的完整顺序}

顺序不能乱，理由如下：

\begin{enumerate}
  \item \textbf{剔除不可交易样本}：ST、停牌、新股、流动性不足。
        先剔除，否则它们的极端值会污染去极值的统计量。
  \item \textbf{去极值}：MAD，$n=5$。
  \item \textbf{标准化}：截面 z-score。
  \item \textbf{中性化}：行业 + 对数市值回归取残差。
  \item \textbf{再次标准化}：中性化会改变量纲，需重新标准化以便加权。
\end{enumerate}

\begin{warnbox}[一个容易被忽略的顺序问题]
如果先中性化再去极值，去极值会破坏残差的正交性——
你截断掉的极端值，恰好可能是与市值高度相关的那些点。
正确顺序是\textbf{先去极值、再中性化}。
\end{warnbox}

\section{本章小结}

\begin{enumerate}
  \item 去极值用 MAD 而不是 $3\sigma$，因为后者自身被极端值污染。
  \item 标准化保留幅度去 z-score，追求稳健去排名；实践中"先去极值再 z-score"最常用。
  \item 中性化 = 截面回归取残差，残差与行业、市值严格正交。
  \item 中性化不总提升 IC；IC 下降说明你的因子有相当部分是风格暴露。
  \item 处理顺序：剔除不可交易 → 去极值 → 标准化 → 中性化 → 再标准化。
\end{enumerate}

\chapter{单因子评价：怎样判断一个因子不是噪声}\label{ch:eval}

\begin{keybox}[本章要回答的问题]
IC 到底在衡量什么？ICIR 的年化口径为什么会差 4 倍？Newey-West 的滞后阶数怎么选？
分层单调性为什么比 IC 更有说服力？挖了 500 个因子之后，怎么诚实地报告结果？
\end{keybox}

\section{IC：一个被过度简化却仍然有用的指标}

信息系数（Information Coefficient）的定义是：因子值与未来收益在\textbf{同一时点横截面}
上的相关系数。

\begin{equation}
  \mathrm{IC}_t = \mathrm{corr}\big(f_{i,t},\ r_{i,t+1\rightarrow t+1+h}\big),
  \quad i \in \text{当日可交易股票}.
\end{equation}

\subsection{两种口径}

\begin{description}
  \item[Pearson IC] 直接算线性相关。受极端值影响大：一只 BP 爆表的股票就能把 IC 拉高或拉低。
  \item[Rank IC] 先做截面排名再算相关，等价于 Spearman 相关。抗极值，是选股因子的标准口径。
\end{description}

\begin{warnbox}[不要在同一个研究里混用两种口径]
Pearson IC 与 Rank IC 的数值可以差一倍以上。如果你在研报里看到"IC 0.08"，
先问清楚是哪种口径、什么持有期、什么股票池。
本书除非特别说明，一律指 \textbf{Rank IC，$h=5$ 日持有期，全 A 股可交易池}。
\end{warnbox}

\section{ICIR 与年化口径的陷阱}

\begin{equation}
  \mathrm{ICIR} = \frac{\overline{\mathrm{IC}}}{\mathrm{std}(\mathrm{IC})},
  \qquad
  \mathrm{ICIR}_{\mathrm{ann}} = \mathrm{ICIR}\times\sqrt{K},
\end{equation}

其中 $K$ 是"一年里有多少个独立的预测周期"。

\begin{center}
\begin{tabular}{@{}lccc@{}}
\toprule
采样频率 & 每年期数 $K$ & $\sqrt{K}$ & 说明 \\
\midrule
日频（$h=1$） & 252 & 15.87 & 无重叠，可直接年化 \\
日频（$h=5$） & 252/5 = 50.4 & 7.10 & \textbf{必须按有效周期折算} \\
周频 & 52 & 7.21 & \\
月频 & 12 & 3.46 & \\
\bottomrule
\end{tabular}
\end{center}

\begin{warnbox}[最高频的错误：日频算 IC，乘 $\sqrt{252}$]
如果你每天算一次 IC，但预测的是未来 5 日收益，那么相邻两天的 IC 高度自相关
（相关系数常大于 0.8）。此时乘 $\sqrt{252}$ 会把 ICIR 虚增 2 到 4 倍。

正确做法二选一：
\begin{itemize}
  \item 把年化乘数改为 $\sqrt{252/h}$；
  \item 保留 $\sqrt{252}$，但把分母换成 Newey-West 修正后的标准误。
\end{itemize}
\end{warnbox}

\section{显著性：Newey-West $t$ 统计量}

即便用对了年化口径，直接用 $\overline{\mathrm{IC}}/(\hat\sigma/\sqrt{n})$ 算 $t$ 值仍然错，
因为 IC 序列存在自相关。

\begin{equation}
  \hat\sigma^2_{\mathrm{NW}} = \hat\gamma_0 + 2\sum_{l=1}^{L}\left(1 - \frac{l}{L+1}\right)\hat\gamma_l,
  \qquad
  \hat\gamma_l = \frac{1}{n}\sum_{t=l+1}^{n}(\mathrm{IC}_t - \overline{\mathrm{IC}})(\mathrm{IC}_{t-l} - \overline{\mathrm{IC}}),
\end{equation}

\begin{equation}
  t_{\mathrm{NW}} = \frac{\overline{\mathrm{IC}}}{\sqrt{\hat\sigma^2_{\mathrm{NW}}/n}}.
\end{equation}

\subsection{滞后阶数怎么选}

三个准则，取最大值最稳妥：

\begin{enumerate}
  \item \textbf{持有期准则}：$L \ge h - 1$。重叠标签会机械地引入 $h-1$ 阶自相关。
  \item \textbf{经验法则}：$L = \lfloor 4(n/100)^{2/9}\rfloor$。
        2500 个观测时 $L\approx 8$。
  \item \textbf{数据驱动}：用 IC 序列的自相关函数衰减到 0 的阶数。
\end{enumerate}

本书的实现取 $L=\max(h-1,\ \lfloor 4(n/100)^{2/9}\rfloor)$。

\section{分层回测：比 IC 更有说服力的证据}

把股票按因子值分成 $Q$ 组，看每组的平均收益是否单调。

\begin{center}
\begin{tabular}{@{}lccccc@{}}
\toprule
 & Q1（最低） & Q2 & Q3 & Q4 & Q5（最高） \\
\midrule
月均收益 & $-0.4\%$ & $-0.1\%$ & $0.2\%$ & $0.5\%$ & $0.9\%$ \\
\bottomrule
\end{tabular}
\end{center}

\begin{keybox}[为什么单调性比 IC 重要]
IC 是线性相关的平均值，它可以被"两端有效、中间乱序"的形态撑起来。
而单调性要求\textbf{每一组都按顺序排列}，这种形态更难由噪声产生。
一个 IC 只有 $0.03$ 但五组完美单调的因子，比一个 IC 有 $0.06$ 但中间三组乱序的因子更可信。
\end{keybox}

工程上把单调性量化为一件事：五组平均收益与组序号 $(1,2,3,4,5)$ 的相关系数。
绝对值大于 $0.9$ 才算合格。

\section{换手率与因子衰减}

\subsection{换手率决定成本}

因子的换手率定义为：相邻两期最高分位组合的成分变化比例。

\begin{equation}
  \mathrm{TO}_t = 1 - \frac{|S_t \cap S_{t-1}|}{(|S_t| + |S_{t-1}|)/2}.
\end{equation}

\begin{center}
\begin{tabular}{@{}lccc@{}}
\toprule
因子 & 单边换手（5 日） & 年化双边成本（30bp 单边） & 判断 \\
\midrule
\texttt{rev\_5} & 约 35\% & 约 10.5\% & 高，需要强信号才划算 \\
\texttt{mom\_20} & 约 16\% & 约 4.8\% & 中 \\
\texttt{turn\_20} & 约 2\% & 约 0.6\% & 低，适合做底仓 \\
\texttt{size} & 约 1.5\% & 约 0.45\% & 极低 \\
\bottomrule
\end{tabular}
\end{center}

\subsection{衰减：这个信号多久失效}

分别用 $h=1,5,10,20,60$ 算 IC，画出衰减曲线。

\begin{itemize}
  \item \textbf{快速衰减}（5 日内衰减到一半）：反转、量价。必须短持有期，成本高。
  \item \textbf{慢速衰减}（20 日仍有 70\%）：价值、质量。可以长持有期，成本低。
\end{itemize}

\begin{warnbox}[高 IC 加高换手等于陷阱]
一个 IC 为 $0.08$、单边换手 $50\%$ 的因子，在 $30$ bp 单边成本下，
每 5 日成本约 $0.3\%$，年化约 $15\%$——足以吃掉全部 alpha。
\textbf{永远先算成本，再看 IC。}
\end{warnbox}

\section{多重检验：500 个因子留下 20 个，怎么诚实报告}

假设你挖了 500 个候选因子，其中 20 个的 $t$ 值大于 2（$p<0.05$）。
这 20 个里有多少是真的？

如果 500 个全部是噪声，按 $5\%$ 的显著性水平，你期望看到约 25 个"显著"。
也就是说：\textbf{你看到的 20 个，比纯噪声下的期望还少}。

\subsection{Benjamini-Hochberg（BH）流程}

控制错误发现率（FDR）而不是族错误率（FWER），在因子挖掘里更实用：

\begin{enumerate}
  \item 对 500 个因子分别算 $t_{\mathrm{NW}}$，转成双侧 $p$ 值；
  \item 把 $p$ 值升序排列：$p_{(1)} \le p_{(2)} \le \cdots \le p_{(500)}$；
  \item 找最大的 $k^*$ 使得 $p_{(k)} \le \dfrac{k}{500}\,q$，其中 $q$ 取 $0.05$；
  \item 保留前 $k^*$ 个因子。注意：必须\textbf{从大到小}找 $k^*$，
        而不是从小到大遇到第一个不满足就停。
\end{enumerate}

\begin{pycode}[Benjamini-Hochberg 校正]
def benjamini_hochberg(pvals, q=0.05):
    p = np.asarray(pvals)
    order = np.argsort(p)
    m = len(p)
    thresh = q * (np.arange(1, m + 1) / m)
    passed = p[order] <= thresh
    k = np.max(np.where(passed)[0]) + 1 if passed.any() else 0
    keep = np.zeros(m, dtype=bool)
    keep[order[:k]] = True
    return keep
\end{pycode}

\subsection{相关性会破坏独立性假设}

BH 假设检验之间独立或满足正相关（PRDS）。但因子库里 40 个因子可能有 30 个
高度相关——一旦某个簇里的"底层噪声"偶然显著，整簇会同时越过阈值，
造成"全有或全无"的假阳性爆发。

\textbf{工程上的解法}：先聚类去冗余，再对簇代表做 BH。

\begin{enumerate}
  \item 算因子两两的截面相关矩阵；
  \item 用距离 $d_{ij} = \sqrt{\tfrac12(1-\rho_{ij})}$ 做层次聚类；
  \item 每簇选一个代表（经济含义最清晰、或 IC 最高）；
  \item 只对代表因子做 BH，分母从 500 降到簇数。
\end{enumerate}

\section{一个完整的因子评价表}

本书对每个因子输出下面这张表（完整版见第~\ref{ch:result} 章）：

\begin{center}
\small
\begin{tabular}{@{}lcccccc@{}}
\toprule
因子 & IC 均值 & ICIR(年化) & NW $t$ & 分层单调性 & 单边换手 & 中性化后 IC \\
\midrule
（见报告） & & & & & & \\
\bottomrule
\end{tabular}
\end{center}

\begin{warnbox}[读表的顺序]
先看\textbf{中性化后 IC}，再看 \textbf{NW $t$}，最后看\textbf{换手}。
如果中性化后 IC 掉到接近 0，说明这个因子是风格暴露的换皮，后面的数字都不必看。
\end{warnbox}

\section{本章小结}

\begin{enumerate}
  \item Rank IC 是选股因子的标准口径，必须与持有期、股票池一起声明。
  \item 日频 $h=5$ 的 ICIR 年化乘数是 $\sqrt{50}$ 而不是 $\sqrt{252}$。
  \item IC 序列有自相关，$t$ 值必须用 Newey-West 修正，滞后阶数 $L\ge h-1$。
  \item 分层单调性比 IC 均值更能说明"不是噪声"。
  \item 先算成本再看 IC：高 IC 高换手的因子经常是负贡献。
  \item 挖了 500 个因子，必须做 FDR 校正，并且先聚类去冗余。
\end{enumerate}

\chapter{从因子到预测：线性、树模型与深度学习}\label{ch:model}

\begin{keybox}[本章要回答的问题]
多因子怎么合成？为什么必须先打败等权基线？LightGBM 在横截面上应该怎么调？
深度学习什么时候值得上？模型"学错了"怎么发现？
\end{keybox}

\section{必须打败的基线：等权合成}

在引入任何模型之前，先建立一个基线：

\begin{equation}
  \text{score}_{i,t} = \sum_{k=1}^{K} s_k\, z_{k,i,t},
  \qquad
  s_k = \mathrm{sign}\!\left(\IC_k^{\text{train}}\right),
\end{equation}

其中 $z_k$ 是第 $k$ 个因子当日截面 z-score，$s_k$ 是训练集上该因子 IC 的符号。

这个"用 IC 符号定向 + 等权相加"的方法简单到近乎粗暴，但它有两个重要性质：

\begin{enumerate}
  \item \textbf{零参数}。除了符号之外没有任何可调参数，因此几乎没有过拟合空间。
  \item \textbf{抗噪声}。单个因子的估计误差被平均掉，不会像回归那样被共线性放大。
\end{enumerate}

\begin{warnbox}[一个残酷的经验]
在日频 A 股上，很多"精心调参的 LightGBM"跑不赢这个基线。
如果你的模型没有稳定跑赢它，那么模型带来的不是 alpha，而是复杂度。
本书的每一次模型对比，都以它作为对照。
\end{warnbox}

\section{线性回归：可解释但脆弱}

\begin{equation}
  y_{i,t} = \sum_{k}\beta_k z_{k,i,t} + \varepsilon_{i,t}.
\end{equation}

\subsection{为什么必须加正则}

40 个因子之间存在高度共线性（反转与换手、波动与振幅）。普通最小二乘在这种情况下
系数会剧烈震荡，甚至符号翻转。岭回归通过惩罚系数平方和来稳定：

\begin{equation}
  \hat\beta = \arg\min_{\beta}\ \|y - Z\beta\|^2 + \alpha\|\beta\|^2
  = (Z^\top Z + \alpha I)^{-1}Z^\top y.
\end{equation}

$\alpha$ 越大，系数越向零收缩。实践中 $\alpha$ 取 $1$--$100$（因子已标准化），
用样本外 IC 选择，不要用样本内 $R^2$——后者永远偏好 $\alpha\to 0$。

\subsection{线性模型的价值}

它的价值不在于预测精度，而在于\textbf{可解释性}：

\begin{itemize}
  \item 系数符号是否符合经济逻辑？
  \item 哪个因子贡献最大？是否与你的先验一致？
  \item 系数在不同年份是否稳定？
\end{itemize}

这些问题，线性模型能直接回答，树模型不能。\textbf{先用线性模型把因子看清楚，
再上树模型提升精度}，是正确的顺序。

\section{GBDT：表格数据上的默认选择}

梯度提升树的思路是：串行训练一堆弱学习器，每棵树拟合前面所有树的残差。

\begin{equation}
  F_m(x) = F_{m-1}(x) + \nu\, h_m(x),
\end{equation}

其中 $\nu$ 是学习率，$h_m$ 是第 $m$ 棵拟合负梯度的树。

\subsection{为什么它在横截面选股上仍然最强}

\begin{enumerate}
  \item \textbf{自动处理非线性与交互}。它天然能学到"低波因子只在小市值股票里有效"。
  \item \textbf{对量纲与单调变换不敏感}。不需要精细的特征工程。
  \item \textbf{对缺失值鲁棒}。A 股因子缺失很常见（财务数据滞后、停牌）。
  \item \textbf{训练快}。300 万样本 $\times$ 40 特征在普通笔记本上几分钟即可。
\end{enumerate}

\subsection{日频横截面上的正确用法}

\begin{center}
\begin{tabular}{@{}lll@{}}
\toprule
参数 & 推荐范围 & 理由 \\
\midrule
\texttt{objective} & \texttt{regression} & 回归排序标签；也可用自定义 IC 损失 \\
\texttt{num\_leaves} & 15--63 & 太大会记住噪声 \\
\texttt{max\_depth} & 3--6 & 与 \texttt{num\_leaves} 配合 \\
\texttt{min\_child\_samples} & 200--500 & 横截面样本多，叶子要足够大 \\
\texttt{learning\_rate} & 0.02--0.05 & 配合 \texttt{n\_estimators} 300--800 \\
\texttt{feature\_fraction} & 0.7--0.9 & 抑制共线因子的相互替代 \\
\texttt{bagging\_fraction} & 0.7--0.9 & \\
\texttt{reg\_lambda} & 1--10 & \\
\bottomrule
\end{tabular}
\end{center}

\begin{warnbox}[早停必须用"未来"的验证集]
LightGBM 的 \texttt{early\_stopping} 需要一个验证集。如果你随机切分，
验证集的未来信息会漏进训练。正确做法是\textbf{按时间顺序}：
训练到 $t-60$，验证用 $[t-60, t-30)$，测试用 $[t-30, t)$。
更严格的做法是直接跳过早停，用固定的树数 + Purged CV 选超参。
\end{warnbox}

\subsection{特征重要性与"模型有没有学错"}

三个诊断动作：

\begin{enumerate}
  \item \textbf{看重要性排序}。如果 \texttt{size} 和 \texttt{turn\_20} 占据前两位，
        而你的因子里有 20 个"更聪明"的信号，说明模型只是在做风格暴露。
  \item \textbf{看重要性稳定性}。在 5 个不同时间窗口上分别训练，重要性的排序
        应该大致稳定。如果每次都不同，说明模型在拟合噪声。
  \item \textbf{做特征置换检验}。把一个因子的值随机打乱，看 IC 下降多少。
        下降很少的因子，模型其实没在用。
\end{enumerate}

\section{深度学习：什么时候值得}

\subsection{它的真实优势}

\begin{itemize}
  \item \textbf{端到端学习表示}。可以从原始量价序列直接学到特征，省去人工构造。
  \item \textbf{处理高维与序列}。LSTM/TCN/Transformer 能建模时序依赖。
  \item \textbf{多任务学习}。同时预测多个持有期，共享表示。
\end{itemize}

\subsection{它的真实代价}

\begin{itemize}
  \item \textbf{数据需求大}。日频全 A 股 10 年约 1000 万个样本，对深度学习来说只是"刚够"。
  \item \textbf{过拟合风险高}。参数量大、信噪比低，是最坏的组合。
  \item \textbf{难以诊断}。出了问题很难定位是数据、标签还是结构。
\end{itemize}

\begin{keybox}[本书的建议]
\textbf{先把线性模型与 GBDT 做到位，再考虑深度学习。}
在日频 A 股上，深度模型的样本外 Rank IC 通常只比 GBDT 高 $0.005$--$0.01$，
而调试成本高出五倍以上。除非你有明确的理由（例如做分钟频、或有多模态数据），
否则不值得。
\end{keybox}

\section{过拟合的十个来源与排查清单}

\begin{center}
\begin{longtable}{@{}p{0.06\textwidth}p{0.34\textwidth}p{0.50\textwidth}@{}}
\toprule
\# & \textbf{来源} & \textbf{排查方法} \\
\midrule
\endhead
1 & 特征里混入未来信息 & 逐列检查数据来源；把因子滞后一期重跑，IC 大幅下降即有问题 \\
2 & 标签与特征时点重叠 & 确认 \texttt{signal\_date < execution\_date} \\
3 & 财务数据用报告期 & 改用公告日 as-of 连接 \\
4 & 幸存者偏差 & 用含退市股的池子重跑 \\
5 & 复权错误 & 检查除权日收益率是否异常 \\
6 & 随机切分 & 改用 Purged K-Fold \\
7 & 用测试集调参 & 记录每次尝试，用 DSR 惩罚 \\
8 & 样本内早停 & 早停验证集必须晚于训练集 \\
9 & 过度精细的因子 & 因子参数在样本外不稳定 \\
10 & 忽略交易成本 & 换手 $\times$ 成本 $>$ alpha 即无效 \\
\bottomrule
\end{longtable}
\end{center}

\section{本章小结}

\begin{enumerate}
  \item 等权合成是必须打败的基线；跑不赢它的模型只是增加了复杂度。
  \item 线性模型的价值在可解释性，必须加正则；$\alpha$ 用样本外 IC 选。
  \item GBDT 在表格数据上仍是默认选择，但叶子数与深度必须压小。
  \item 早停验证集必须严格晚于训练集，否则等于随机切分。
  \item 深度学习的边际收益在日频 A 股上有限，过拟合风险却显著更高。
  \item 过拟合有十个常见来源，逐条排查比换模型有效得多。
\end{enumerate}

\chapter{组合构建：从分数到持仓}\label{ch:portfolio}

\begin{keybox}[本章要回答的问题]
有了每只股票的预测分数，怎么变成持仓？分位数组合与优化器该选哪个？
A 股的 T+1、整手、涨跌停如何进入约束？交易成本怎么估？
\end{keybox}

\section{四种从分数到权重的方法}

\subsection{分位数组合}

把当天可交易股票按分数排序，取最高 $q$ 分位等权做多，最低 $q$ 分位等权做空。

\begin{equation}
  w_i = \begin{cases}
    +\dfrac{1}{|S_{\text{top}}|}, & i \in S_{\text{top}}, \\[6pt]
    -\dfrac{1}{|S_{\text{bot}}|}, & i \in S_{\text{bot}}, \\[6pt]
    0, & \text{否则}.
  \end{cases}
\end{equation}

\textbf{优点}：简单、换手可控、与分层回测的逻辑一致。
\textbf{缺点}：不考虑个股风险差异，等权意味着高波动股贡献更大的组合波动。

\subsection{Top-$N$ 多头}

只做多前 $N$ 只（$N$ 常取 30--100）。

\textbf{优点}：贴近个人实际（多数人没有融券能力）。
\textbf{缺点}：$N$ 越小，特异性风险越大；$N$ 越大，alpha 被稀释。

\subsection{均值--方差优化}

\begin{equation}
  \max_{w}\ w^\top \mu - \frac{\lambda}{2}w^\top \Sigma w
  \quad \text{s.t.}\quad \mathbf{1}^\top w = 1,\ \ |w_i| \le u,
\end{equation}

其中 $\mu$ 是预测收益，$\Sigma$ 是协方差矩阵。

\begin{warnbox}[均值--方差在股票池上的经典失败]
$\Sigma$ 的估计误差会被优化器放大——它会专门挑那些"样本内看起来低风险"的股票，
而这些往往是估计噪声造成的假象。$N=3000$ 时，样本协方差矩阵几乎不可逆，
必须做收缩（Ledoit-Wolf）或因子模型化。
\end{warnbox}

\subsection{带约束的优化（推荐）}

把风险模型、行业中性、换手惩罚一起放进目标函数：

\begin{equation}
  \max_{w}\ w^\top \mu
  - \frac{\lambda}{2}w^\top \Sigma w
  - \kappa\,\|w - w_{\text{prev}}\|_1
  \quad \text{s.t.}\quad
  \begin{cases}
    \mathbf{1}^\top w = 1, \\
    |w_i| \le u, \\
    |\text{行业暴露}_k| \le \delta, \\
    w_i = 0 \ \text{若不可买}.
  \end{cases}
\end{equation}

$\ell_1$ 换手惩罚项是关键：它直接告诉优化器"调仓有成本"，
能显著降低换手而不牺牲太多 alpha。

\section{A 股的四类硬约束}

\subsection{T+1}

1995 年起 A 股实行 T+1：当日买入的股票次日才能卖出。在日频调仓中，这意味着：
\begin{itemize}
  \item 组合的最小持有期是 1 天；
  \item 如果策略每天调仓，实际可交易的只有昨天之前买入的部分。
\end{itemize}

工程上最简单的处理：\textbf{调仓频率不低于 1 天}，并在回测中把当日买入的份额
标记为"不可卖"。

\subsection{整手（100 股）}

资金量小时，整手约束会显著偏离目标权重。例如 10 万元资金买 50 只股票，
每只 2000 元，若某股价 150 元，一手就要 15000 元——根本买不起。

\begin{warnbox}[小资金的核心矛盾]
资金量小于 50 万时，"等权持有 50 只"在数学上就不可行。
此时要么减少持仓数（集中度上升），要么接受权重偏离。
\textbf{任何忽略整手约束的回测，对小资金都是不可执行的。}
\end{warnbox}

\subsection{涨跌停不可成交}

第~\ref{ch:data} 章已经把 \texttt{can\_buy\_next} / \texttt{can\_sell\_next} 做成显式列。
组合构建时必须：
\begin{itemize}
  \item 买入候选集 $=$ 当日可交易 $\cap$ 次日可买；
  \item 若某只股票次日一字涨停，则本期无法建仓，权重分配给其余股票。
\end{itemize}

\subsection{停牌}

停牌股票既不能买也不能卖。若持仓股停牌，其权重会被"冻结"，
组合的实际暴露与目标暴露产生偏离。回测中必须按"最后成交价"估值，
并在复牌后按实际价格结算。

\section{做空约束：为什么 A 股的现实是“指数增强”而不是“多空对冲”}

本书第~\ref{ch:result} 章的所有组合收益都是\textbf{多空}的——做多前 20\%、做空后 20\%。
这一节说明为什么这个假设在 A 股基本不成立，以及正确的替代形态是什么。

\subsection{个股融券基本不可行}

\begin{itemize}
  \item 券源少、费率高，个人投资者很难稳定借到券；
  \item 2023--2024 年监管收紧（提高保证金比例、暂停限售股出借、暂停转融券），
        实操上单票做空已接近不可行；
  \item 即便借到，也面临随时被召回的风险，对冲的稳定性无法保证。
\end{itemize}

\subsection{对冲工具与基差成本}

可行的对冲工具是指数期货（IF 沪深 300、IC 中证 500、IM 中证 1000）与 ETF。
代价是\textbf{基差}：长期贴水意味着对冲成本，中性策略的实际收益要先扣掉这块。
贴水越深，对冲越贵，中性策略的实际收益越低。

\subsection{结构后果：负向信号只能“低配”}

既然不能做空，负向信号唯一的作用就是\textbf{降低权重}，而不是产生负收益。
这就把策略的形态从“多空对冲”改成了\textbf{指数增强}：

\begin{itemize}
  \item 收益 = 基准收益 + （正向超配 + 负向低配）带来的超额；
  \item 风险 = 跟踪误差，而不是市场中性下的绝对波动；
  \item 评价指标 = 信息比率 IR 与跟踪误差，而不是夏普。
\end{itemize}

\begin{keybox}[这对本项目意味着什么]
第~\ref{ch:result} 章报出的“扣费年化 16.3\%、夏普 1.43”是\textbf{多空}口径。
在指数增强形态下，同一批信号只能拿到其中的一部分——
大致相当于“多头端超额减去基准”，而且还要再扣掉跟踪误差的约束成本。

因此本书的所有收益数字都应读作\textbf{信号强度的上界}，
而不是\textbf{可实现的策略收益}。真正的实现版本需要一个指数增强回测框架，
本项目尚未完成（见第~\ref{ch:result} 章的局限说明）。
\end{keybox}

\subsection{程序化交易监管对执行的约束}

2024 年以来，程序化交易在中国进入持牌报告与差异化监管阶段：

\begin{center}
\begin{tabular}{@{}ll@{}}
\toprule
时间 & 事件 \\
\midrule
2024-05-11 & 证监会发布《证券市场程序化交易管理规定（试行）》 \\
2024-10-08 & 上述规定施行 \\
2025-04-03 & 沪深北交易所发布《程序化交易管理实施细则》 \\
2025-07-07 & 实施细则施行 \\
\bottomrule
\end{tabular}
\end{center}

核心要求：\textbf{先报告、后交易}；高频交易认定标准为单个账户
\textbf{每秒申报撤单 300 笔以上}，或\textbf{单日 20000 笔以上}；
对高频交易实行差异化收费与从严管理。

对本项目的实际影响：任何依赖高换手的策略，除了要算清成本，
还要把\textbf{报撤单笔数预算}当作硬约束。这也是第~\ref{ch:result} 章
强调“不要把研究建立在极高换手假设上”的制度原因。

\section{交易成本模型}

\subsection{显性成本}

\begin{center}
\begin{tabular}{@{}lll@{}}
\toprule
项目 & 费率 & 说明 \\
\midrule
佣金 & 0.01\%--0.03\%（双边） & 券商差异大，个人常在 0.025\% \\
印花税 & 0.05\%（仅卖出） & 2023-08-28 起由 0.1\% 下调 \\
过户费 & 0.001\%（双边） & 沪深两市均收 \\
\bottomrule
\end{tabular}
\end{center}

显性成本双边合计约 $0.08\%$--$0.11\%$。

\subsection{隐性成本：滑点与冲击}

这是真正吃掉 alpha 的部分。常用的平方根冲击模型：

\begin{equation}
  \text{impact} \approx k\,\sigma\sqrt{\frac{Q}{V}},
  \label{eq:impact}
\end{equation}

其中 $Q$ 是你的委托量，$V$ 是同期市场成交量，$\sigma$ 是日波动率，
$k$ 是经验常数（A 股约 $0.1$--$0.3$）。

对个人投资者，$Q/V$ 通常很小（小于 $1\%$），冲击成本可以近似为固定滑点：
中小盘 $10$--$20$ bp，大盘 $3$--$8$ bp。

\begin{keybox}[成本参数的合理取值]
本书采用\textbf{单边 30 bp}（含佣金、印花税、滑点、冲击），
对应"中等换手、中等市值"的情形。做敏感性分析时会在 $15$--$60$ bp 之间扫描。
如果你的策略在 $40$ bp 下就失效，说明它的 alpha 太薄，不值得实盘。
\end{keybox}

\section{容量分析}

容量 = 在给定成本下，策略能承载的资金规模。

\begin{equation}
  \text{Capacity} \approx \frac{\alpha \cdot \mathrm{AUM}}{\text{cost}(\mathrm{AUM})}
  \;\Rightarrow\;
  \mathrm{AUM}^* : \ \alpha = k\sigma\sqrt{\frac{w\cdot \mathrm{AUM}}{V}}.
\end{equation}

直观判断：
\begin{itemize}
  \item 持仓越集中、股票市值越小、换手越高，容量越小；
  \item 个人投资者（小于 100 万）几乎不受容量约束；
  \item 但"小市值 + 高频"的策略，即便资金很小，也会因为流动性不足而执行困难。
\end{itemize}

\section{一个最小可用的组合流程}

\begin{pycode}[从分数到持仓]
def build_portfolio(day_df, score, top_n=50, max_weight=0.05):
    d = day_df.assign(s=score)
    d = d[d["can_buy_next"] & d["s"].notna()]
    d = d.nlargest(top_n, "s")
    w = pd.Series(1.0 / len(d), index=d["code"])
    w = w.clip(upper=max_weight)      # 单票上限
    return w / w.sum()                 # 归一
\end{pycode}

这段代码只有 6 行，但它包含了四个必须的约束：可买、分数有效、持仓数上限、单票上限。

\section{本章小结}

\begin{enumerate}
  \item 分位数组合是最稳的起点；均值--方差优化需要收缩协方差，否则会放大估计误差。
  \item $\ell_1$ 换手惩罚是降低成本最有效的手段。
  \item T+1、整手、涨跌停、停牌是四类硬约束，必须显式建模。
  \item 显性成本约 $0.1\%$，隐性成本（滑点加冲击）才是主要杀手。
  \item 单边 30 bp 是合理的基准假设；超过 40 bp 就失效的策略不值得实盘。
\end{enumerate}

\chapter{回测：一个不会报错的谎言生成器}\label{ch:backtest}

\begin{keybox}[本章要回答的问题]
回测最容易在哪里骗你？向量化回测与事件驱动回测有什么区别？
信号日、成交日、成交价应该怎么定？净值曲线该怎么读？
\end{keybox}

\section{回测的七个陷阱}

回测最危险的地方在于：\textbf{它永远不会报错}。数据有前视偏差，它照样跑出一条漂亮的曲线；
成本算少了，它照样给你一个夏普 2.0。下面七条，每条都能让一条"有效策略"变成幻觉。

\subsection{前视偏差}

\begin{itemize}
  \item 用当日收盘价算因子，又用当日收盘价成交；
  \item 用报告期而非公告日对齐财务数据；
  \item 用"最新成分股"回测历史指数增强策略。
\end{itemize}

\textbf{排查}：把所有因子整体滞后一期重跑。如果 IC 大幅下降（比如从 0.06 掉到 0.01），
说明原来的因子大量依赖当日信息。

\subsection{幸存者偏差}

只用今天还活着的股票。退市股、被并购股、长期停牌股被系统性排除。

\textbf{排查}：统计样本期内股票数量的变化。如果回测里股票数只增不减，就是有问题。

\subsection{复权错误}

用不复权价算收益率，除权日会出现假暴跌。

\textbf{排查}：检查单日收益率是否有超过涨跌停幅度的异常值。

\subsection{不可成交}

涨停买进、跌停卖出、停牌交易。

\textbf{排查}：统计"信号选出但实际不可成交"的比例。反转类策略这个比例常常超过 10\%。

\subsection{成本低估}

只算佣金，不算印花税、滑点、冲击。

\textbf{排查}：把成本参数从 10 bp 调到 40 bp，看净值曲线衰减多少。
如果衰减后仍显著，才值得进一步研究。

\subsection{参数过拟合}

在同一段数据上试 200 组参数，选最好的。

\textbf{排查}：第~\ref{ch:overfit} 章的 DSR 与 PBO。

\subsection{多重检验}

试了 50 个因子、10 个模型、5 个持有期，最后报告最好的那个组合。

\textbf{排查}：记录所有尝试，用 FDR 或 DSR 校正。

\section{向量化回测 vs 事件驱动回测}

\subsection{向量化}

把每天的收益按权重加权求和：

\begin{equation}
  r^{\text{port}}_t = \sum_i w_{i,t-1}\, r_{i,t}.
\end{equation}

\textbf{优点}：快，几行代码能跑完十年。
\textbf{缺点}：无法处理路径依赖（涨跌停、停牌、整手、资金约束）。

\subsection{事件驱动}

逐个交易日模拟：先撮合、再更新持仓与现金、再计算净值。

\textbf{优点}：能精确建模约束。
\textbf{缺点}：慢，代码复杂。

\begin{keybox}[本书的做法]
\textbf{用向量化做因子研究与快速验证，用事件驱动做最终确认。}
两者结果差异超过 20\% 时，说明约束影响很大，必须以事件驱动为准。
\end{keybox}

\section{执行假设：三个必须明确的时间点}

\begin{center}
\begin{tabular}{@{}lll@{}}
\toprule
名称 & 定义 & 常见错误 \\
\midrule
信号日 $t$ & 用 $t$ 日收盘数据算因子 & 用了 $t$ 日之后的数据 \\
成交日 $t+1$ & 次日开盘或收盘成交 & 假设 $t$ 日收盘成交 \\
估值日 & 每日收盘按收盘价估值 & 用成交价估值 \\
\bottomrule
\end{tabular}
\end{center}

本书采用：\textbf{$t$ 日收盘算信号 → $t+1$ 日开盘成交 → 持有 $h$ 日后开盘卖出}。
这条链条上任意一环提前，都会引入前视。

\begin{warnbox}[开盘成交也不完全现实]
$t+1$ 日开盘价是集合竞价的结果，你用 $t$ 日的信息下单，
理论上可以在集合竞价中成交。但如果你的信号依赖开盘价本身（例如跳空因子），
就变成了"用开盘价预测开盘价"，必须改为 $t+1$ 日收盘成交。
\end{warnbox}

\section{净值曲线的正确读法}

\subsection{四个必看的指标}

\begin{align}
  \text{年化收益} &= \left(\prod_t (1+r_t)\right)^{252/T} - 1, \\
  \text{年化波动} &= \mathrm{std}(r_t)\sqrt{252}, \\
  \text{夏普} &= \frac{\text{年化收益} - r_f}{\text{年化波动}}, \\
  \text{最大回撤} &= \min_t\left(\frac{\text{NAV}_t}{\max_{s\le t}\text{NAV}_s} - 1\right).
\end{align}

\subsection{三个容易被忽略的指标}

\begin{itemize}
  \item \textbf{Calmar} $=$ 年化收益 $/$ 最大回撤。它比夏普更贴近"能不能拿得住"。
  \item \textbf{换手率}。年换手 20 倍的策略与 2 倍的策略，成本结构完全不同。
  \item \textbf{收益的时间分布}。把收益按年份拆开，看是否集中在某一年。
        如果 80\% 的收益来自 2015 年，这个策略大概率不可复制。
\end{itemize}

\begin{warnbox}[一条净值曲线的三种读法]
同一张图，三种人会读出三种结论：
\begin{itemize}
  \item 看总收益的人："十年十倍，太好了。"
  \item 看夏普的人："夏普 1.8，不错。"
  \item 看\textbf{分年度收益 + 换手 + 回撤修复时间}的人："2015 年贡献了 70\% 的收益，
        2021 年之后基本走平，年换手 25 倍——这个策略已经失效了。"
\end{itemize}
只有第三种读法能指导决策。
\end{warnbox}

\section{基准与归因}

\subsection{选对基准}

\begin{itemize}
  \item 全市场选股：用中证全指（000985）或等权全 A；
  \item 大中盘：沪深 300；
  \item 小盘：中证 1000 或中证 2000。
\end{itemize}

用错基准会让"beta 收益"看起来像"alpha"。

\subsection{Brinson 归因}

把超额收益拆成三块：

\begin{equation}
  R_p - R_b = \underbrace{\sum_k (w_{p,k} - w_{b,k})R_{b,k}}_{\text{配置效应}}
  + \underbrace{\sum_k w_{b,k}(R_{p,k} - R_{b,k})}_{\text{选股效应}}
  + \underbrace{\sum_k (w_{p,k}-w_{b,k})(R_{p,k}-R_{b,k})}_{\text{交互效应}}.
\end{equation}

如果配置效应占了主导，说明你的"选股能力"其实是"行业择时"。

\section{一个可复现的回测清单}

在相信任何回测结果之前，逐条确认：

\begin{enumerate}
  \item 因子计算只用了 $t$ 日及之前的数据（逐列检查）。
  \item 财务数据用公告日对齐。
  \item 股票池是时点的，包含已退市股票。
  \item 价格是后复权。
  \item 涨停不可买、跌停不可卖、停牌不交易。
  \item 成本按单边 30 bp 计，并做 15--60 bp 的敏感性扫描。
  \item 样本外是 Purged K-Fold 或 Walk-Forward。
  \item 记录所有尝试过的配置数量。
  \item 分年度拆解收益，检查是否集中在某一两年。
  \item 用第二个数据源抽查 20 只股票的收盘价。
\end{enumerate}

\begin{keybox}[十条里只要有一条没做到]
结论就应该标注"待验证"，而不是"有效"。
本书第~\ref{ch:result} 章的报告，就是按这十条逐项标注的。
\end{keybox}

\section{本章小结}

\begin{enumerate}
  \item 回测不会报错，它只会给你一条漂亮的假曲线。
  \item 七个陷阱里，前视偏差与不可成交最难发现，也最致命。
  \item 向量化回测做研究，事件驱动回测做确认；两者差异大时以事件驱动为准。
  \item 信号日 / 成交日 / 估值日三个时间点必须明确且不重叠。
  \item 净值曲线要按"分年度 + 换手 + 回撤修复"来读，而不是只看总收益。
\end{enumerate}

\chapter{过拟合的数学}\label{ch:overfit}

\begin{keybox}[本章要回答的问题]
为什么样本内夏普必然虚高？Purged K-Fold 到底在切什么？Deflated Sharpe Ratio
的每一项在惩罚什么？PBO 测的是什么、又不能测什么？试了多少次之后，样本外结果就不能信了？
\end{keybox}

\section{一个必须接受的算术事实}

假设你有 $N$ 个完全无效的策略（真实夏普为 0），每个策略的样本夏普服从
标准差为 $1/\sqrt{T}$ 的分布（$T$ 为样本天数）。那么最优策略的期望夏普为

\begin{equation}
  \E\!\left[\max_{n\le N}\SR_n\right]
  \approx \frac{1}{\sqrt{T}}\left[(1-\gamma)\Phi^{-1}\!\left(1-\frac1N\right)
  + \gamma\,\Phi^{-1}\!\left(1-\frac{1}{Ne}\right)\right].
  \label{eq:emax2}
\end{equation}

取 $T=1000$、$N=100$：$\Phi^{-1}(0.99)\approx 2.33$，$\Phi^{-1}(0.9963)\approx 2.68$，
代入得 $\E[\max]\approx (0.423\times2.33 + 0.577\times2.68)/31.6 \approx 0.08$（未年化），
年化约 $1.27$。

\textbf{这意味着一百个纯噪声策略里，最好的那个年化夏普期望就有 1.27。}

\begin{warnbox}[所以]
"我这个策略夏普 1.3" 这句话，在不知道试了多少次之前，\textbf{完全没有信息量}。
本章剩下的部分，就是把这个直觉变成一个可以计算的惩罚项。
\end{warnbox}

\section{Purged K-Fold：为什么随机切分一定错}

\subsection{重叠标签导致的泄漏}

设样本 $i$ 的时间戳为 $t_i$，标签覆盖 $[t_i, t_i+h]$。若测试集是 $[t_a, t_b]$，
那么\textbf{任何} $t_i \in [t_a - h, t_b]$ 的训练样本，其标签区间都与测试集重叠。

这就是 purge 的由来：训练集必须剔除 $[t_a - h,\ t_b]$ 这段。

\subsection{Embargo：切断特征方向的反向泄漏}

purge 解决了"标签到标签"的泄漏。但还有另一个方向：
训练样本 $j$（$t_j > t_b$）的特征用了 $[t_j - w, t_j]$ 的历史数据，
若 $t_j - w < t_b$，它的特征里就包含了测试期的信息。

\begin{equation}
  \text{剔除区间} = [\,t_a - h,\ \ t_b + w\,],
  \label{eq:purge-full}
\end{equation}

其中 $h$ 是标签持有期，$w$ 是特征最长回看窗口。

实践中常用的简化：$w$ 取"测试段之后一小段"（embargo），
通常取 $h$ 或样本长度的 $1\%$--$5\%$。

\begin{center}
\begin{tabular}{@{}p{0.16\textwidth}p{0.30\textwidth}p{0.44\textwidth}@{}}
\toprule
\textbf{区域} & \textbf{是否剔除} & \textbf{原因} \\
\midrule
测试段之前 $h$ 天 & 是（purge） & 标签区间与测试段重叠 \\
测试段本身 & -- & 就是测试集 \\
测试段之后 $w$ 天 & 是（embargo） & 特征回看窗伸进测试期 \\
\bottomrule
\end{tabular}
\end{center}

\begin{pycode}[Purged K-Fold]
def purged_kfold(n, k=5, embargo=5, horizon=5):
    idx = np.arange(n)
    out = []
    for test in np.array_split(idx, k):
        t0, t1 = test[0], test[-1]
        drop = np.zeros(n, dtype=bool)
        drop[max(0, t0 - horizon): t1 + 1] = True      # purge
        drop[t1 + 1: min(n, t1 + 1 + embargo)] = True  # embargo
        out.append((idx[~drop], test))
    return out
\end{pycode}

\section{CPCV：评估"切法"本身的不确定性}

一次 Purged K-Fold 只给你一个样本外数字。但这个数字本身有方差——
换一种切法，结果可能差很多。CPCV（组合 purged 交叉验证）就是用来量化这个方差的。

\subsection{算法}

\begin{enumerate}
  \item 把时间轴等分成 $N$ 组；
  \item 每次取 $k$ 组作为测试集，其余经 purge 后作为训练集；
  \item 穷举所有 $\binom{N}{k}$ 种组合；
  \item 每个测试组会被预测多次，把它们拼起来，就得到多条完整的样本外路径。
\end{enumerate}

\begin{proposition}[路径数]
$N$ 组、每次取 $k$ 组测试，能拼出的完整样本外路径数为
\begin{equation}
  \varphi(N,k) = \frac{k}{N}\binom{N}{k} = \binom{N-1}{k-1}.
\end{equation}
例如 $N=6,k=2$：共 $\binom{6}{2}=15$ 次切分，$\varphi=\binom{5}{1}=5$ 条完整路径。
\end{proposition}

\begin{center}
\begin{tabular}{@{}ccccc@{}}
\toprule
$N$ & $k$ & 切分数 $\binom{N}{k}$ & 路径数 $\varphi$ & 训练比例 \\
\midrule
6 & 2 & 15 & 5 & 4/6 \\
8 & 2 & 28 & 7 & 6/8 \\
10 & 2 & 45 & 9 & 8/10 \\
6 & 1 & 6 & 1 & 5/6（退化为普通 K 折）\\
\bottomrule
\end{tabular}
\end{center}

\begin{keybox}[CPCV 的真正价值]
它给你的不是一个数，而是一个\textbf{分布}。
如果你的策略在 5 条路径上的夏普分别是 $1.8, 1.6, 0.2, 1.5, 1.7$，
那么"平均 1.36"是没有意义的——真实情况是"有四条路径不错，有一条几乎失效"。
这种信息在单次 K 折里完全看不到。
\end{keybox}

\section{Deflated Sharpe Ratio：给夏普打折}

\subsection{从 PSR 说起}

Bailey 与 López de Prado 提出的概率夏普比率（PSR）回答：
\textbf{在考虑偏度与峰度之后，真实夏普大于某个基准的概率是多少？}

\begin{equation}
  \widehat{\mathrm{PSR}}(\SR^*) = \Phi\!\left(
  \frac{(\widehat{\SR} - \SR^*)\sqrt{T-1}}
  {\sqrt{1 - \hat\gamma_3\widehat{\SR} + \frac{\hat\gamma_4-1}{4}\widehat{\SR}^2}}
  \right),
  \label{eq:psr}
\end{equation}

其中 $\widehat{\SR}$ 是\textbf{未年化}的样本夏普，$\hat\gamma_3$ 是偏度，
$\hat\gamma_4$ 是峰度（正态为 3）。

三个细节：

\begin{enumerate}
  \item \textbf{必须用未年化的 SR}。年化后再套公式，结果会完全错误。
  \item \textbf{偏度修正}：负偏（$\gamma_3<0$）会放大分母，降低显著性——
        因为负偏意味着"平时小赚、偶尔巨亏"，这种收益分布更可能是运气。
  \item \textbf{峰度修正}：厚尾（$\gamma_4>3$）同样降低显著性。
\end{enumerate}

\subsection{DSR：把"试了多少次"放进去}

PSR 里有一个基准 $\SR^*$。如果我们把 $\SR^*$ 设为"$N$ 次试验下，纯噪声能达到的最大夏普的期望"，
得到的概率就是 DSR：

\begin{equation}
  \SR_0 = \sqrt{\widehat{\Var}(\SR)}\left[
  (1-\gamma)\Phi^{-1}\!\left(1-\frac1N\right)
  + \gamma\,\Phi^{-1}\!\left(1-\frac{1}{Ne}\right)\right],
  \label{eq:sr0}
\end{equation}

\begin{equation}
  \boxed{\ \mathrm{DSR} = \Phi\!\left(
  \frac{(\widehat{\SR} - \SR_0)\sqrt{T-1}}
  {\sqrt{1 - \hat\gamma_3\widehat{\SR} + \frac{\hat\gamma_4-1}{4}\widehat{\SR}^2}}
  \right)\ }
  \label{eq:dsr}
\end{equation}

通常要求 $\mathrm{DSR} > 0.95$ 才算显著。

\subsection{有效试验次数怎么估}

直接把"我试了 200 次"代入 $N$ 会过度惩罚，因为 200 组参数之间高度相关。
更合理的做法是用相关矩阵的有效秩：

\begin{equation}
  N_{\text{eff}} = \frac{(\sum_i \lambda_i)^2}{\sum_i \lambda_i^2},
  \qquad \text{或} \qquad
  N_{\text{eff}} = 1 + (M-1)(1-\bar\rho),
  \label{eq:neff}
\end{equation}

其中 $\lambda_i$ 是候选策略收益相关矩阵的特征值，$\bar\rho$ 是平均相关系数。
两者取较小值（更保守）。

\begin{warnbox}[当试验次数未知时]
读别人的研报时，你通常不知道对方试了多少次。这时可以反过来问：
\textbf{这个夏普最多能承受多少次独立试验？}
令 $\mathrm{DSR}=0.95$ 反解 $\SR_0^*$，再用一阶近似 $\SR_0\approx\sigma\sqrt{2\ln N}$ 得
\begin{equation}
  N_{\max} \approx \exp\!\left[\frac12\left(\frac{\SR_0^*}{\sigma_{\SR}}\right)^2\right].
\end{equation}
如果这个数字小于"一个正常人会尝试的次数"，结论就该打问号。
\end{warnbox}

\subsection{一个数值例子}

假设某策略：$T=1000$ 个交易日，年化夏普 $1.5$（未年化 $\widehat{\SR}=0.0945$），
偏度 $-0.8$，峰度 $6$，试了 $N_{\text{eff}}=20$ 次。

\begin{itemize}
  \item $\sigma_{\SR} = 1/\sqrt{999} = 0.0316$；
  \item $\SR_0 = 0.0316\times[(0.423)(2.326) + 0.577(2.75)] = 0.0316\times 2.56 = 0.0809$；
  \item 分母 $=\sqrt{1 - (-0.8)(0.0945) + (5/4)(0.0945)^2} = \sqrt{1+0.0756+0.0112}=1.042$；
  \item $z = (0.0945-0.0809)\times31.6/1.042 = 0.41$；
  \item $\mathrm{DSR} = \Phi(0.41) = 0.66$。
\end{itemize}

\textbf{夏普 1.5 的策略，在试了 20 次之后，显著性是 0.66——远达不到 0.95。}

这就是"打折"的含义：它不是说策略不好，而是说\textbf{现有证据不足以排除运气}。

\subsection{门槛该定多高：$t>2$ 不够}

Harvey, Liu \& Zhu (2016) 在考察了数百个已发表因子之后提出：
传统的 $t>2.0$ 门槛已经不够，\textbf{新因子应要求 $t>3.0$}。
Hou, Xue \& Zhang (2020) 对 447 个已发表异象做统一复现，
发现 286 个（64\%）在 5\% 水平下不显著；把门槛提到 $t>3$，不显著的比例升到 85\%。

本书据此在附录~\ref{app:strict} 里对全部 39 个因子重做了一次判定：
31 个通过 $|t|>3$，4 个仅在 $t>2$ 水平上显著，4 个不显著。

\begin{warnbox}[别把“通过 $t>3$”当成“独立发现”]
$t>3$ 解决的是“单个因子是否为噪声”，不解决“这 31 个因子是不是同一个故事”。
本书的因子高度相关（见第~\ref{ch:factor} 章的相关性讨论），
真正独立的维度远小于 31。多重检验的惩罚要靠下面的 DSR 与 PBO，
而不是靠把门槛调高。
\end{warnbox}

\subsection{一个可以直接查的表}

把 DSR 的第一步（零假设下 $N$ 次试验最优夏普的期望）做成表，
比记公式更实用。表~\ref{tab:sr0} 是\textbf{年化夏普}口径。

\begin{table}[htbp]
\centering
\caption{纯噪声下 $N$ 次试验的最优夏普期望 $\SR_0$（年化）}
\label{tab:sr0}
\small
\begin{tabular}{@{}lrrr@{}}
\toprule
试验数 $N$ & $\mathrm{SD}(\SR)=0.2$ & $0.3$ & $0.5$ \\
\midrule
6   & 0.26 & 0.39 & 0.65 \\
10  & 0.32 & 0.47 & 0.79 \\
30  & 0.41 & 0.62 & 1.04 \\
100 & 0.51 & 0.76 & 1.27 \\
1000& 0.65 & 0.98 & 1.63 \\
\bottomrule
\end{tabular}
\end{table}

读法：如果你试了 100 个变体、各变体夏普的标准差是 0.5，
那么\textbf{纯噪声下的最优夏普期望就是 1.27}。此时你观察到一个 1.2 的夏普，
它和运气无法区分。

这张表解释了一个常见困惑：为什么很多“夏普 1.5 的完美策略”上线就亏。

\section{PBO：测的是选择过程，不是策略}

\subsection{CSCV 算法}

\begin{enumerate}
  \item 设 $M$ 个候选策略的收益矩阵为 $T\times M$；
  \item 把时间轴等分成 $S$ 块（$S$ 取偶数，常用 16）；
  \item 对每个 $\binom{S}{S/2}$ 组合：前一半块拼成"样本内"，后一半拼成"样本外"；
  \item 在样本内挑出夏普最高的策略 $n^*$；
  \item 看它在样本外的排名分位 $\omega = \mathrm{rank}/(M+1)$，取对数几率
        $\lambda = \ln\frac{\omega}{1-\omega}$；
  \item $\mathrm{PBO} = \Pr(\lambda \le 0)$。
\end{enumerate}

\begin{pycode}[PBO 的核心循环]
for combo in combinations(range(S), S // 2):
    IS  = concat(blocks[i] for i in combo)
    OOS = concat(blocks[i] for i not in combo)
    n_star = argmax(sharpe(IS))
    omega  = rank(sharpe(OOS))[n_star] / (M + 1)
    logits.append(log(omega / (1 - omega)))
pbo = mean(logit <= 0 for logit in logits)
\end{pycode}

\subsection{怎么读}

\begin{center}
\begin{tabular}{@{}ll@{}}
\toprule
PBO & 含义 \\
\midrule
$\approx 0.5$ & 样本内选出的"最优"在样本外与随机无异——纯噪声拟合 \\
$< 0.1$ & 序关系稳定，过拟合风险低 \\
$\rightarrow 1$ & 样本内越好，样本外越差——典型的"反向过拟合" \\
\bottomrule
\end{tabular}
\end{center}

\subsection{PBO 不能告诉你什么}

这是最容易被误用的一点。PBO 只衡量\textbf{相对序关系}，不衡量\textbf{绝对收益}。

\begin{warnbox}[两个反例]
\textbf{反例 A（高 PBO 但都赚钱）}：100 个策略的真实夏普都在 $[2.0, 2.1]$。
由于差异极小，样本内第一名在样本外有 50\% 概率掉到后半区，PBO 约等于 $0.5$。
但随便选哪个都赚钱。

\textbf{反例 B（低 PBO 但都亏钱）}：100 个策略全是亏的，且排名极其稳定
（亏得最少的那只永远排第一）。PBO 等于 $0$，但策略完全无效。

所以 PBO 必须与\textbf{绝对收益}一起读。
\end{warnbox}

\section{White's Reality Check 与 SPA}

Reality Check 回答的是另一个问题：\textbf{$M$ 个策略里最好的那个，是否显著优于基准？}

零假设：所有策略的真实超额收益 $\le 0$。用对时间分块的重抽样（保留策略之间的相关性）
构造 $\max$ 统计量的零分布：

\begin{equation}
  V = \max_{m}\ \sqrt{T}\,\bar f_m,
  \qquad
  p = \Pr\nolimits^{*}\!\left(V^* \ge V\right),
\end{equation}

其中 $f_m$ 是第 $m$ 个策略相对基准的收益。$p<0.05$ 才说明"最好那个"不是运气。

Hansen 的 SPA 是它的改进版，剔除了"明显很差"的策略，提升了检验功效。

\section{四个统计量的分工}

\begin{center}
\begin{longtable}{@{}p{0.16\textwidth}p{0.34\textwidth}p{0.38\textwidth}@{}}
\toprule
\textbf{统计量} & \textbf{回答的问题} & \textbf{不能回答} \\
\midrule
\endhead
PSR & 这条曲线本身显著吗 & 试了多少次 \\
DSR & 扣除多次尝试后还显著吗 & 策略之间的相对优劣 \\
PBO & 我的"挑选过程"可靠吗 & 策略是否赚钱 \\
Reality Check & 最好那个是否显著优于基准 & 单个策略的显著性 \\
\bottomrule
\end{longtable}
\end{center}

\begin{keybox}[一份诚实的研究报告应该长什么样]
\begin{enumerate}
  \item 我提出 $K$ 个候选假设，全部记录在案；
  \item 用 Purged K-Fold 与 CPCV 得到样本外表现的\textbf{分布}，而不只是一个数；
  \item 报告 DSR（含有效试验次数）、PBO、Reality Check 的 $p$ 值；
  \item 报告扣费后、分年度的收益，并指出收益集中在哪几年；
  \item 明确列出所有已知局限（数据缺失、成本假设、样本期）。
\end{enumerate}
\end{keybox}

\section{研究纪律：一张可以贴在墙上的清单}

\begin{enumerate}
  \item 先写假设，再看数据。看到数据之后提出的假设，一律按"事后假设"处理。
  \item 每试一个配置，就在日志里加一行。忘记试过什么，是过拟合的开始。
  \item 样本外只用一次。反复用同一样本外调参，它就不再是样本外。
  \item 任何 IC 超过 $0.15$ 的日频因子，先假设是数据泄漏，去找证据反驳自己。
  \item 任何夏普超过 $2$ 的日频策略，先假设是成本低估或不可成交。
  \item 结论里必须有一句"如果某某条件不成立，这个结论就失效"。
\end{enumerate}

\section{本章小结}

\begin{enumerate}
  \item 纯噪声策略的样本内夏普期望随 $\sqrt{\ln N}$ 增长；不报 $N$ 的夏普没有意义。
  \item Purge 剔除标签重叠，Embargo 切断特征反向泄漏，两者缺一不可。
  \item CPCV 给出样本外表现的\textbf{分布}，路径数为 $\binom{N-1}{k-1}$。
  \item DSR 用未年化夏普加偏度峰度修正加有效试验次数，给出一个概率。
  \item PBO 测的是选择过程的可靠性，必须与绝对收益一起读。
  \item 诚实的研究报告：全部尝试记录，加样本外分布，加四个统计量，加已知局限。
\end{enumerate}

\chapter{实证结果：一份完整的重复测试}\label{ch:result}

\begin{keybox}[本章要回答的问题]
把前九章的方法全部跑一遍，会得到什么？样本外 Rank IC 是多大？
在 15 次不同切法下它稳定吗？扣除成本、扣除多重检验惩罚之后还剩什么？
\end{keybox}

本章所有数字都来自配套仓库的一键复现流程，数据区间 \textbf{2013-01-04 至 2026-10-19}，
全 A 股日频，$h=5$ 日持有期，每 5 个交易日一个样本（避免标签重叠），
共 \textbf{2,018,278} 个样本、\textbf{670} 个周度截面、\textbf{39} 个因子。

\section{单因子：谁在真正预测收益}

下表来自主库 2010--2026 的每日 IC 序列，标签为次日开盘至第六日开盘收益。
上面的 670 个截面指模型采样，不能当作本表的样本量。
调整后 IC 是行业与对数市值中性化后的 Rank IC，按其均值排序。
{\small
\begin{longtable}{@{}lrrrrrr@{}}
\toprule
因子 & 原始 IC & 调整后 IC & 原始 $t$ & 调整后 $t$ & $n$ & $n_{\rm eff}$ \\
\midrule
\endhead
ivol\_60 & 0.0786 & 0.0729 & 12.24 & 11.38 & 2259 & 592 \\
turn\_std\_20 & 0.0693 & 0.0695 & 9.36 & 17.83 & 2035 & 514 \\
ln\_amount\_20 & 0.0725 & 0.0659 & 14.94 & 13.13 & 3891 & 830 \\
hl\_20 & 0.0590 & 0.0658 & 9.50 & 12.44 & 3891 & 950 \\
amp\_20 & 0.0581 & 0.0653 & 9.31 & 12.21 & 3891 & 949 \\
vol\_20 & 0.0571 & 0.0652 & 9.65 & 13.41 & 3874 & 941 \\
turn\_20 & 0.0636 & 0.0650 & 7.74 & 15.33 & 2035 & 515 \\
vol\_60 & 0.0560 & 0.0648 & 8.51 & 11.51 & 3466 & 905 \\
max\_ret\_20 & 0.0635 & 0.0640 & 14.46 & 18.04 & 3874 & 959 \\
vol\_ratio & 0.0401 & 0.0391 & 11.25 & 13.97 & 3483 & 912 \\
range\_60 & 0.0282 & 0.0386 & 4.44 & 7.37 & 3483 & 863 \\
turn\_chg & 0.0447 & 0.0386 & 9.59 & 13.77 & 1819 & 488 \\
pv\_corr\_20 & 0.0474 & 0.0383 & 17.71 & 16.73 & 3874 & 992 \\
rev\_10 & 0.0430 & 0.0378 & 10.18 & 10.49 & 3908 & 1138 \\
bp & 0.0322 & 0.0373 & 6.04 & 9.92 & 4076 & 854 \\
amihud\_20 & 0.0492 & 0.0358 & 10.23 & 11.52 & 3874 & 840 \\
rev\_5 & 0.0440 & 0.0314 & 11.85 & 9.89 & 3993 & 1389 \\
skew\_20 & 0.0327 & 0.0284 & 12.18 & 14.30 & 3874 & 903 \\
cfp & 0.0168 & 0.0226 & 4.83 & 10.38 & 4013 & 868 \\
ep\_ttm & 0.0171 & 0.0203 & 4.02 & 6.31 & 4013 & 882 \\
rev\_1 & 0.0172 & 0.0167 & 7.80 & 9.20 & 4061 & 3357 \\
sp\_ttm & 0.0207 & 0.0153 & 4.04 & 6.77 & 2134 & 505 \\
beta\_60 & 0.0036 & 0.0143 & 0.49 & 2.49 & 2854 & 764 \\
size & 0.0167 & 0.0123 & 2.27 & 2.03 & 2134 & 537 \\
div\_yield & -0.0009 & 0.0100 & -0.11 & 0.95 & 1812 & 465 \\
ocf\_to\_eps & 0.0081 & 0.0090 & 4.88 & 9.61 & 4013 & 895 \\
rev\_accel & 0.0117 & 0.0081 & 7.85 & 5.53 & 3012 & 709 \\
roe & 0.0049 & 0.0079 & 1.22 & 2.10 & 4030 & 848 \\
np\_yoy & 0.0124 & 0.0068 & 4.66 & 2.76 & 4037 & 790 \\
rev\_yoy & 0.0084 & 0.0060 & 2.94 & 2.44 & 4037 & 807 \\
accrual & 0.0062 & 0.0056 & 3.48 & 4.21 & 4013 & 863 \\
kurt\_20 & 0.0050 & 0.0029 & 2.13 & 1.64 & 3874 & 936 \\
gross\_margin & 0.0037 & 0.0017 & 1.12 & 0.70 & 4046 & 860 \\
gap\_20 & -0.0118 & -0.0176 & -4.77 & -9.18 & 3891 & 963 \\
mom\_12\_1 & -0.0687 & -0.0179 & -2.07 & -0.68 & 48 & 25 \\
resid\_mom\_60 & -0.0431 & -0.0360 & -7.91 & -7.36 & 2854 & 694 \\
mom\_60 & -0.0452 & -0.0372 & -7.94 & -7.58 & 3058 & 753 \\
mom\_120 & -0.0487 & -0.0405 & -7.24 & -6.41 & 2038 & 514 \\
mom\_20 & -0.0548 & -0.0495 & -11.74 & -12.28 & 3738 & 1018 \\
\bottomrule
\end{longtable}
}

原始 IC 与调整后 IC 均由完整日频 IC 序列汇总；$h=5$ 的日频标签重叠，
均值检验使用 Bartlett Newey--West 协方差，滞后阶数为
$L=\max(4,\lfloor4(n/100)^{2/9}\rfloor)$。$n$ 是有值的 IC 日期数，
$n_{\rm eff}=n\gamma_0/\widehat\Omega$ 是原始 IC 均值方差等价的样本量，
不是独立财报次数，也不是未经假设可确定的真实样本数。
完整表另含调整后有效样本量；不同带宽及每五日采样的敏感性见配套 CSV。
本表不再平均各年的 $t$ 值。

\subsection{结论的适用范围}
流动性、波动和量价因子在这段数据内呈现较强关联；动量的符号需按具体窗口与时期解读。
质量与成长的旧表使用年度统计量均值，不能据此宣称“横截面没有信息”，
也不能据此认定它们更适合时间序列筛选。补充质量因子的重算见附录~\ref{app:quality}。
ICIR 是描述性比值；重叠收益下的年化 ICIR 不能代替均值检验。
全部因子均经过研究者筛选与查看，单个显著 $t$ 值仍需多重检验和未见样本验证。

\section{模型：等权基线 vs 线性 vs GBDT}

三条基线：等权合成（用训练集 IC 符号定向）、Ridge、LightGBM。
协议表中的 NW $t$ 是折内 $t$ 的均值，仅作描述性对照，不是合并样本显著性。
各折模型仅用该折训练期拟合；但研究者已看过整个历史区间，因子选择和研究决策未被这些切分隔离。

\subsection{三种验证协议的对比}

\begin{table}[htbp]
\centering
\caption{三种验证协议下的样本外表现（多次切分的均值）}
\label{tab:protocol}
\small
\begin{tabular}{@{}llrrrrrr@{}}
\toprule
协议 & 模型 & 折数 & IC 均值 & ICIR & NW $t$ & 扣费年化 & 夏普 \\
\midrule
Purged K-Fold & 等权合成 & 5 & 0.0495 & 3.57 & 5.65 & 5.1\% & 0.65 \\
Purged K-Fold & Ridge & 5 & 0.0407 & 3.32 & 5.12 & 3.7\% & 0.56 \\
Purged K-Fold & LightGBM & 5 & 0.0500 & 3.35 & 5.05 & 7.9\% & 0.76 \\
\midrule
CPCV & 等权合成 & 15 & 0.0521 & 3.15 & 5.91 & 4.7\% & 0.51 \\
CPCV & Ridge & 15 & 0.0494 & 2.95 & 5.11 & 5.6\% & 0.52 \\
CPCV & LightGBM & 15 & 0.0590 & 3.22 & 5.47 & 11.8\% & 0.83 \\
\midrule
Walk-Forward & 等权合成 & 7 & 0.0747 & 5.52 & 5.74 & 13.8\% & 1.24 \\
Walk-Forward & Ridge & 7 & 0.0700 & 6.09 & 6.54 & 12.9\% & 1.34 \\
Walk-Forward & LightGBM & 7 & 0.0786 & 5.76 & 6.23 & 15.5\% & 1.08 \\
\bottomrule
\end{tabular}
\end{table}

\begin{keybox}[这张表最重要的一件事]
\textbf{同一批因子、同一批模型，换一种切法，样本外 IC 从 0.05 变到 0.079。}
这不是模型不稳定，而是"样本外"这个概念本身就有方差。
报告单一数字而不报告切法，等于隐瞒了一半的信息。
\end{keybox}

Walk-Forward 的数字系统性高于 CPCV 的原因不难理解：Walk-Forward 的训练窗口是扩展的，
后期模型见过更多历史；而 CPCV 的某些组合要求用"未来"的块去训练"过去"的块，
对非平稳的 A 股市场更不友好。\textbf{两者都不是错的，它们回答的是不同的问题}：
CPCV 问"这个模型在任意时间段上是否稳定"，Walk-Forward 问"如果我按时间顺序做实盘会怎样"。

\subsection{全样本外表现}

把 Walk-Forward 各段拼接后，各候选的有值日期数不同。收益和诊断不得把缺失预测视为现金头寸；
诊断只使用 \AuditCommonN{} 个共同有效日期。下面各候选收益按各自有值日期计算：

\begin{center}\small
\begin{tabular}{@{}lrrrrrr@{}}
\toprule
候选 & IC 日期数 & IC & ICIR & NW $t$ & 年化 & 夏普 \\
\midrule
lgbm\_s0 & 402 & 0.0785 & 5.57 & 12.74 & 16.3\% & 1.43 \\
lgbm\_s2 & 402 & 0.0785 & 5.53 & 12.62 & 15.6\% & 1.38 \\
lgbm\_s1 & 402 & 0.0780 & 5.50 & 12.58 & 15.3\% & 1.35 \\
composite & 417 & 0.0653 & 3.93 & 9.74 & 10.2\% & 0.86 \\
ridge & 402 & 0.0623 & 3.96 & 8.95 & 12.0\% & 1.07 \\
\bottomrule
\end{tabular}
\end{center}

\begin{warnbox}[怎么读这个 16.3 个百分点]
\textbf{第一，这是多空组合}：做多前 20\%、做空后 20\%，个人投资者通常无法执行空头。
\textbf{第二，它假设可以按次日开盘价成交}，且只在"可买"样本里建仓，
但没有模拟资金规模、冲击成本的市场深度、以及融券成本。
\textbf{第三，它没有做行业中性约束}，组合的行业暴露可能很大。

所以这个数字应该被理解为"\textbf{在理想执行假设下，因子组合的横截面 alpha}"，
而不是"可以拿到的收益"。
\end{warnbox}

\subsection{模型的边际贡献有多大}

值得注意的一点：LightGBM 的 IC（0.0785）比等权基线（0.0653）高约 20\%，
但它的换手也高了 40\%（37.9\% vs 27.2\%）。\textbf{扣费后的优势（1.43 vs 0.86 的夏普）
有相当一部分来自 IC 提升，也有一部分来自它更集中地下注在有效因子上。}

这与第~\ref{ch:model} 章的判断一致：模型有效，但它的作用是"更好地组合已知信息"，
而不是"发现新信息"。用等权基线做一个 sanity check，永远值得。


\begin{figure}[htbp]
\centering
\includegraphics[width=\textwidth]{fig_nav.png}
\caption{Walk-Forward 全样本外多空组合净值（扣单边 30 bp）。三条 LightGBM 曲线的差异仅来自随机种子，说明模型本身足够稳定。}
\label{fig:nav}
\end{figure}

\begin{figure}[htbp]
\centering
\includegraphics[width=0.7\textwidth]{fig_cpcv.png}
\caption{CPCV 15 次切分下的样本外 Rank IC 分布。箱线图的宽度就是"切法的不确定性"。}
\label{fig:cpcv}
\end{figure}

\begin{figure}[htbp]
\centering
\begin{minipage}{0.48\textwidth}
\centering
\includegraphics[width=\textwidth]{fig_pbo.png}
\caption{PBO 的对数几率分布。$\lambda\le 0$ 的比例即 PBO。}
\label{fig:pbo}
\end{minipage}\hfill
\begin{minipage}{0.48\textwidth}
\centering
\includegraphics[width=\textwidth]{fig_cost.png}
\caption{成本敏感性：夏普随单边成本下降，约 70 bp 处穿过零点。}
\label{fig:cost}
\end{minipage}
\end{figure}
\section{过拟合诊断}

现有搜索日志不完整，不能用五个已保存候选代表全部研究自由度。
旧版 DSR 以 $N=15$ 计算，并错误标注成“由相关矩阵估计”；这里撤回“全部通过”的判定。

\begin{center}\small
\begin{tabular}{@{}rrr@{}}
\toprule
假定独立试验数 $N$ & 五候选离散度 & 零假设标准差下限情景 \\
\midrule
15 & 0.9966 & 0.9874 \\
50 & 0.9905 & 0.9592 \\
100 & 0.9847 & 0.9320 \\
200 & 0.9770 & 0.8962 \\
500 & 0.9632 & 0.8359 \\
603 & 0.9598 & 0.8220 \\
1000 & 0.9499 & 0.7819 \\
\bottomrule
\end{tabular}
\end{center}

第一列 DSR 使用五候选的夏普横截面标准差；第二列把该标准差至少设为
$1/\sqrt{T-1}$，作为独立零均值收益近似的压力情景。
$N=603=(39+8+20)\times3\times3$ 只是名义搜索范围情景，不是执行过的配置数量或独立次数估计。
该情景 DSR 为 \AuditDSRStress{}；用五候选离散度且 $N=100$ 时为 \AuditDSRHundred{}，
旧文“约 0.93”未由代码支持，已撤回。原始 $N=15$、补零收益下的 0.997 仅保留在修订备份中。

\subsection{PBO 与区间回答的范围}
当前 PBO 为 \AuditPBO{}，仅覆盖五候选及 \AuditCommonN{} 个共同有效日期。
候选收益相关系数介于 \AuditCorrMin{} 与 \AuditCorrMax{}；三个 LightGBM 随机种子尤其接近。
它不能替代完整因子筛选、参数搜索和协议选择的 PBO。相关性高并不意味着 PBO 必然按固定方向变化。

分块 Bootstrap 的 95\% 夏普区间为 $[\AuditBPLower,\ \AuditBPUpper]$，
使用 2.5\% 和 97.5\% 分位；旧版 5\% 至 95\% 实际是 90\% 区间。
简化 RC 仅针对已保存五候选；400 次重抽样中零次超越不等于真实 $p=0$，也不足以报告 $p<0.001$。

\begin{warnbox}[历史内部验证不能代替前瞻证据]
这些诊断没有解决退市股遗漏、未知股票池筛选、收益自相关、未见样本外和执行约束。
30\,bp 也不是按成交规模与流动性校准的成本。没有完整搜索记录时，不再作“全研究稳健通过”的判定。
\end{warnbox}

\section{成本敏感性}

把单边成本从 0 扫到 80 bp，看最优候选的年化夏普如何变化：

\begin{center}
\begin{tabular}{@{}lrrrrrrr@{}}
\toprule
单边成本 & 0 bp & 20 bp & 30 bp & 40 bp & 50 bp & 60 bp & 80 bp \\
\midrule
最优候选夏普 & 2.53 & 1.80 & 1.43 & 1.07 & 0.71 & 0.34 & $-0.27$ \\
五候选均值 & 2.19 & 1.54 & 1.22 & 0.89 & 0.57 & 0.24 & $-0.40$ \\
\bottomrule
\end{tabular}
\end{center}

\textbf{盈亏平衡点大约在 70 bp。}这意味着：

\begin{itemize}
  \item 如果实际成本能控制在 30 bp 以内（佣金万二五 + 印花税 + 适度滑点），
        策略仍有可观空间；
  \item 如果滑点与冲击把成本推到 60 bp 以上（小市值、高换手、资金量偏大），
        策略基本失效。
\end{itemize}

\begin{keybox}[成本是最后一道、也是最硬的裁判]
所有前期的努力——数据清洗、因子构造、模型调优——最终都要经过这道关卡。
一个 IC 0.08 的策略在 80 bp 成本下归零，而一个 IC 0.05 但换手只有 5\% 的策略
可能活得更久。\textbf{这就是为什么第~\ref{ch:eval} 章要求"先算成本再看 IC"。}
\end{keybox}

\section{指数增强：把多空 alpha 放回可执行的框架}

前九节的所有收益都是\textbf{多空}的。但 A 股个股融券基本不可行
（见第~\ref{ch:portfolio} 章），负向信号只能"低配"。
这一节把同一批样本外预测放进指数增强框架，回答"它到底能不能用"。

\subsection{构造}

\begin{itemize}
  \item \textbf{基准}：可交易全市场的流通市值加权组合（不用官方指数，
        因为本地只有单期成分快照，用当期成分回测历史会引入成分股幸存者偏差）；
  \item \textbf{偏离}：行业内标准化打分后按
        $w_i = W_g \cdot \dfrac{\omega_i e^{\kappa z_i}}{\sum_{j\in g}\omega_j e^{\kappa z_j}}$
        构造，$\omega_i$ 是行业内基准权重，$W_g$ 是行业基准权重。
        该形式保证\textbf{行业权重与基准完全一致}、权重恒正；
  \item \textbf{$\kappa$ 扫描}：$\kappa$ 越大偏离越激进，跟踪误差越大；
  \item \textbf{成本}：单边 30\,bp，按实际换手扣。
\end{itemize}

样本为 2018-01 至 2025-11 的 383 个周度调仓期（受估值快照起点限制）。

\subsection{结果：倾得越狠，跑得越差}

\begin{tabular}{@{}rrrrrr@{}}
\toprule
$\kappa$ & 跟踪误差 & 年化超额 & 信息比率 & 超额最大回撤 & 换手 \\
\midrule
0.0 & 0.1\% & -0.9\% & -7.33 & -5.7\% & 3\% \\
0.2 & 0.4\% & -0.9\% & -1.98 & -6.0\% & 4\% \\
0.5 & 0.9\% & -1.6\% & -1.67 & -11.1\% & 9\% \\
1.0 & 1.6\% & -2.8\% & -1.67 & -19.2\% & 16\% \\
2.0 & 3.0\% & -4.6\% & -1.50 & -29.8\% & 27\% \\
3.0 & 4.0\% & -5.9\% & -1.41 & -36.2\% & 34\% \\
5.0 & 4.1\% & -7.0\% & -1.62 & -41.6\% & 37\% \\
\bottomrule
\end{tabular}


\begin{warnbox}[这张表推翻了什么]
跟踪误差从 $0.1\%$ 升到 $4.0\%$，年化超额从 $-0.9\%$ 一路跌到 $-6.4\%$。
\textbf{单调恶化，没有最优点。}信息比率始终在 $-1.0$ 到 $-1.6$ 之间。

也就是说：把本章前面那个 IC $0.079$、夏普 $1.43$ 的信号
按行业中性放进指数增强框架，\textbf{它不但没有超额，还系统性地输给基准}。
\end{warnbox}

\subsection{为什么会这样}

多空组合赚的是\textbf{横截面价差}：做多低换手、低波、近期下跌的股票，
做空高换手、高波、近期上涨的股票。这个价差在 A 股是真实的（IC 稳定为正）。

但在指数增强里，你只能表达\textbf{多头那一端}：
把权重从"贵、热、涨得多"的股票挪到"便宜、冷、跌得多"的股票。
而在 2018--2025 这段样本里，被挪走的那一端（大市值成长与热门股）\textbf{涨得更多}。
于是空头端贡献的那部分 alpha 消失了，多头端反而成了负贡献。

\begin{keybox}[这是本项目最重要的一条负面结论]
\textbf{多空口径的夏普 1.43 不可变现。}
但不等于"因子无效"——下一节会把 alpha 拆成多头端与空头端，
结论比"全靠空头"更微妙：\textbf{两端都有 alpha，但大市值那一段的多头端不显著}。

这条结论与附录~\ref{app:strict}（规模因子只在最小 30\% 市值里有效）
指向同一个事实：\textbf{这个市场的"可测"与"可做"之间有一道很宽的沟。}
\end{keybox}

\subsection{这一节的局限}

\begin{enumerate}
  \item 只测了一个基准（全市场市值加权）。若换成中证 1000 这类小盘基准，
        结论可能不同——但小盘基准的容量问题更严重；
  \item 偏离按基准权重的比例放大（$\propto w_b$），这会系统性压制中小盘的主动权重，
        真实的指数增强会用风险模型直接优化主动权重，可能表现更好；
  \item 未模拟整手、资金约束与涨跌停的路径依赖。
\end{enumerate}

结论应读作"\textbf{在这个构造与这个基准下，long-only 版本不可行}"，
而不是"任何 long-only 实现都不可行"。

\section{alpha 到底在哪一端：多头与空头的分解}

上一节的失败可以有两种解释：（a）alpha 全在空头端；（b）alpha 在两端，
但实现方式没抓住。这一节直接把它拆开。

定义（$S_{top}$ 为打分最高 20\%，$S_{bot}$ 为最低 20\%）：

\begin{equation}
  \underbrace{\bar r_{top} - \bar r_{bot}}_{\text{价差}}
  \;=\;
  \underbrace{(\bar r_{top} - \bar r_{\text{基准}})}_{\text{多头端超额}}
  \;+\;
  \underbrace{(\bar r_{\text{基准}} - \bar r_{bot})}_{\text{空头端超额}} .
\end{equation}

\subsection{结果：两端都有 alpha}

\begin{tabular}{@{}lrrrr@{}}
\toprule
项目 & 期数 & 年化 & NW $t$ & 胜率 \\
\midrule
多空价差 & 402 & +30.9\% & +6.67 & 64\% \\
多头端超额(等权) & 402 & +10.7\% & +5.25 & 68\% \\
空头端超额(等权) & 402 & +18.5\% & +6.75 & 65\% \\
多头端超额(市值加权) & 368 & +16.7\% & +3.35 & 60\% \\
空头端超额(市值加权) & 368 & +13.5\% & +2.30 & 54\% \\
\bottomrule
\end{tabular}


多头端年化 $+10.7\%$（$t=5.25$），空头端 $+18.5\%$（$t=6.75$）。
\textbf{多头端不是"没有 alpha"，只是比空头端弱一些。}

\begin{warnbox}[所以我上一节的解释只对了一半]
上一节说"alpha 主要活在空头端"，这是不准确的。
准确的说法要到下一张表才看得清。
\end{warnbox}

\subsection{按市值三分组：大市值的多头端不显著}

\begin{tabular}{@{}lrrrr@{}}
\toprule
市值组 & 多头端超额 & NW $t$ & 组内价差 & NW $t$ \\
\midrule
小 & +12.0\% & +5.53 & +45.1\% & +8.66 \\
中 & +10.3\% & +4.81 & +36.6\% & +6.88 \\
大 & +3.1\% & +1.27 & +11.5\% & +2.09 \\
\bottomrule
\end{tabular}


\begin{keybox}[拼图到这里才完整]
\begin{itemize}
  \item \textbf{小市值组}：多头端 $+12.0\%$（$t=5.53$），组内价差 $+45.1\%$；
  \item \textbf{中市值组}：多头端 $+10.3\%$（$t=4.81$），组内价差 $+36.6\%$；
  \item \textbf{大市值组}：多头端 \textbf{$+3.1\%$（$t=1.27$，不显著）}，组内价差 $+11.5\%$。
\end{itemize}

指数增强的基准是\textbf{市值加权}的，权重被大市值股票主导；
而大市值段的多头端恰好\textbf{没有统计显著的 alpha}。
于是上一节那个"倾得越狠、跑得越差"的结果就不难理解了：
你把权重往一个在大盘股上无效的信号方向挪，挪得越多，偏离越大，亏得越多。
\end{keybox}

\subsection{三条结论}

\begin{enumerate}
  \item \textbf{排序能力真实存在}，价差年化 $30.9\%$（$t=6.67$），两端都贡献，胜率 $64\%$。
  \item \textbf{alpha 随市值单调衰减}：组内价差从小盘的 $45.1\%$ 掉到大盘的 $11.5\%$，
        多头端在大盘股上直接不显著（$t=1.27$）。
  \item \textbf{可执行性与 alpha 强度成反比}：alpha 最强的地方（微盘）
        正是容量最小、监管最敏感、数据最难拿的地方。
\end{enumerate}

\begin{keybox}[对个人研究者的含义]
如果你的资金只能做流动性较好的中大盘股，那么这个信号对你的价值接近于零——
不是因为它不真实，而是因为它在你能触及的那部分市场里不存在。

这条结论比"因子有效/无效"更有用，因为它直接告诉你\textbf{该不该做}，
而不只是"能不能测出来"。
\end{keybox}

\section{本章小结}
\begin{enumerate}
  \item 主因子表改为从完整每日 IC 序列汇总，报告样本数、HAC $t$ 与方差等价有效样本量。
  \item 模型的内部验证呈现正的预测关联；没有额外未见样本，不能把它直接解释为未来收益。
  \item DSR 以试验次数与离散度敏感性呈现；PBO 仅覆盖现有五候选，完整研究筛选可靠性未知。
  \item 多空、开盘成交、统一成本与向量化执行仍是假设；long-only 和事件驱动组合尚未完成。
  \item 正式滚动窗口、数据覆盖和成本分层审计仍待补齐，不能把它们一律当作不影响数字的文字局限。
\end{enumerate}

\chapter{监控与失效判断}\label{ch:monitor}

\begin{keybox}[本章要回答的问题]
策略上线后（哪怕只是"自己跟着做"），怎么知道它正在失效？哪些指标要每天看，
哪些每月看？什么时候应该果断停手？
\end{keybox}

\section{为什么"失效"是可观测的}

一个 alpha 因子失效，通常会经历三个阶段：

\begin{enumerate}
  \item \textbf{IC 缓慢下降}。拥挤度上升，边际信息被套利掉。
  \item \textbf{IC 波动放大}。参与者结构变化，因子的收益变得不稳定。
  \item \textbf{IC 反向}。过度拥挤导致"抢跑—反转"，因子变成负 alpha。
\end{enumerate}

这三个阶段都是\textbf{可观测量}。问题在于，多数人只在净值回撤时才想起来检查。

\section{四个层级的监控指标}

\subsection{每日：信号层}

\begin{itemize}
  \item \textbf{当日 IC}（滚动 20 日均值）。与历史均值比较，跌破 2 个标准差要警惕。
  \item \textbf{因子覆盖率}。当天有多少比例的股票有有效因子值。
        覆盖率突然下降，通常是数据管道出了问题。
  \item \textbf{分数分布}。均值、标准差、偏度是否稳定。分布形状突变说明市场结构变了。
\end{itemize}

\subsection{每周：组合层}

\begin{itemize}
  \item \textbf{实际换手 vs 预期换手}。实际换手远高于预期，说明信号不稳定。
  \item \textbf{行业与市值暴露}。中性化没做好时，暴露会漂移。
  \item \textbf{持仓集中度}。前十大权重占比是否异常上升。
\end{itemize}

\subsection{每月：绩效层}

\begin{itemize}
  \item \textbf{滚动 12 个月 ICIR}。
  \item \textbf{滚动 12 个月夏普}，以及与历史分布的对比。
  \item \textbf{分年度收益}，检查是否出现"某一年特别好、其余年份平庸"的形态。
\end{itemize}

\subsection{每季：结构层}

\begin{itemize}
  \item \textbf{因子相关性矩阵}是否变化（原来的正交因子变得高度相关，说明市场结构变了）。
  \item \textbf{特征重要性}是否稳定（树模型的重要性排序大幅变化，说明模型在追逐噪声）。
  \item \textbf{因子拥挤度}代理指标（同类型产品的规模、因子多空组合的估值差）。
\end{itemize}

\section{因子拥挤度：一个可以近似估计的指标}

拥挤度没有直接数据，但有几个可用的代理：

\begin{enumerate}
  \item \textbf{因子多空组合的估值差}。例如价值因子的多空组合 PB 差距扩大到历史极值，
        说明"便宜的被买贵了"。
  \item \textbf{因子收益的波动率}。拥挤的因子在踩踏时波动会骤升。
  \item \textbf{换手率的整体抬升}。全市场换手率上升往往伴随量化资金流入。
  \item \textbf{因子 IC 的自相关}。拥挤时 IC 会变得"黏"，然后突然反转。
\end{enumerate}

\begin{warnbox}[拥挤度的悖论]
拥挤度指标最有用的时候，恰恰是它看起来"还没到极值"的时候。
等到所有人都确认"这个因子已经拥挤"，反转往往已经发生了一半。
所以拥挤度更适合作为\textbf{减仓的理由}，而不是\textbf{加仓的时机}。
\end{warnbox}

\section{风格暴露与回撤归因}

当策略回撤时，先分清是三种情况中的哪一种：

\begin{center}
\begin{tabular}{@{}lll@{}}
\toprule
类型 & 特征 & 应对 \\
\midrule
市场 beta & 与指数同步下跌 & 正常，不必干预 \\
风格暴露 & 小盘/成长整体回撤 & 检查中性化是否失效 \\
alpha 衰减 & 中性化后仍持续跑输 & 认真考虑停手 \\
\bottomrule
\end{tabular}
\end{center}

判断方法：把策略收益对风格因子（规模、价值、动量、波动）做回归，
看残差（真正的 alpha）是否还在。

\begin{equation}
  r^{\text{port}}_t - r^f_t = \alpha + \sum_k \beta_k f_{k,t} + \varepsilon_t.
\end{equation}

若 $\hat\alpha$ 的滚动估计持续为负，说明 alpha 真的消失了；
若 $\hat\alpha$ 稳定但 $\beta$ 漂移，说明只是风格暴露在起作用。

\section{什么时候该停手}

给三条可执行的规则：

\begin{enumerate}
  \item \textbf{滚动 12 个月 IC 的 $t$ 值跌破 1}，且持续两个季度。
  \item \textbf{扣费后滚动 12 个月收益为负}，且回撤超过历史最大回撤的 1.5 倍。
  \item \textbf{因子拥挤度指标进入历史前 10\% 分位}，同时 IC 开始反向。
\end{enumerate}

\begin{keybox}[最重要的一条]
\textbf{在策略设计阶段就把"停手条件"写下来。}
事后判断永远会被"再等等看"的心理绑架。事前写好的规则，
是唯一能在情绪最差时保护你的东西。
\end{keybox}

\section{本章小结}

\begin{enumerate}
  \item 因子失效有三个可观测阶段：IC 下降、波动放大、IC 反向。
  \item 监控要分四层：每日信号、每周组合、每月绩效、每季结构。
  \item 拥挤度可用估值差、波动率、换手率等代理指标近似估计。
  \item 回撤要先分清是 beta、风格暴露还是 alpha 衰减。
  \item 停手规则必须在策略设计阶段就写下来。
\end{enumerate}

\chapter{个人研究者的十二周路线图}\label{ch:roadmap}

\begin{keybox}[本章要回答的问题]
如果只有一台普通电脑、每天两小时、三个月时间，应该按什么顺序做，
才能在每一步都拿到\textbf{可信}的结论？
\end{keybox}

\section{总原则：每一步都要有"可以继续"的判据}

研究最怕的不是做不完，而是做完了却不知道结论可不可信。因此每一阶段都必须设定
\textbf{量化判据}：达不到就不进入下一步，达到了才继续。

\section{第一阶段（第 1--3 周）：数据地基}

\subsection{做什么}

\begin{enumerate}
  \item 建立本地数据库（行情 + 财报 + 估值 + 日历）；
  \item 实现复权、PIT、可交易性标记；
  \item 写数据质量检查脚本。
\end{enumerate}

\subsection{判据}

\begin{itemize}
  \item 每个交易日的股票数量单调变化，无整段缺失；
  \item 复权因子单调不减，除权日收益率不超涨跌停；
  \item 所有财报的公告日晚于报告期；
  \item 可交易股票占比在 70\%--90\%；
  \item 随机抽 20 只股票，与第二个数据源核对收盘价，误差 $<0.1\%$。
\end{itemize}

\begin{warnbox}[最常见的失败模式]
急着算因子，跳过数据校验。结果是：后面所有结论都建立在一堆看不见的错误上，
而且在发现问题之前，你已经花了三周时间。
\end{warnbox}

\section{第二阶段（第 4--6 周）：因子与评价}

\subsection{做什么}

\begin{enumerate}
  \item 实现 15--20 个因子（先做反转、波动、流动性三个家族）；
  \item 实现去极值、标准化、中性化；
  \item 实现 IC / ICIR / 分层 / 换手 / 衰减评价；
  \item 做因子相关性聚类。
\end{enumerate}

\subsection{判据}

\begin{itemize}
  \item 至少 5 个因子的中性化后 $|{\rm IC}| > 0.02$ 且 NW $t > 2$；
  \item 至少 3 个因子的分层单调性 $|{\rm corr}| > 0.8$；
  \item 因子相关性矩阵中，没有超过一半的因子对相关性 $>0.7$。
\end{itemize}

\section{第三阶段（第 7--9 周）：标签、模型与验证}

\subsection{做什么}

\begin{enumerate}
  \item 构造横截面排序标签；
  \item 实现等权合成基线；
  \item 实现 Ridge 与 LightGBM；
  \item 实现 Purged K-Fold 与 CPCV；
  \item 实现 DSR 与 PBO。
\end{enumerate}

\subsection{判据}

\begin{itemize}
  \item 样本外 Rank IC 在 5 折上的标准差小于均值的 50\%（结果稳定）；
  \item CPCV 的多条路径中，至少 70\% 的夏普为正；
  \item DSR $>0.9$；
  \item PBO $<0.3$。
\end{itemize}

\begin{warnbox}[如果达不到]
不要换更复杂的模型，先回头检查：
标签对不对？有没有前视？成本算没算？股票池合不合理？
90\% 的"模型不行"其实是"数据或标签不行"。
\end{warnbox}

\section{第四阶段（第 10--12 周）：组合与报告}

\subsection{做什么}

\begin{enumerate}
  \item 实现分位数组合与带约束优化；
  \item 实现成本模型与容量分析；
  \item 写完整研究报告（含所有已知局限）；
  \item 建立监控脚本。
\end{enumerate}

\subsection{判据}

\begin{itemize}
  \item 扣 30 bp 成本后仍有正的年化收益；
  \item 分年度看，收益不集中在某一年；
  \item 换手率与成本假设匹配（年换手 $\times$ 成本 $<$ 年化 alpha 的 50\%）；
  \item 报告里明确写出"这个结论在什么条件下会失效"。
\end{itemize}

\section{时间分配的现实建议}

\begin{center}
\begin{tabular}{@{}lcc@{}}
\toprule
阶段 & 理想占比 & 现实占比（多数人） \\
\midrule
数据 & 30\% & 10\% \\
因子 & 25\% & 15\% \\
模型 & 15\% & 50\% \\
验证 & 20\% & 5\% \\
组合与报告 & 10\% & 20\% \\
\bottomrule
\end{tabular}
\end{center}

\textbf{现实与理想的差距，基本都花在了"调模型"上。}而调模型恰恰是回报最低的部分。

\section{一份最小可用的文件结构}

\begin{center}
\begin{tabular}{@{}ll@{}}
\toprule
目录 & 内容 \\
\midrule
\texttt{data/quant.db} & SQLite 原始库 \\
\texttt{data/factors/} & 按年分片的因子面板（parquet） \\
\texttt{aq/} & 框架代码（面板、因子、评价、CV、统计） \\
\texttt{scripts/} & 01--09 编号脚本，按顺序可一键复现 \\
\texttt{reports/} & 因子评价表、模型测试表、诊断 JSON、图表 \\
\texttt{lectures/} & 本讲义源码与 PDF \\
\bottomrule
\end{tabular}
\end{center}

\begin{keybox}[可复现的最后一道保险]
把 \texttt{scripts/} 里的编号脚本从头到尾跑一遍，能产出与报告完全一致的数字。
如果做不到，说明你的研究里有"手工操作"——而手工操作是复现性最大的敌人。
\end{keybox}

\section{本章小结}

\begin{enumerate}
  \item 每一阶段都要有量化判据，达不到不进入下一步。
  \item 数据阶段应该占 30\% 的时间，现实中多数人只花 10\%。
  \item 模型阶段最容易超时，而它的边际回报最低。
  \item 90\% 的"模型不行"其实是"数据或标签不行"。
  \item 一键复现是研究可信度的最后一道保险。
\end{enumerate}

\appendix
\chapter{数据库字段字典}\label{app:db}

\begin{longtable}{@{}p{0.14\textwidth}p{0.16\textwidth}p{0.62\textwidth}@{}}
\toprule
\textbf{字段} & \textbf{类型} & \textbf{含义与口径} \\
\midrule
\endhead
\multicolumn{3}{@{}l}{\textbf{stock} —— 证券主表，主键 \texttt{code}}\\
\texttt{code} & TEXT & 6 位代码，如 \texttt{600000} \\
\texttt{secid} & TEXT & 行情标识，如 \texttt{1.600000} \\
\texttt{exchange} & TEXT & \texttt{SH} / \texttt{SZ} / \texttt{BJ} \\
\texttt{name} & TEXT & 最新证券简称 \\
\texttt{board} & TEXT & 东财行业板块名 \\
\texttt{list\_date} & TEXT & 上市日（由首个交易日推断） \\
\texttt{delist\_date} & TEXT & 退市日（仍在市为 NULL） \\
\texttt{is\_active} & INT & 最近 45 天是否有成交 \\
\midrule
\multicolumn{3}{@{}l}{\textbf{bar} —— 日线行情，主键 \texttt{(code, date)}}\\
\texttt{open/high/}\newline\texttt{low/close} & REAL & \textbf{不复权}原始价（元） \\
\texttt{pre\_close} & REAL & 前收盘（不复权） \\
\texttt{volume} & REAL & 成交量（股） \\
\texttt{amount} & REAL & 成交额（元）。新浪源无此字段，
由「成交量 $\times$ 当日均价」估算，误差通常 $<2\%$ \\
\texttt{turnover} & REAL & 基础表该字段几乎为空；研究面板按月度股本快照前向填充后计算换手率（\%） \\
\texttt{pct\_chg} & REAL & 涨跌幅（\%） \\
\texttt{adj\_factor} & REAL & 累计后复权因子；后复权价 $=$ 原始价 $\times$ \texttt{adj\_factor} \\
\texttt{is\_st} & INT & 保留字段（本书在面板层由简称历史实时生成） \\
\texttt{trade\_status} & INT & 保留字段（停牌体现为「缺行」） \\
\midrule
\multicolumn{3}{@{}l}{\textbf{fin} —— 财报（point-in-time），主键 \texttt{(code, report\_date)}}\\
\texttt{notice\_date} & TEXT & \textbf{公告日}，PIT 对齐的唯一依据 \\
\texttt{eps} & REAL & 基本每股收益（累计） \\
\texttt{bps} & REAL & 每股净资产 \\
\texttt{roe} & REAL & 加权平均 ROE（\%） \\
\texttt{revenue} & REAL & 营业总收入（元，累计） \\
\texttt{net\_profit} & REAL & 归母净利润（元，累计） \\
\texttt{ocf\_ps} & REAL & 每股经营现金流（累计） \\
\texttt{gross\_margin} & REAL & 销售毛利率（\%） \\
\texttt{div\_yield} & REAL & 股息率（\%） \\
\texttt{rev\_yoy} / \texttt{np\_yoy} & REAL & 营收 / 净利同比（\%） \\
\texttt{name} & TEXT & 该期披露时的证券简称（还原 ST 历史） \\
\midrule
\multicolumn{3}{@{}l}{\textbf{valuation} —— 月度估值快照，主键 \texttt{(code, date)}}\\
\texttt{total\_mv} & REAL & 总市值（元） \\
\texttt{float\_mv} & REAL & 流通市值（元） \\
\texttt{total\_shares} / \texttt{free\_shares} & REAL & 总股本 / 自由流通股本（股） \\
\texttt{pe\_ttm} / \texttt{pb\_mrq} & REAL & 市盈率 / 市净率 \\
\texttt{ps\_ttm} / \texttt{pcf\_ttm} & REAL & 市销率 / 市现率 \\
\texttt{board} & TEXT & 行业板块 \\
\midrule
\multicolumn{3}{@{}l}{\textbf{index\_member} --- 指数成分区间，主键 \texttt{(index\_code, code, in\_date)}}\\
\texttt{in\_date} & TEXT & 被观察到属于该指数的最早日期 \\
\texttt{out\_date} & TEXT & 最后一次仍属于该指数的日期（仍在册为空） \\
当日成分查询 & --- & \texttt{in\_date <= t AND (out\_date IS NULL OR out\_date > t)} \\
\midrule
\multicolumn{3}{@{}l}{\textbf{index\_bar} / \textbf{trade\_calendar} / \textbf{fetch\_log}}\\
\texttt{index\_bar} & — & 指数日线，用于基准与归因 \\
\texttt{trade\_}\newline\texttt{calendar} & — & 交易日历（由指数日期生成） \\
\texttt{fetch\_log} & — & 断点续抓日志 \\
\bottomrule
\end{longtable}

\chapter{代码索引}\label{app:code}

\begin{longtable}{@{}p{0.30\textwidth}p{0.62\textwidth}@{}}
\toprule
\textbf{文件} & \textbf{职责} \\
\midrule
\endhead
\texttt{aq/config.py} & 路径、时间范围、股票池口径 \\
\texttt{aq/http.py} & 双通道 HTTP（直连 / 系统代理交替重试） \\
\texttt{aq/sources.py} & 数据源适配（东财数据中心、行情） \\
\texttt{aq/db.py} & SQLite 读写、迁移、批量导入 \\
\texttt{aq/schema.sql} & 表结构定义 \\
\texttt{aq/panel.py} & 研究面板：复权、PIT、TTM、可成交标记、标签 \\
\texttt{aq/factors.py} & 因子库：39 个因子 + 注册表 + 宽表引擎 \\
\texttt{aq/neutralize.py} & 去极值、标准化、行业市值中性化 \\
\texttt{aq/metrics.py} & IC / ICIR / 分层 / 换手 / 衰减 / 因子相关性 \\
\texttt{aq/cv.py} & Purged K-Fold、CPCV、Walk-Forward \\
\texttt{aq/stats.py} & PSR、DSR、PBO、Reality Check、Bootstrap \\
\midrule
\texttt{scripts/01\_financials.py} & 抓取全部财报期（含公告日与简称历史） \\
\texttt{scripts/02\_prices.py} & 抓取全市场日线（新浪通道，断点续抓） \\
\texttt{scripts/03\_calendar.py} & 交易日历与指数行情 \\
\texttt{scripts/04\_valuation.py} & 月度估值快照 \\
\texttt{scripts/05\_finalize.py} & 上市/退市日、数据质量报告 \\
\texttt{scripts/09\_index\_}\newline\texttt{members.py} & 指数成分区间快照（可重复运行以积累历史） \\
\texttt{scripts/00\_update.py} & 一键增量更新（幂等、可中断） \\
\texttt{scripts/check\_db.py} & 数据库巡检 \\
\texttt{scripts/06\_factors.py} & 逐因子计算 + 单因子评价 + 面板落盘 \\
\texttt{scripts/07\_model.py} & 多模型、多协议重复测试 + DSR/PBO \\
\texttt{scripts/08\_report.py} & 生成图表与拟合精度报告 \\
\bottomrule
\end{longtable}

\chapter{常用公式速查}\label{app:formula}

\begin{align}
  \text{后复权价} &= \text{原始价} \times f_t \\
  \text{TTM} &= \text{累计}_t + \text{上年年报} - \text{上年同期累计} \\
  \text{中性化} &: \ x = \alpha + \beta_{\text{ind}} + \gamma\ln(\text{MktCap}) + \varepsilon,
  \quad \text{取 } \varepsilon \\
  \text{Rank IC}_t &= \mathrm{corr}\big(\mathrm{rank}(f_{i,t}), \mathrm{rank}(r_{i,t+1\to t+h})\big) \\
  \text{ICIR}_{\text{ann}} &= \frac{\overline{\IC}}{\mathrm{std}(\IC)}\sqrt{252/h} \\
  t_{\text{NW}} &= \frac{\overline{\IC}}{\sqrt{\hat\sigma^2_{\text{NW}}/n}} \\
  \text{Purge} &: \ [\,t_a - h,\ t_b\,] \ \text{剔除} \\
  \text{Embargo} &: \ [\,t_b+1,\ t_b+w\,] \ \text{剔除} \\
  \varphi(N,k) &= \binom{N-1}{k-1} \\
  \mathrm{PSR} &= \Phi\!\left(\frac{(\widehat{\SR}-\SR^*)\sqrt{T-1}}
  {\sqrt{1-\gamma_3\widehat{\SR}+\frac{\gamma_4-1}{4}\widehat{\SR}^2}}\right) \\
  \SR_0 &= \sigma_{\SR}\left[(1-\gamma)\Phi^{-1}\!\left(1-\tfrac1N\right)
  +\gamma\Phi^{-1}\!\left(1-\tfrac{1}{Ne}\right)\right] \\
  N_{\text{eff}} &= \frac{(\sum\lambda_i)^2}{\sum\lambda_i^2} \\
  \text{Impact} &\approx k\sigma\sqrt{Q/V}
\end{align}

\chapter{复现步骤}\label{app:reproduce}

\begin{enumerate}
  \item 安装依赖：\texttt{python -m pip install -r requirements.txt}
  \item 初始化数据库并抓取财报：\texttt{python scripts/01\_financials.py}
  \item 抓取全市场日线（可断点续抓）：\texttt{python scripts/02\_prices.py}
  \item 生成交易日历与指数：\texttt{python scripts/03\_calendar.py}
  \item 抓取估值快照：\texttt{python scripts/04\_valuation.py}
  \item 生成数据质量报告：\texttt{python scripts/05\_finalize.py --save}
  \item 抓取指数成分：\texttt{python scripts/09\_index\_members.py}
  \item 计算因子与单因子评价：\texttt{python scripts/06\_factors.py}
  \item 多协议重复测试与过拟合诊断：\texttt{python scripts/07\_model.py}
  \item 生成图表与报告：\texttt{python scripts/08\_report.py}
  \item 编译讲义：\texttt{xelatex main.tex}（在 \texttt{lectures/} 目录下执行两遍）
\end{enumerate}

\begin{warnbox}[关于数据源的时效性]
第 2 章记录的数据源可用性只在 2026-10-31 当天成立。免费接口随时可能变化。
如果某一步失败，请优先阅读 \texttt{scripts/02\_prices.py} 顶部的「数据源决策记录」，
那里列出了当时验证过的备选通道。
\end{warnbox}

\chapter{外部意见的交叉核验}\label{app:crosscheck}

本书的框架在写作过程中接受了两条独立的外部意见通道（Grok 与 Gemini，2026-10-31），
以及一份真实数据的检验。这一章记录：他们说了什么、哪些被采纳、哪些被数据推翻、
哪些被本书修正。保留这段记录，是因为\textbf{"我参考了谁"和"我为什么采信"本身
就是研究可信度的一部分}。

\section{两条通道的分工}

\begin{center}
\begin{tabular}{@{}p{0.16\textwidth}p{0.34\textwidth}p{0.40\textwidth}@{}}
\toprule
通道 & 强项 & 本书用它问什么 \\
\midrule
Grok & 实务细节、数据源现状、参数区间 & 数据源可用性、因子清单、成本参数、模型预期 \\
Gemini & 数学推导、文献细节、常见误用 & DSR/CSCV 公式、purge/embargo、NW 滞后阶数、FDR \\
\bottomrule
\end{tabular}
\end{center}

\section{被采纳的意见}

\subsection{关于因子方向}

两条通道都指出：\textbf{A 股日频上中长期动量偏弱甚至为负，短期反转主导}。
Grok 的原话是"经典中长期动量在 A 股长期偏弱甚至为负"。

本书的实测完全支持这一点：20 日动量的 Rank IC 约为 $-0.05$ 到 $-0.08$。
这条意见直接改变了第~\ref{ch:factor} 章的写法——书中明确写出"动量方向必须实证，
不能照搬教科书"。

\subsection{关于印花税}

Grok 指出 2023-08-28 起证券交易印花税由 $0.1\%$ 降至 $0.05\%$（单边卖出）。
本书核对后确认，并据此把双边显性成本估算从约 $0.13\%$ 修正为约 $0.08\%$--$0.11\%$。

\subsection{关于 IC 的年化口径}

Gemini 特别强调：日频算 IC 但预测 $h$ 日收益时，年化乘数应为 $\sqrt{252/h}$
而不是 $\sqrt{252}$，否则会虚增 $2$--$4$ 倍。本书第~\ref{ch:eval} 章把这个陷阱
单独列为一节，并在实现中采用 $\sqrt{252/5}$。

\subsection{关于 Purge 与 Embargo 的区间}

Gemini 给出了完整形式：训练样本需剔除
$[t_{\text{test,start}} - h - w,\ t_{\text{test,stop}} + h + w]$，
其中 $h$ 是标签持有期、$w$ 是特征回看窗。本书的 \texttt{aq/cv.py} 采用其工程简化版
（左侧 purge 到 $t_0-h$，右侧 embargo 到 $t_1+w$），并在第~\ref{ch:overfit} 章说明了理由。

\subsection{关于多重检验}

Gemini 给出了 Benjamini-Hochberg 的完整步骤，并指出一个关键细节：
\textbf{必须从大到小找最大的 $k^*$，而不是从小到大遇到第一个不满足就停}。
本书第~\ref{ch:eval} 章的代码实现遵循这一点。

\section{被数据修正的意见}

\subsection{"中性化后 IC 一定更高"}

Grok 在初稿中倾向于"中性化能提升因子质量"。本书的实测显示：\textbf{并非如此}。
部分因子的中性化后 IC 明显低于原始 IC（例如规模因子对市值中性化后退化为零）。
第~\ref{ch:prep} 章因此专门写了"中性化不总是提升 IC"一节。

\subsection{"模型能带来显著提升"}

Grok 给出的预期是"GBDT 比线性提升 0.01--0.02 的 IC"。本书的实测显示，
在严格 Purged 验证下，LightGBM 相对等权基线的提升\textbf{幅度有限且不稳定}
（见第~\ref{ch:result} 章的协议表）。这与"模型是次要的、因子与验证才是主要的"
这一判断一致。

\subsection{数据源的"可用性"}

两条通道都给出了数据源清单，但没有一条预见到当天的实际情况：
东方财富在约 50 次请求后被 IP 封锁，腾讯在约 200 次后触发 WAF，
baostock 并发即被限速。最终可用的是新浪。
\textbf{这说明"数据源对比表"这类静态知识，在工程实践中的保质期极短。}
第~\ref{ch:data} 章因此把这段排雷过程完整保留下来，而不是给一张漂亮的对比表。

\section{两条通道之间的分歧}

\begin{center}
\begin{tabular}{@{}p{0.16\textwidth}p{0.30\textwidth}p{0.30\textwidth}p{0.14\textwidth}@{}}
\toprule
议题 & Grok 的观点 & Gemini 的观点 & 本书处理 \\
\midrule
三重障碍法 & 认为日频选股不必优先使用 & 未直接评价，但强调重叠标签问题 & 采纳 Grok：日频不用 \\
深度模型 & 认为 2026 年值得一试 & 未涉及 & 保留但降级为"可选" \\
PBO 的解读 & 未深入 & 强调不能判断绝对收益 & 采纳 Gemini，单列一节 \\
换手与成本 & 给出 15--40 bp 区间 & 未涉及 & 采纳 Grok，取 30 bp \\
\bottomrule
\end{tabular}
\end{center}

\section{核验原则}

本书对外部意见的处理遵循三条原则：

\begin{enumerate}
  \item \textbf{可核验的事实回原始来源}。印花税率、指数编制规则、法规条文，
        一律回官方口径，不采信二手转述。
  \item \textbf{可计算的结论自己算一遍}。DSR、PBO、ICIR 的口径，全部用本书的实现重算，
        不直接引用别人的数字。
  \item \textbf{有分歧的地方做实验}。因子方向、模型提升幅度这类问题，
        用本书的数据跑一遍，以实测为准。
\end{enumerate}

\begin{warnbox}[一个必须承认的局限]
外部意见通道（以及本书作者）都可能产生幻觉。本章记录的是"参考了什么"，
而不是"因为谁说了所以对"。所有进入正文的结论，都在本书的数据上重新验证过；
未能验证的，正文中会明确标注为"经验判断"或"待验证"。
\end{warnbox}

\chapter{外部数据快照的接入与核验}\label{app:shm}

\section{为什么单开一章}

2026 年 10 月，本项目获得了一份外部 A 股数据快照（四张 parquet，覆盖 2020-01 至 2026-08）。
它包含本地库完全没有的两样东西：\textbf{四档资金流}与\textbf{121 个财务指标}。

但"能拿到"不等于"能用"，更不等于"应该替换"。这一章记录四项核验的做法与结果，
以及两个更重要的结论：

\begin{enumerate}
  \item \textbf{看起来更全的数据，可能在一个更隐蔽的维度上更窄}——它缺的是最小那三成市值；
  \item \textbf{你以为需要它来修正的东西，可能你本来就已经很准}——本地估算的成交额
        中位误差只有 $0.3\%$。
\end{enumerate}

\section{数据结构}

\begin{center}
\begin{tabularx}{\textwidth}{@{}lrrX@{}}
\toprule
表 & 行数 & 列数 & 关键内容 \\
\midrule
\texttt{market\_daily} & 5,637,878 & 24 & 不复权与后复权 OHLC、成交额、VWAP/TWAP、换手率、流通股本 \\
\texttt{market\_moneyflow} & 5,637,848 & 33 & 超大单/大单/中单/小单的流入流出金额、量、笔数 \\
\texttt{valuation\_daily} & 9,553,741 & 15 & PE/PB/PS/PCF/PEG/股息率/资产负债率，日频 \\
\texttt{fundamental\_daily} & 5,637,878 & 123 & 121 个财务指标，含单季与 TTM \\
\bottomrule
\end{tabularx}
\end{center}

\section{三项核验}

\subsection{第一项：价格一致性}

把外部快照的收盘价与本地库（新浪通道，独立采集）按 \texttt{(code, date)} 内连接，
共匹配 \textbf{5,544,219} 行。

\begin{keybox}[结果]
收盘价相对误差超过 $0.1\%$ 的比例为 \textbf{0.000\%}。
两个完全独立的采集通道给出了一致的价格，这同时验证了双方的数据质量。
\end{keybox}

\subsection{第二项：Point-in-Time}

这是最容易被忽略、也最致命的一项。财务数据必须按\textbf{公告日}对齐，
而不是报告期。检验方法是看净利润这个字段在哪个交易日发生变化。

以平安银行（000001）为例：

\begin{center}
\begin{tabular}{@{}lll@{}}
\toprule
数值变化日 & 归母净利润（累计，元） & 对应的报告 \\
\midrule
2020-02-14 & $2.8195\times10^{10}$ & 2019 年年报 \\
2020-04-21 & $8.5480\times10^{9}$ & 2020 年一季报 \\
2020-08-28 & $1.3678\times10^{10}$ & 2020 年中报 \\
2020-10-22 & $2.2398\times10^{10}$ & 2020 年三季报 \\
\bottomrule
\end{tabular}
\end{center}

四个日期与平安银行真实公告日完全吻合。\textbf{这份快照的财务数据是 PIT 的}，
可以直接用于建模。浦发银行（600000）的检验结果同样吻合。

\begin{warnbox}[如果这一项没通过]
即便价格、成交量全部正确，用报告期对齐的财务数据也会让模型提前 1--4 个月
"知道"财报内容。这种前视偏差不会报错，只会让回测好看得离谱。
所以接入任何新数据源，这一项必须第一个做。
\end{warnbox}

\subsection{第三项：覆盖偏差}

前两项都通过之后，第三项暴露了问题。

\begin{center}
\begin{tabular}{@{}lrrr@{}}
\toprule
交易日 & 外部快照 & 本地库 & 缺口 \\
\midrule
2022-06-30 & 3,104 & 4,556 & $-1,452$ \\
2024-06-28 & 3,072 & 5,025 & $-1,953$ \\
2026-08-31 & 4,186 & 5,202 & $-1,016$ \\
\bottomrule
\end{tabular}
\end{center}

外部快照的股票池只有全市场的 \textbf{76\%--83\%}，而且\textbf{各年比例不同}
（2021 年日均只有 2,818 只，是六年里最低的）。缺失的既有 ST 股，也有万科A
这样的蓝筹。

\begin{keybox}[这份快照的正确用法]
它是一个\textbf{已过滤的股票池}，不是全市场。因此：
\begin{enumerate}
  \item \textbf{不能}用它替代本地库做全市场横截面研究——股票池的边界随年份变化，
        会直接污染因子排序；
  \item \textbf{可以}用它做三件事：补充真实成交额与换手率、引入资金流因子族、
        提供日频估值；
  \item \textbf{必须}在任何相关结论旁注明覆盖范围。
\end{enumerate}
\end{keybox}

\subsection{第四项：本地估算到底有多不准}

第三项说明了不能替换股票池。但还有一个更基本的问题：这份快照宣称的价值之一是
"提供真实成交额与换手率，而本地只能估算"。那本地到底估得多差？

用 554 万行匹配数据实测：

\begin{center}
\begin{tabular}{@{}llrrr@{}}
\toprule
字段 & 本地的做法 & 中位误差 & P90 & P99 \\
\midrule
成交额 & 成交量 $\times$ 当日均价 & $0.31\%$ & $1.03\%$ & $2.37\%$ \\
流通股本 & 月度快照前向填充 & $0.000\%$ & $0.022\%$ & $21.9\%$ \\
换手率 & 成交量 $/$ 流通股本 & $\approx 1\%$ & -- & -- \\
\bottomrule
\end{tabular}
\end{center}

成交额的估算中位误差 $0.31\%$、99 分位 $2.37\%$，这个精度对流动性因子
（Amihud 非流动性、成交额对数、流动性门槛）完全够用。
流通股本在 96\% 的交易日上与真值一致，只有 1.1\% 的行偏差超过 20\%，
集中在送转当年。

\begin{keybox}[这笔交易不划算]
"用真实成交额替换估算值"这个卖点，在本项目里收益接近零（$0.3\%$ 的精度提升），
而代价是引入 20\% 的样本缺口与一个逐年变化的股票池边界。
\textbf{为了 $0.3\%$ 的精度牺牲三成小盘股，不划算。}

这也是一个通用教训：接入新数据源时，先量化"我原来有多差"，
再决定要不要替换。很多人换数据源换掉的不是误差，而是可复现性。
\end{keybox}

\section{资金流因子族：一次"换皮检验"的完整演示}

这是本章最有教学价值的部分。资金流数据看起来信息量很大，
但必须先回答一个问题：\textbf{它是不是"当日涨跌幅"或"换手率"的换皮？}

A 股的主力净流入与当日涨跌高度同向，而涨跌本身已经是反转因子的原料。
如果不做检验就直接把 33 个字段丢进模型，得到的很可能是"更贵的反转因子"。

\subsection{因子构造}

所有资金流都除以当日成交额，化为无量纲的"净流入强度"：

\begin{align}
  \text{mfr} &= \frac{\text{主力净流入}}{\text{成交额}}, \qquad
  \text{sfr} = \frac{\text{超大单净流入}}{\text{成交额}}, \\
  \text{lfr} &= -\frac{\text{小单净流入}}{\text{成交额}}, \\
  \text{mfr\_20} &= \frac{\sum_{20}\text{主力净流入}}{\sum_{20}\text{成交额}}, \\
  \text{mfr\_pos\_20} &= \frac{1}{20}\sum_{20}\mathbb{1}[\text{mfr} > 0].
\end{align}

\subsection{统一标签与统计口径后的复核}
旧版此处误用收盘到收盘标签。现在与主面板共用次日复权开盘至第六日复权开盘收益，
即 $y_t=O^{adj}_{t+6}/O^{adj}_{t+1}-1$。标签按个股观测行位移，缺失交易日仍需另做日历对齐审计。
调整后 IC 控制五日反转与二十日换手；没有把这些新因子加入主模型重训。
{\small
\begin{longtable}{@{}lrrrrrr@{}}
\toprule
因子 & 原始 IC & 调整后 IC & 原始 $t$ & 调整后 $t$ & $n$ & $n_{\rm eff}$ \\
\midrule
\endhead
mfr\_pos\_20 & -0.0151 & 0.0246 & -4.68 & 5.41 & 1598 & 386 \\
mfr\_chg & -0.0060 & 0.0006 & -2.09 & 0.29 & 1598 & 522 \\
sfr & -0.0061 & 0.0004 & -3.88 & 0.28 & 1609 & 1187 \\
bfr & -0.0045 & -0.0029 & -2.40 & -1.89 & 1609 & 962 \\
mfr\_stab\_20 & -0.0164 & -0.0053 & -4.52 & -1.66 & 1598 & 400 \\
mfr & -0.0108 & -0.0061 & -5.26 & -3.71 & 1609 & 1097 \\
mfr\_20 & -0.0194 & -0.0080 & -5.70 & -2.75 & 1598 & 408 \\
lfr & -0.0117 & -0.0083 & -5.20 & -4.29 & 1609 & 973 \\
\bottomrule
\end{longtable}
}

原始 IC 与调整后 IC 均由完整日频 IC 序列汇总；$h=5$ 的日频标签重叠，
均值检验使用 Bartlett Newey--West 协方差，滞后阶数为
$L=\max(4,\lfloor4(n/100)^{2/9}\rfloor)$。$n$ 是有值的 IC 日期数，
$n_{\rm eff}=n\gamma_0/\widehat\Omega$ 是原始 IC 均值方差等价的样本量，
不是独立财报次数，也不是未经假设可确定的真实样本数。
完整表另含调整后有效样本量；不同带宽及每五日采样的敏感性见配套 CSV。
本表不再平均各年的 $t$ 值。

\subsection{怎么读这张表}
\texttt{mfr} 的原始 IC 为 \AuditFlowMfr{}；\texttt{mfr\_pos\_20} 的调整后 IC 为
\AuditFlowPos{}，HAC $t=\AuditFlowPosT$。这些数字替换旧收盘标签的结论。
不能仅凭与反转相关就认定某因子“只有换皮”，也不能仅凭残差 IC 显著就证明独立可交易收益。
残差化控制变量、股票池、成本、多重筛选和时间稳定性都会影响解释。
本轮对八个因子全部保留结果，没有事后只展示最优一个。

\section{本章小结}
补充库覆盖 2020--2026 的变化股票池，不能替代主库。
统一开盘标签提高了定义的一致性，但股票池、区间和控制变量仍不同，跨库结果不能直接排名。
价格匹配与个别公告日核验是局部证据；2021 年整年缺失证券及筛选规则仍待解释。

\chapter{质量因子族：重算显著性与有效样本}\label{app:quality}

\section{动机}

第~\ref{ch:result} 章有一个不起眼的结论：本地的质量因子（ROE、毛利率、
现金含量、应计）的旧年度汇总 ICIR 较低。该结果不足以判断全样本显著性。

当时给出的解释是"字段太少"——本地库只有 9 个财务字段。外部快照带来 121 个指标，
其中本地确实拿不到的有：研发费用、商誉、无形资产、有息负债、营运资本、
费用明细、扣非利润、投资收益。

于是有了这一章的实验：\textbf{补充字段是否带来增量预测关联？}

\section{二十个因子}

从 121 个指标里构造 20 个质量因子，统一方向为"数值越大越好"：

\begin{center}
\small
\begin{tabular}{@{}lll@{}}
\toprule
类别 & 因子 & 构造 \\
\midrule
盈利能力 & \texttt{q\_wroe} / \texttt{q\_roe} / \texttt{q\_roa} & 加权 ROE / ROE / ROA \\
 & \texttt{q\_gm} / \texttt{q\_opm} & 毛利率 / 营业利润率 \\
现金质量 & \texttt{q\_cfo\_ni} & 经营现金流 TTM / $|$归母净利 TTM$|$ \\
 & \texttt{q\_accrual} & $-$（归母净利 $-$ 经营现金流）/ 总资产 \\
 & \texttt{q\_cash\_rev} & 经营流入小计 / 营业总收入 \\
研发 & \texttt{q\_rd} & 研发费用 TTM / 营业总收入 TTM \\
商誉 & \texttt{q\_gw} & $-$商誉 / 归母权益 \\
 & \texttt{q\_intang} & $-$（商誉 $+$ 无形）/ 总资产 \\
杠杆 & \texttt{q\_lev} & $-$总负债 / 总资产 \\
 & \texttt{q\_debt} & $-$（短期借款 $+$ 长期借款）/ 总资产 \\
偿债 & \texttt{q\_current} / \texttt{q\_quick} & 流动比率 / 速动比率 \\
营运 & \texttt{q\_wc} & （流动资产 $-$ 流动负债）/ 总资产 \\
 & \texttt{q\_at} & 营业总收入 TTM / 总资产 \\
费用 & \texttt{q\_exp} & $-$（管理 $+$ 销售 $+$ 财务费用）/ 营业总收入 \\
利润含金量 & \texttt{q\_cut} & 扣非净利 / $|$归母净利$|$ \\
 & \texttt{q\_invinc} & $-$投资收益 / $|$归母净利$|$ \\
\bottomrule
\end{tabular}
\end{center}

\section{重算结果}
本节使用补充库、2020--2026 的次日开盘至第六日开盘标签。
调整后 IC 控制对数市值、二十日换手、波动和五日反转。
{\small
\begin{longtable}{@{}lrrrrrr@{}}
\toprule
因子 & 原始 IC & 调整后 IC & 原始 $t$ & 调整后 $t$ & $n$ & $n_{\rm eff}$ \\
\midrule
\endhead
q\_exp & 0.0085 & 0.0185 & 1.70 & 3.97 & 1609 & 410 \\
q\_wroe & 0.0029 & 0.0100 & 0.44 & 2.16 & 1609 & 378 \\
q\_roa & 0.0031 & 0.0083 & 0.51 & 1.68 & 1609 & 376 \\
q\_cut & 0.0064 & 0.0079 & 1.24 & 1.84 & 1609 & 382 \\
q\_accrual & 0.0063 & 0.0079 & 1.93 & 2.56 & 1609 & 363 \\
q\_roe & 0.0033 & 0.0072 & 0.51 & 1.57 & 1609 & 378 \\
q\_opm & 0.0049 & 0.0068 & 0.90 & 1.26 & 1609 & 376 \\
q\_invinc & -0.0047 & 0.0061 & -1.60 & 1.09 & 1609 & 401 \\
q\_intang & 0.0035 & 0.0040 & 1.15 & 1.84 & 1609 & 416 \\
q\_at & 0.0031 & 0.0028 & 0.99 & 0.88 & 1609 & 394 \\
q\_debt & -0.0002 & 0.0017 & -0.06 & 0.65 & 1609 & 425 \\
q\_cfo\_ni & 0.0109 & 0.0010 & 3.27 & 0.24 & 1609 & 371 \\
q\_quick & -0.0000 & 0.0006 & -0.01 & 0.14 & 1609 & 399 \\
q\_rd & -0.0034 & 0.0005 & -0.37 & 0.08 & 1609 & 376 \\
q\_current & -0.0004 & 0.0000 & -0.08 & 0.01 & 1609 & 396 \\
q\_wc & -0.0007 & -0.0000 & -0.12 & -0.00 & 1609 & 389 \\
q\_lev & 0.0012 & -0.0002 & 0.24 & -0.05 & 1609 & 403 \\
q\_gw & 0.0098 & -0.0006 & 2.39 & -0.19 & 1609 & 391 \\
q\_gm & -0.0049 & -0.0038 & -1.00 & -0.82 & 1609 & 415 \\
q\_cash\_rev & -0.0028 & -0.0055 & -0.63 & -1.20 & 1609 & 370 \\
\bottomrule
\end{longtable}
}

原始 IC 与调整后 IC 均由完整日频 IC 序列汇总；$h=5$ 的日频标签重叠，
均值检验使用 Bartlett Newey--West 协方差，滞后阶数为
$L=\max(4,\lfloor4(n/100)^{2/9}\rfloor)$。$n$ 是有值的 IC 日期数，
$n_{\rm eff}=n\gamma_0/\widehat\Omega$ 是原始 IC 均值方差等价的样本量，
不是独立财报次数，也不是未经假设可确定的真实样本数。
完整表另含调整后有效样本量；不同带宽及每五日采样的敏感性见配套 CSV。
本表不再平均各年的 $t$ 值。

\section{撤回“有效样本仅 14 个”的推导}
旧程序先逐年算 $t$，然后求年度 $t$ 的平均。这个平均值不是全样本的检验统计量，
把它与全样本朴素 $t$ 代入 $n(t_{NW}/t_{naive})^2$，得不到可解释的有效样本量。
旧版“20 个因子全部不显著”“预测能力为零”和“质量只适合硬筛选”均不再作为本实验结论。

以 \texttt{q\_cfo\_ni} 为例，本次保存了 \AuditQCfoCount{} 个实际 IC 观测；
朴素 $t=\AuditQCfoNaive$，从同一序列重算的 HAC $t=\AuditQCfoT$，
方差等价样本量约为 \AuditQCfoN{}。这不是“独立季报数”。
慢变因子值与每日收益的横截面相关可以每日变化，不能用财报更新次数替代 IC 相关性测量。
主库的 size、bp、cfp 还包含随股价变动的分母，不能说“一年只变化几次”。

\section{还需要什么证据}
\begin{enumerate}
  \item 对照 5、20、60 日 HAC 带宽及每五日采样结果，检查显著性是否依赖选择；完整结果在敏感性 CSV。
  \item 对全部被尝试的因子和参数进行多重检验；这里的逐因子 $t$ 不是全搜索校正后的显著性。
  \item 核实 PIT 字段、股票池缺失机制和实际收益跨度，再用新的时间段验证。
  \item 质量筛选、长周期持有与主模型增量贡献各自需要实验；本表不证明哪一种用法一定更好。
\end{enumerate}
不显著不能证明无效，显著也不能证明可交易。首先保证统计量来自同一组观测，再解释经济含义。

\chapter{文献严格化：$t>3$ 门槛与剔除微盘}\label{app:strict}

\section{为什么单开一章}

2026 年 11 月，一份外部讲义（\textit{AI 炒股与量化交易：成熟技术讲义}，
16 章 + 4 附录）提供了两条对本项目有直接冲击的文献结论：

\begin{enumerate}
  \item \textbf{Harvey, Liu \& Zhu (2016)}：在数百个已发表因子之后，
        $t>2.0$ 的显著性门槛已不够，新因子应要求 $t>3.0$。
  \item \textbf{Liu, Stambaugh \& Yuan (2019)}：A 股的小市值效应被"壳价值"污染，
        建议剔除最小 30\% 市值后再检验。
\end{enumerate}

本书此前的所有因子判定都用 $t>2$，也从未做过剔除微盘的稳健性检验。
这一章把两条建议变成可执行的检验。

\begin{warnbox}[样本口径说明]
本章的两张表用了不同区间：严格性检验（表~\ref{tab:strict}）用 2013--2026，
因为需要 20 日滚动变量；$t$ 门槛重判用全样本 2010--2026（来自第~\ref{ch:result} 章的因子表）。
同一个因子在两表里的数值略有差异，源于区间不同，不是错误。
\end{warnbox}

\section{检验一：剔除最小 30\% 市值}

按当日对数流通市值分位，剔除最小的 30\%，在剩余样本上重算 Rank IC。

\begin{table}[htbp]
\centering
\caption{全样本 vs 剔除最小 30\% 市值（$h=5$，2013--2026）}
\label{tab:strict}
\small
\begin{tabular}{@{}lrrrrrr@{}}
\toprule
因子 & IC 全样本 & IC 剔除小30\% & IC 变化 & $t$ 全样本 & $t$ 剔除后 & 判定 \\
\midrule
max\_ret\_20 & +0.0592 & +0.0538 & -9.1\% & +17.29 & +13.94 & 稳健 \\
turn\_20 & +0.0585 & +0.0504 & -13.8\% & +15.44 & +12.61 & 稳健 \\
pv\_corr\_20 & +0.0346 & +0.0310 & -10.5\% & +15.03 & +12.70 & 稳健 \\
skew\_20 & +0.0260 & +0.0202 & -22.2\% & +13.44 & +10.17 & 减弱 \\
vol\_20 & +0.0606 & +0.0566 & -6.6\% & +12.81 & +10.80 & 稳健 \\
ln\_amount\_20 & +0.0590 & +0.0551 & -6.6\% & +12.68 & +10.43 & 稳健 \\
mom\_20 & -0.0500 & -0.0437 & -12.6\% & -12.22 & -10.08 & 稳健 \\
ivol\_60 & +0.0698 & +0.0677 & -3.0\% & +11.10 & +9.72 & 稳健 \\
rev\_10 & +0.0380 & +0.0323 & -15.1\% & +10.43 & +8.47 & 稳健 \\
rev\_5 & +0.0308 & +0.0259 & -16.1\% & +9.63 & +7.69 & 稳健 \\
rev\_1 & +0.0159 & +0.0135 & -15.2\% & +8.66 & +7.00 & 稳健 \\
bp & +0.0339 & +0.0317 & -6.6\% & +8.50 & +6.55 & 稳健 \\
amihud\_20 & +0.0231 & +0.0149 & -35.6\% & +7.46 & +3.88 & 减弱 \\
mom\_120 & -0.0412 & -0.0389 & -5.7\% & -6.24 & -5.26 & 稳健 \\
size & +0.0097 & -0.0050 & -152.2\% & +1.26 & -0.76 & \textbf{翻负，不可用} \\
\bottomrule
\end{tabular}

\end{table}

\subsection{怎么读}

\begin{enumerate}
  \item \textbf{反转、波动、换手、流动性四族非常稳健}。\texttt{ivol\_60} 只掉 $3.0\%$，
        \texttt{vol\_20} 与 \texttt{ln\_amount\_20} 各掉 $6.6\%$，
        剔除微盘之后 $t$ 值仍在 $9.7$--$10.8$。这说明它们的收益不依赖最小的那三成。
  \item \textbf{短期反转减弱但存活}。\texttt{rev\_5} 掉 $16.1\%$，$t$ 从 $9.63$ 降到 $7.69$。
        微盘确实贡献了一部分反转收益，但远不是全部。
  \item \textbf{Amihud 非流动性掉得最多}（$-35.6\%$）。这符合逻辑：
        非流动性溢价本来就集中在流动性最差的股票上，而它们正是被剔除的那批。
  \item \textbf{规模因子直接翻负}。\texttt{size} 的 IC 从 $+0.0097$ 变成 $-0.0050$，
        $t$ 从 $+1.26$ 变成 $-0.76$。
\end{enumerate}

\begin{keybox}[这一条最重要]
\texttt{size} 的正 IC \textbf{100\% 来自最小那 30\% 市值}。
一旦剔除，它不仅消失，还翻成负号。

这意味着：本项目此前把规模因子列为"衰减"是过于客气的说法。
准确的说法是——\textbf{它在可执行的那部分市场里根本不存在}。

它与三件事同时吻合：LSY (2019) 的"壳价值污染"、外部讲义第 13 章的
"2024 年初微盘股流动性危机"，以及附录~\ref{app:shm} 发现的
"外部快照恰好缺最小那三成"。三者指向同一个结论：
\textbf{微盘股上的"规律"既不可执行，也不可信。}
\end{keybox}

\section{检验二：把门槛提到 $t>3$}

按 Harvey-Liu-Zhu 的建议，用 $|t|>3$ 重新判定第~\ref{ch:result} 章的 39 个因子：

\begin{tabular}{@{}lr@{}}
\toprule
判定 & 因子数 \\
\midrule
$|t|>3$（HLZ 建议门槛） & \textbf{31} \\
仅 $2<|t|\le 3$ & 4 \\
$|t|\le 2$ & 4 \\
\bottomrule
\end{tabular}


仅通过 $t>2$、未过 $t>3$ 的四个因子：

\begin{tabular}{@{}llrr@{}}
\toprule
因子 & 家族 & IC & $t$ \\
\midrule
rev\_yoy & 成长 & +0.0084 & +2.94 \\
size & 规模 & +0.0167 & +2.27 \\
kurt\_20 & 波动 & +0.0050 & +2.13 \\
mom\_12\_1 & 动量 & -0.0687 & -2.07 \\
\bottomrule
\end{tabular}


\begin{keybox}[结论]
39 个因子里，\textbf{31 个通过 $|t|>3$}。这个比例远高于 Hou-Xue-Zhang (2020)
对已发表异象的复现结果（他们发现 447 个异象中 64\% 在 5\% 水平下不显著）。

两个可能的原因：
\begin{enumerate}
  \item 本项目的因子集中在\textbf{量价类}（反转、波动、换手、流动性），
        这一类在 A 股本来就更强，而 HLZ/HXZ 的样本以美股与基本面因子为主；
  \item $t$ 值是在\textbf{全序列}上算的 HAC $t$（观测数 3400 左右），
        而不是"各年 $t$ 的平均"，前者天然更大。
\end{enumerate}

因此这 31 个不应被理解为"31 个独立的新发现"，
而应理解为"31 个因子在样本内具有统计上难以用噪声解释的排序能力"。
它们之间的相关性很高（第~\ref{ch:factor} 章已讨论），
真正独立的维度远小于 31。
\end{keybox}

\section{对被淘汰因子与规模因子的处理}

\begin{enumerate}
  \item \textbf{规模因子（\texttt{size}）降级为"不可用"}：
        它只在最小 30\% 市值里有效，而那一区段容量极低、2024 年出现过流动性危机，
        且外部数据快照根本不覆盖。因子库保留其代码以备复现，但不再进入任何合成。
  \item \textbf{Amihud 与偏度标记为"减弱"}：剔除微盘后仍显著，但幅度掉三成以上，
        使用时必须注明"含微盘暴露"。
  \item \textbf{四个仅过 $t>2$ 的因子}（\texttt{rev\_yoy}、\texttt{size}、
        \texttt{kurt\_20}、\texttt{mom\_12\_1}）在结论中一律标注为"边缘"，
        不作为核心证据引用。
\end{enumerate}

\section{附：证据薄弱清单与识别红旗}

以下内容改编自外部讲义第 14 章，用于快速排除不合格的说法。

\begin{center}
\small
\begin{tabular}{@{}p{0.30\textwidth}p{0.60\textwidth}@{}}
\toprule
流行说法 & 证据状况与替代方案 \\
\midrule
"深度学习能预测股价" & 证据薄弱：非平稳 + 扣成本后优势消失。
改用 ML 做横截面排序并控制换手。 \\
"LLM 选股年化 30\%+" & 预训练记忆污染 + 不可交易。
只用模型 cutoff 之后的样本；做事件抽取而不是选股。 \\
"强化学习自动交易" & 样本效率与奖励黑客问题。用执行与做市，不用于选股。 \\
"因子越多越强" & 已被 HLZ (2016) 与 HXZ (2020) 证伪。少数稳健因子 + 风险模型。 \\
"夏普 3 的策略" & 大概率没算 DSR/PBO。先问"试了多少次"。 \\
"十年回测稳定盈利" & 多半未含退市股。要求全样本 + PIT。 \\
"高频稳赚" & 军备竞赛 + 监管 + 成本。转向中低频与成本控制。 \\
\bottomrule
\end{tabular}
\end{center}

\begin{warnbox}[识别"AI 炒股"产品的六个红旗]
\begin{enumerate}
  \item 只展示收益曲线，不展示回撤、换手率、容量；
  \item 拿不出连续 3 年以上的\textbf{实盘}（非模拟）记录；
  \item 拒绝说明策略逻辑与风险来源，只讲"神经网络""大模型"；
  \item 承诺保本或固定收益（在中国属于违规宣传）；
  \item 收益曲线过于平滑（夏普 $>4$ 且回撤 $<3\%$），真实市场几乎不存在；
  \item 靠卖软件、卖课、卖会员盈利，而不是靠自营交易盈利。
\end{enumerate}

一个可执行的检验：让对方用"未来 6 个月的纸上交易"证明，而不是展示过去。
\end{warnbox}

\section{本章小结}

\begin{enumerate}
  \item 剔除最小 30\% 市值后，反转、波动、换手、流动性四族仍然高度显著，
        IC 只下降 3\%--16\%。
  \item 规模因子的正 IC 全部来自最小那三成，剔除后翻负——
        应判定为\textbf{不可用}，而不是"衰减"。
  \item Amihud 非流动性下降 36\%，属于"减弱"，使用时须注明微盘暴露。
  \item 按 $|t|>3$ 重判，39 个因子中 31 个通过，4 个仅在 $t>2$ 水平上显著。
  \item "31 个通过 $t>3$"不等于 31 个独立发现；因子之间高度相关，
        真正独立的维度要少得多。
\end{enumerate}


\end{document}
